% \iffalse meta-comment
%
% Copyright (C) 2021 Universite Laval
%
% This file may be distributed and/or modified under the conditions
% of the LaTeX Project Public License, either version 1.3c of this
% license or (at your option) any later version. The latest version
% of this license is in:
%
%   http://www.latex-project.org/lppl.txt
%
% and version 1.3c or later is part of all distributions of LaTeX
% version 2006/05/20 or later.
%
% This work has the LPPL maintenance status `maintained'.
%
% The Current Maintainer of this work is Universite Laval
% <ulthese-dev@bibl.ulaval.ca>.
%
% This work consists of the files ulthese.dtx and ulthese.ins
% and the derived files listed in the README.md file.
%
% \fi
%
% \iffalse
%<*dtx>
\ProvidesFile{ulthese.dtx}
%</dtx>
%<class>\NeedsTeXFormat{LaTeX2e}[2009/09/24]
%<class>\ProvidesClass{ulthese}%
%<*class>
  [2021/08/09 v5.3a Universite Laval thesis and memoir class]
%</class>
%<*driver>
\documentclass[11pt,english,french]{ltxdoc}
  \usepackage[utf8]{inputenc}
  \usepackage[T1]{fontenc}
  \usepackage{natbib}
  \usepackage{babel}
  \usepackage[autolanguage]{numprint}
  \usepackage{microtype}
  \usepackage[scaled=0.92]{helvet}
  \usepackage[scaled=1.02]{inconsolata}
  \usepackage[sc]{mathpazo}
  \usepackage{fontawesome}
  \usepackage{metalogo}
  \usepackage{tabularx,booktabs,caption}
  \DisableCrossrefs
  \CodelineNumbered
  \RecordChanges
  \GlossaryPrologue{\section*{Historique des versions}%
    \addcontentsline{toc}{section}{Historique des versions}}

  %% Couleurs
  \usepackage{xcolor}
  \definecolor{link}{rgb}{0,0.4,0.6}   % ~RoyalBlue de dvips
  \definecolor{url}{rgb}{0.6,0,0}      % rouge foncé
  \definecolor{citation}{rgb}{0,0.5,0} % vert foncé

  %% Liste description alignée à gauche
  \usepackage{enumitem}
  \setlist[description]{leftmargin=*,align=left}

  %% Environnement pour remarques
  \usepackage{amsthm}
  \theoremstyle{definition}
  \newtheorem*{rem}{Remarque}

  %% Hyperliens
  \usepackage{hyperref}
  \hypersetup{%
    pdfauthor = {Faculté des études supérieures et postdoctorales},
    pdftitle = {Guide de l'utilisateur de la classe ulthese},
    colorlinks = {true},
    linktocpage = {true},
    urlcolor = {url},
    linkcolor = {link},
    citecolor = {citation},
    pdfpagemode = {UseOutlines},
    pdfstartview = {Fit},
    bookmarksopen = {true},
    bookmarksnumbered = {true},
    bookmarksdepth = {subsubsection}}
  \addto\extrasfrench{%
    \def\appendixautorefname{annexe}%
    \def\tableautorefname{tableau}%
    \def\subsectionautorefname{section}%
  }

  %% Paramétrage de babel
  \frenchbsetup{%
    StandardItemizeEnv=true,       % listes standards
    ThinSpaceInFrenchNumbers=true, % espace fine dans les nombres
    og=«, fg=»                     % « et » sont les guillemets
  }
  \def\frenchtablename{{\scshape Tab.}}

  %% Commandes de mise en forme spéciales
  \newcommand{\class}[1]{\textsf{#1}}
  \newcommand{\pkg}[1]{\textbf{#1}}

  %% Commandes pour les liens vers la documentation
  \newcommand{\link}[2]{\href{#1}{#2~\raisebox{-0.2ex}{\faExternalLink}}}
  \newcommand{\doc}[2][documentation]{\link{#2}{#1}}

  %% Francisation d'un mot clé de l'historique des commandes
  \def\generalname{Général}

\begin{document}
  \DocInput{ulthese.dtx}
\end{document}
%</driver>
% \fi
% \CheckSum{914}
% \DoNotIndex{\',\^,\`,\ ,\ae}
% \DoNotIndex{\RequirePackage,\ExecuteOptions,\ProcessOptions}
% \DoNotIndex{\newcommand,\newcommand*}
% \DoNotIndex{\setlength}
% \changes{4.4}{2017-05-16}{Prise en charge des bibliographies
%   multiples pour les thèses et mémoires par articles.}
% \changes{4.4}{2017-05-16}{Documentation révisée, principalement ce
%   qui touche à la préparation de la bibliographie.}
% \changes{4.3}{2017-03-01}{Nouvelles options pour les titres de
%   docteur en musique, maître en architecture, maître en
%   ergothérapie, maître en physiothérapie et maître en
%   psychoéducation.}
% \changes{4.3}{2017-03-01}{Modifications à la composition des sigles de
%   grades.}
% \changes{4.3}{2017-03-01}{Vérification de la compatibilité entre le
%   grade et les options multifacultaire, cotutelle, bidiplomation et
%   extension.}
% \changes{4.3}{2017-03-01}{Améliorations à la documentation.}
% \changes{4.2}{2016-03-29}{Ajout de l'option «examen».}
% \changes{4.1}{2016-02-13}{Paquetage geometry déclaré incompatible
%   avec la classe.}
% \changes{4.0}{2015-06-12}{Nouvelles règles de présentation matérielle
%   de la FESP: recto seulement, page frontispice.}
% \changes{3.1}{2014-05-23}{Prise en charge de la maîtrise en bidiplomation.}
% \changes{3.0a}{2014-03-24}{Modifications et corrections à la
%   documentation, notamment relativement à la configuration de natbib.}
% \changes{3.0}{2014-01-06}{Déclaration du grade en option de la
%   classe.}
% \changes{3.0}{2014-01-06}{Moteur XeLaTeX supporté.}
% \changes{3.0}{2014-01-06}{Ajout de l'option nobabel.}
% \changes{2.1}{2013-01-16}{Utilisation transparente de la police
%   Helvetica pour la page de titre.}
% \changes{2.0}{2013-01-13}{Traitement automatique des longs titres.}
% \changes{1.0b}{2012-11-11}{Ajouts et corrections mineures dans la
%   documentation.}
% \changes{1.0a}{2012-10-17}{Précisions dans la documentation.}
% \changes{1.0}{2012-09-30}{Version initiale.}
% \GetFileInfo{ulthese.dtx}
% \title{\class{ulthese}: la classe pour les thèses et mémoires de
%   l'Université Laval\thanks{Ce document décrit la classe
%   \class{ulthese}~\fileversion, datée du \filedate.}}
% \author{Faculté des études supérieures et postdoctorales\thanks{%
%   Cette classe et sa documentation ont été rédigées par Vincent
%   Goulet~(Faculté des sciences et de génie) avec la collaboration de
%   Koassi D'Almeida~(Faculté des études supérieures et postdoctorales) et
%   de Pierre Lasou~(Bibliothèque).}}
% \date{}
% \maketitle
%
% \section{Introduction}
%
% La classe \class{ulthese} permet de composer avec {\LaTeX} ou
% {\XeLaTeX} des thèses et mémoires immédiatement conformes aux règles
% générales de présentation matérielle de la Faculté des études
% supérieures et postdoctorales (FESP) de l'Université Laval. Ces
% règles définissent principalement la présentation des pages de titre
% des thèses et mémoires ainsi que la disposition du texte sur la
% page. La classe en elle-même est donc relativement simple.
%
% Cependant, \class{ulthese} est basée sur la classe \class{memoir}
% \citep{memoir}, une extension de la classe standard \class{book}
% facilitant à plusieurs égards la préparation de documents d'allure
% professionnelle dans {\LaTeX}. La classe \class{memoir} incorpore
% d'office plus de 30 des paquetages les plus populaires\footnote{%
% Consulter la section~18.24 de la documentation de \class{memoir}
% pour la liste ou encore le journal de la compilation (\emph{log})
% d'un document utilisant la classe \class{ulthese}.}. %
% L'intégralité des fonctionnalités de \class{memoir} se retrouve donc
% dans \class{ulthese}.
%
% La classe \class{memoir} fait partie des distributions {\LaTeX}
% modernes; elle devrait donc être installée et disponible sur votre
% système. Elle est livrée avec une %
% \doc{http://texdoc.net/pkg/memoir} %
% exhaustive: le guide de l'utilisateur fait près de 600~pages!
% N'hésitez pas à vous y référer pour réaliser une mise en page
% particulière.
%
% L'autre compagnon naturel de la présente documentation est
% \doc[Rédaction avec {\LaTeX}]{http://texdoc.net/pkg/formation-latex-ul} ^^A
% \citep{formation-latex-ul}, la formation {\LaTeX} de l'Université Laval.
% Vous y trouverez des informations additionnelles sur l'utilisation des
% classes \class{ulthese} et \class{memoir}.
%
% \section{Démarrage rapide (pour les impatients)}
% \label{sec:utilisation:rapide}
%
% La classe est livrée avec les distributions {\TeX}~Live, Mac{\TeX} et
% MiK{\TeX}. Si vous utilisez l'une ou l'autre de ces distributions et
% qu'elle est à jour\footnote{%
%   Validez à l'aide de l'assistant de mise à jour de votre
%   distribution.}, %
% vous devriez pouvoir utiliser \class{ulthese} sans autre
% intervention.
%
% Il est recommandé de segmenter tout document d'une certaine ampleur
% dans des fichiers |.tex| distincts pour chaque partie ---
% habituellement un fichier par chapitre. Le document complet est alors
% composé à l'aide d'un fichier maître qui contient le préambule
% {\LaTeX} et un ensemble de commandes \cmd{\include} pour réunir les
% parties dans un tout.
%
% La classe \class{ulthese} est livrée avec un ensemble de gabarits sur
% lesquels se baser pour:
% \begin{itemize}
% \item le document maître de divers types de thèses et de mémoires
%   (standard, par articles, en cotutelle, en bidiplomation);
% \item les fichiers des parties les plus usuelles (résumés français
%   et anglais, avant-propos, introduction, chapitres, conclusion,
%   etc.).
% \end{itemize}
% Les noms des fichiers devraient permettre de facilement identifier
% leur contenu (c'est là une bonne pratique: le nom |rappels.tex|
% parle de lui-même et résiste mieux aux changements à l'ordre des
% chapitres que |chapitre1.tex|).
%
% Dans {\TeX}~Live, les gabarits sont classés avec la documentation de
% la classe. Pour les utiliser, copiez les fichiers appropriés dans votre
% dossier de travail.
%
% Pour débuter la rédaction, renommez le gabarit de document maître
% approprié d'après votre numéro de dossier. Par exemple, l'étudiante
% dont le numéro de dossier est 111555555 et qui entame la rédaction
% d'une thèse par articles renommera le fichier
% \begin{quote}
%   |gabarit-doctorat-articles.tex|
% \end{quote}
% en
% \begin{quote}
%   |111555555.tex|.
% \end{quote}
%
% Les gabarits comportent des commentaires succincts pour vous guider
% dans la préparation de votre document. La \autoref{sec:gabarits}
% fournit les informations détaillées.
%
% \section{Installation}
%
% Cette section explique comment installer la classe \class{ulthese}
% si elle n'est pas disponible sur votre système ou si la version
% n'est pas à jour.
%
% La classe est distribuée sous forme d'une archive |ulthese.zip| via
% le réseau de sites \emph{Comprehensive {\TeX} Archive Network} (CTAN):
% \begin{quote}
%   \url{https://www.ctan.org/pkg/ulthese}
% \end{quote}
%
% L'installation de la classe consiste à créer le fichier
% |ulthese.cls| et les gabarits |.tex| à partir du code source
% documenté se trouvant dans le fichier |ulthese.dtx|. Il est
% recommandé de simplement créer ces fichiers dans le dossier de
% travail de la thèse ou du mémoire.
%
% Pour procéder à l'installation, décompressez l'archive |ulthese.zip|
% dans votre dossier de travail, puis compilez avec {\LaTeX} le fichier
% |ulthese.ins| en exécutant la commande suivante depuis une invite
% de commande:
% \begin{quote}
%   |latex ulthese.ins|
% \end{quote}
% Vous pouvez aussi ouvrir le fichier |ulthese.ins| dans votre éditeur
% de texte favori et lancer depuis celui-ci la compilation avec
% {\LaTeX}, pdf{\LaTeX}, {\XeLaTeX} ou un autre moteur {\TeX}.
%
% \changes{5.3a}{2021-08-09}{Ajout à la documentation de la
%   disponibilité d'une archive pour les personnes qui ne disposent
%   pas d'une installation de {\TeX} locale.}
% \begin{rem}
%   Si vous ne disposez pas d'une version de {\TeX} sur votre poste de
%   travail --- vraisemblablement parce que vous utilisez une
%   plateforme de rédaction en ligne --- téléchargez
%   l'archive |ulthese-sans-tex-local.zip| depuis le
%   \link{https://gitlab.com/vigou3/ulthese/-/releases}{dépôt du
%   projet}. Celle-ci contient tous les fichiers
%   normalement créés lors de l'installation de la classe.
% \end{rem}
%
% \section{Utilisation}
%
% La classe est compatible avec les moteurs {\LaTeX} traditionnels
% ainsi qu'avec le plus récent moteur {\XeLaTeX}.
%
% On charge la classe avec la commande
% \begin{quote}
% |\documentclass|\oarg{options}|{ulthese}|
% \end{quote}
% Les marges, l'interligne et la numérotation des pages sont adaptées
% aux règles de présentation matérielle de la FESP. Les sections
% suivantes décrivent les options et les commandes définies par la
% classe.
%
% \changes{5.0}{2018-04-24}{La classe ne crée plus de page de titre
%   suite à une nouvelle directive de la FESP.}
% \begin{rem}
%   Depuis la version 5.0, la classe ne crée plus de page de titre
%   par défaut, sa production étant désormais prise en charge par la
%   FESP. Néanmoins, les fonctionnalités de composition d'une page
%   de titre demeurent disponibles dans la classe. Les commandes
%   pertinentes sont regroupées à l'\autoref{sec:pagetitre}.
% \end{rem}
%
% \subsection{Options de la classe}
% \label{sec:utilisation:options}
%
% Cette section passe en revue les \meta{options} que l'on peut
% spécifier au chargement de la classe. Les commandes mentionnées
% ci-dessous font quant à elles l'objet de la
% \autoref{sec:utilisation:commandes}.
%
% \begin{DescribeMacro}{PhD,MSc,MA, ...}
%   Identifie le type de grade; consulter le \autoref{tab:grades} pour
%   la liste complète des options et les grades correspondants. La
%   déclaration d'un type de grade est obligatoire, sauf en cas de
%   l'utilisation de l'option |projet| (voir ci-dessous).
% \end{DescribeMacro}
%
% \begin{table}[t]
%   \centering
%   \caption{Options de la classe pour la déclaration du grade et libellés
%     correspondants}
%   \label{tab:grades}
%   \begin{tabularx}{0.85\linewidth}{lX}
%     \toprule
%     Option & Nom du grade (sigle) \\
%     \midrule
%     |LLD|
%       & Docteur en droit (LL.~D.) \\
%     |DMus|
%       & Docteur en musique (D.~Mus.) \\
%     |DPsy|
%       & Docteur en psychologie (D.~Psy.) \\
%     |DThP|
%       & Docteur en théologie pratique (D.~Th.~P.) \\
%     |PhD|
%       & Philosophi{\ae} doctor (Ph.~D.) \\
%     \addlinespace[6pt]
%     |MATDR|
%       & Maître en aménagement du territoire et développement régional
%       (M.ATDR) \\
%     |MArch|
%       & Maître en architecture
%       (M.~Arch.) \\
%     |MA|
%       & Maître ès arts (M.A.) \\
%     |LLM|
%       & Maître en droit (LL.~M.) \\
%     |MErg|
%       & Maître en ergothérapie (M.~Erg.) \\
%     |MMus|
%       & Maître en musique (M.~Mus.) \\
%     |MPht|
%       & Maître en physiothérapie (M.~Pht.) \\
%     |MSc|
%       & Maître ès sciences (M.~Sc.) \\
%     |MScGeogr|
%       & Maître en sciences géographiques (M.~Sc.~géogr.) \\
%     |MServSoc|
%       & Maître en service social (M.~Serv.~soc.) \\
%     |MPsEd|
%       & Maître en psychoéducation (M.~Ps.~éd.) \\
%     \bottomrule
%   \end{tabularx}
% \end{table}
%
% \begin{DescribeMacro}{cotutelle}
%   Identifie une thèse effectuée en cotutelle avec une autre
%   université. Valide uniquement avec un grade de doctorat.
% \end{DescribeMacro}
%
% \begin{DescribeMacro}{bidiplomation}
%   Identifie un mémoire effectué en bidiplomation avec une autre
%   université. Valide uniquement avec un grade de maîtrise.
% \end{DescribeMacro}
%
% \begin{DescribeMacro}{examen}
%   Identifie un examen de doctorat. Cette option permet d'utiliser la
%   classe pour la rédaction d'un examen de doctorat respectant les
%   règles de présentation matérielle de la FESP. Elle a pour effet de
%   changer l'appellation «Thèse» sur la page de titre pour «Examen
%   de doctorat». Elle supprime également la page frontispice.
%   L'option n'est compatible qu'avec l'une des options de grade de
%   doctorat.
% \end{DescribeMacro}
%
% \changes{5.0}{2018-05-09}{Ajout d'une option «essai» pour les
%   maîtrises analogue à «examen» pour les doctorats.}
% \changes{5.2}{2019-05-06}{L'option «essai» n'empêche plus
%   l'utilisation d'une page frontispice.}
% \begin{DescribeMacro}{essai}
%   Identifie un essai de maîtrise. Cette option permet d'utiliser la
%   classe pour la rédaction d'un essai de maîtrise respectant les
%   règles de présentation matérielle de la FESP. Elle a pour effet de
%   changer l'appellation «Maîtrise» sur la page de titre pour
%   «Essai». L'option n'est compatible qu'avec l'une des options de
%   grade de maîtrise. Pour un essai de maîtrise, l'utilisation
%   simultanée de l'option |article| de la classe \class{memoir} peut
%   s'avérer tout indiquée.
% \end{DescribeMacro}
%
% \changes{5.2}{2019-05-06}{Ajout d'une option «projet» pour les
%   cours de projet de recherche.}
% \begin{DescribeMacro}{projet}
%   Identifie un projet de recherche. Cette option permet d'utiliser
%   la classe pour la rédaction d'un rapport de projet de recherche
%   respectant les règles de présentation matérielle de la FESP. Elle
%   a pour effet d'inscrire l'appellation «Projet de recherche» sur la
%   page de titre et de supprimer la mention du grade. L'utilisation
%   de cette option rend d'ailleurs non obligatoire la déclaration de
%   grade qui, si elle est présente, est simplement ignorée. Comme
%   pour un essai de maîtrise, l'utilisation simultanée de l'option
%   |article| de la classe \class{memoir} peut s'avérer tout indiquée
%   avec les cette option.
% \end{DescribeMacro}
%
% \changes{5.3}{2019-11-29}{Ajout d'une option «stage» pour les
%   rapports de stage.}
% \begin{DescribeMacro}{stage}
%   Identifie un rapport de stage. L'effet de cette option est
%   identique à celui de l'option |projet| à deux exceptions près:
%   l'inscription sur la page de titre est «Rapport de stage» et la
%   page frontispice est supprimée.
% \end{DescribeMacro}
%
% \begin{DescribeMacro}{10pt,11pt,12pt}
%   Sélectionne une taille de police de 10, 11 ou 12~points. Par
%   défaut la classe utilise une police de 11~points. Ces options
%   n'ont aucun effet sur la taille des polices des pages de titre.
% \end{DescribeMacro}
%
% \changes{4.4}{2017-05-15}{Ajout de l'option bibchapitre.}
% \changes{4.4}{2017-05-15}{Ajout de l'option bibsection.}
% \begin{DescribeMacro}{bibchapitre}
%   \begin{DescribeMacro}{bibsection}
%     Permettent de composer une bibliographie distincte après chaque
%     fichier inséré dans le document avec \cmd{\include},
%     habituellement un chapitre; voir la
%     \autoref{sec:bibliographie:multiples}. Ces options sont
%     principalement utiles pour les thèses et mémoires par articles.
%     La bibliographie se présente sous la forme d'un chapitre non
%     numéroté avec |bibchapitre|, et sous la forme d'une section
%     numérotée avec |bibsection|.
%   \end{DescribeMacro}
% \end{DescribeMacro}
%
% \begin{DescribeMacro}{nonatbib}
%   Empêche le chargement du paquetage \pkg{natbib}. La paquetage est
%   normalement chargé par la classe; voir la
%   \autoref{sec:bibliographie}. L'option |nonatbib| permet d'empêcher
%   le chargement pour modifier les options du paquetage ou en cas de
%   conflit avec un autre paquetage de mise en forme de la
%   bibliographie.
% \end{DescribeMacro}
%
% \begin{DescribeMacro}{nobabel}
%   Empêche le chargement du paquetage \pkg{babel}. La classe utilise
%   par défaut ce paquetage pour le traitement des langues dans le
%   document; voir la \autoref{sec:langues}. L'option |nobabel| permet
%   d'empêcher son chargement si un autre paquetage devait être
%   utilisé --- on pense ici principalement à \pkg{polyglossia} pour
%   un document produit avec le moteur {\XeLaTeX}.
% \end{DescribeMacro}
%
% \changes{5.1}{2018-09-30}{Ajout de l'option nohyperref.}
% \begin{DescribeMacro}{nohyperref}
%   Empêche le chargement du paquetage \pkg{hyperref}. L'interaction
%   de ce paquetage avec les autres est parfois --- voire souvent ---
%   délicate. La classe se charge de le charger en tout dernier tel
%   que généralement recommandé. L'option |nohyperref| permet
%   d'empêcher son chargement s'il est absolument nécessaire de
%   charger d'autres paquetages avant \pkg{hyperref}. Si vous utilisez
%   cette option, vous devez obligatoirement charger \pkg{hyperref}
%   dans le préambule du document.
% \end{DescribeMacro}
%
% \begin{DescribeMacro}{english, french, ...}
%   Déclare les langues utilisées dans le document. Ces options sont
%   transférées au paquetage \pkg{babel} (dans la mesure où |nobabel|
%   n'est pas spécifié, bien entendu). Le libellé des langues devrait
%   donc correspondre aux options de \pkg{babel}. La dernière langue
%   spécifiée est la langue active par défaut dans le document.
% \end{DescribeMacro}
%
% Toute autre option sera passée à la classe \class{memoir} dont,
% entre autres, le format du papier. Le format lettre nord-américain
% (option |letterpaper|) est utilisé par défaut. Si votre thèse doit
% être imprimée en format international A4, utilisez l'option
% |a4paper|. La classe \class{memoir} est toujours chargée avec
% l'option |oneside|.
%
% \subsection{Commandes de la classe}
% \label{sec:utilisation:commandes}
%
% La classe \class{ulthese} définit quelques nouvelles commandes
% servant principalement à créer la page frontispice et des éléments
% des pages liminaires. Le sommaire des commandes se trouve dans le
% \autoref{tab:commandes} et les descriptions détaillées, ci-dessous.
%
% \begin{table}[t]
%   \centering
%   \caption{Sommaire des commandes de la classe \class{ulthese}.
%     Celles marquées d'une étoile $^\star$ sont obligatoires.}
%   \label{tab:commandes}
%   \begin{tabularx}{\linewidth}{lX}
%     \toprule
%     Commande & Usage \\
%     \midrule
%     \cmd{\titre}$^\star$ & titre principal du document \\
%     \cmd{\soustitre} & sous-titre du document \\
%     \cmd{\auteur}$^\star$ & nom complet de l'auteur \\
%     \cmd{\programme}$^\star$ & nom officiel du programme d'études \\
%     \cmd{\direction}$^\star$ & nom du directeur (directrice) de recherche \\
%     \cmd{\codirection} & noms des codirecteurs de recherche \\
%     \cmd{\frontispice}$^\star$ & production de la page frontispice \\
%     \cmd{\dedicace} & dédicace du document \\
%     \cmd{\epigraphe} & épigraphe du document \\
%     \bottomrule
%   \end{tabularx}
% \end{table}
%
% \begin{DescribeMacro}{\titre}
%   Titre principal de la thèse ou du mémoire. Ne pas utiliser la
%   commande \cmd{\title} de {\LaTeX} pour ce faire.
%
%   Vous devez couper manuellement avec |\\| ou \cmd{\newline} un
%   titre très long. Par exemple, la déclaration d'un titre d'une
%   seule ligne est:
%   \begin{quote}
%     |\titre{Ceci est un titre d'une seule ligne}|
%   \end{quote}
%   Pour un titre de deux lignes, ce sera plutôt:
%   \begin{quote}
%     |\titre{Ceci est la première ligne d'un long titre \\| \\
%     |       et ceci est la seconde}|
%   \end{quote}
% \end{DescribeMacro}
%
% \begin{DescribeMacro}{\soustitre}
%   Sous-titre de la thèse ou du mémoire, le cas échéant. Les remarques
%   sur un long titre principal s'appliquent également au sous-titre.
% \end{DescribeMacro}
%
% \begin{DescribeMacro}{\auteur}
%   Nom complet de l'auteur de la thèse ou du mémoire, sous la forme
%   |Prénom Nom| avec seulement des majuscules initiales. Ne pas
%   utiliser la commande \cmd{\author} de {\LaTeX} pour le nom de
%   l'auteur.
% \end{DescribeMacro}
%
% \begin{DescribeMacro}{\programme}
%   Nom complet officiel du programme d'études comme «Doctorat en
%   informatique» ou «Maîtrise en mathématiques». Si le programme
%   comporte une majeure, séparer sa mention de celle du programme
%   principal par un tiret demi-quadratin (obtenu avec |--|).
% \end{DescribeMacro}
%
% \begin{DescribeMacro}{\direction}
%   Nom complet du directeur ou de la directrice de recherche, sous la
%   forme |Prénom Nom| avec seulement des majuscules initiales, suivi
%   d'une virgule et de la mention «directeur de recherche» ou
%   «directrice de recherche».
%
%   Les thèses en cotutelle comportent un directeur ou une directrice
%   de recherche et un directeur ou une directrice de cotutelle.
%   Séparer chaque mention par |\\|, comme ceci:
%   \begin{quote}
%     |\direction{Prénom Nom, directrice de recherche \\| \\
%     |           Prénom Nom, directeur de cotutelle}|
%   \end{quote}
%
%   De même, les maîtrises en bidiplomation comptent deux directeurs
%   ou directrices de recherche. Séparer chaque mention tel que
%   mentionné ci-dessus.
% \end{DescribeMacro}
%
% \begin{DescribeMacro}{\codirection}
%   Comme la commande \cmd{\direction}, mais pour le ou les
%   codirecteurs de recherche, s'il y a lieu. Lorsqu'il y a plus d'un
%   codirecteur de recherche, séparer chaque mention par |\\|, comme
%   ceci:
%   \begin{quote}
%     |\codirection{Prénom Nom, directrice de recherche \\| \\
%     |             Prénom Nom, directeur de recherche}|
%   \end{quote}
% \end{DescribeMacro}
%
% \begin{DescribeMacro}{\frontispice}
%   Création de la page frontispice. De toutes les commandes
%   ci-dessus, c'est la seule qui doit se trouver dans le corps du
%   document plutôt que dans le préambule.
% \end{DescribeMacro}
%
% \begin{DescribeMacro}{\dedicace}
%   Ajout d'une dédicace («À mes parents», «À Camille») à la thèse ou
%   au mémoire. La dédicace est disposée seule sur une page liminaire,
%   à une dizaine de lignes de la marge du haut et alignée à droite.
%   Par défaut, elle est composée en italique.
% \end{DescribeMacro}
%
% \begin{DescribeMacro}{\epigraphe}
%   Ajout d'une épigraphe au début du document. Comme la dédicace,
%   l'épigraphe est disposée seule sur une page liminaire, à une
%   dizaine de lignes de la marge du haut et alignée à droite. La
%   commande accepte deux arguments, soit le texte de la citation et
%   son auteur ou la source, dans l'ordre.
%
%   Pour ajouter une épigraphe au début d'un ou plusieurs chapitres,
%   utiliser directement la commande \cmd{\epigraph} de
%   \class{memoir}, sur laquelle \cmd{\dedicace} et \cmd{\epigraphe}
%   sont d'ailleurs basées.
% \end{DescribeMacro}
%
% \subsection{Citations}
% \label{sec:citations}
%
% {\LaTeX} offre deux environnements pour les citations dans le
% texte: |quote| et |quotation|.
%
% \begin{DescribeEnv}{quote}
%   L'environnement |quote| sert pour les citations «courtes»,
%   quelques lignes au plus. Dans la classe, le texte est alors placé
%   en retrait des marges normales de 10~mm à gauche et à droite.
% \end{DescribeEnv}
%
% \begin{DescribeEnv}{quotation}
%   L'environnement |quotation|, quant à lui, doit être utilisé pour
%   les citations «longues», celles qui peuvent s'étendre sur plus de
%   cinq lignes ou, surtout, plus d'un paragraphe. Dans la classe, le
%   texte est alors toujours placé en retrait de 10~mm, mais également
%   à interligne simple. De plus, un espace vertical sépare les
%   paragraphes, le cas échéant, afin de bien les distinguer les uns
%   des autres.
% \end{DescribeEnv}
%
% \subsection{Interligne}
% \label{sec:interligne}
%
% \begin{DescribeMacro}{\OnehalfSpacing}
%   L'espacement d'un interligne et demi utilisé dans la classe est
%   obtenu avec la commande \cmd{\OnehalfSpacing} de \class{memoir}.
%   L'interligne simple est automatiquement rétabli pour les pages
%   de titre, la table des matières, la liste des tableaux, la liste
%   des figures et les longues citations (\autoref{sec:citations}).
% \end{DescribeMacro}
%
% \begin{DescribeMacro}{\SingleSpacing}
%   Si ce devait être nécessaire ailleurs dans le document, la
%   commande \cmd{\SingleSpacing} permet de passer à l'interligne simple.
% \end{DescribeMacro}
%
% \subsection{Hyperliens}
% \label{sec:hyperliens}
%
% Lorsque le document est produit avec pdf{\LaTeX} ou {\XeLaTeX},
% toutes les références internes et externes sont automatiquement
% transformées en hyperliens cliquables. Ceux-ci apparaissent dans une
% teinte de bleu assez foncée tout à la fois visible en couleur et peu
% contrastante si le document est imprimé en noir et blanc.
%
% \begin{DescribeMacro}{\hypersetup}
%   Vous pouvez ajouter des options de configuration pour les
%   hyperliens ou les métadonnées du document PDF avec la commande
%   \cmd{\hypersetup} du paquetage \class{hyperref}.
% \end{DescribeMacro}
%
% \subsection{Autres paquetages chargés}
% \label{sec:paquetages}
%
% Outre \class{memoir}, la classe \class{ulthese} charge quelques
% paquetages qui peuvent aussi vous être utiles. Il n'est donc pas
% nécessaire de charger de nouveau les paquetages suivants:
% \begin{description}
% \item[\normalfont\pkg{babel}] \citep{babel} gestion des documents
%   rédigés dans une ou plusieurs langues autres que l'anglais (si
%   l'option |nobabel| de la classe est absente; voir aussi la
%   \autoref{sec:langues});
% \item[\normalfont\pkg{numprint}] \citep{numprint} requis par la
%   commande \cmd{\nombre} de \pkg{babel}; le paquetage est donc
%   chargé uniquement si \pkg{babel} l'est. Permet de composer
%   automatiquement des nombres avec un séparateur toutes les trois
%   positions (une espace en français);
% \item[\normalfont\pkg{natbib}] \citep{natbib} gestion de la
%   bibliographie (si l'option |nonatbib| de la classe est absente;
%   voir aussi la \autoref{sec:bibliographie});
% \item[\normalfont\pkg{chapterbib}] \citep{chapterbib} gestion de
%   bibliographies multiples (si l'une ou l'autre des options
%   |bibchapitre| ou |bibsection| est spécifiée;
%   voir aussi la \autoref{sec:bibliographie:multiples});
% \item[\normalfont\pkg{hyperref}] \citep{hyperref} permet de
%   transformer toutes les références en hyperliens cliquables lorsque
%   le document est produit avec pdf{\LaTeX} ou {\XeLaTeX} (si
%   l'option |nohyperref| est spécifiée, vous devez obligatoirement
%   charger le paquetage vous-même);
% \item[\normalfont\pkg{etoolbox}] \citep{etoolbox} outils
%   additionnels pour la conception de classes et de paquetages
%   (utilisé pour le chargement de \pkg{hyperref});
% \item[\normalfont\pkg{fontspec}] \citep{fontspec} gestion des
%   polices OpenType sous {\XeLaTeX} (chargé avec ce moteur
%   seulement);
% \item[\normalfont\pkg{graphicx}] \citep{graphicx} insertion et
%   manipulation de graphiques;
% \item[\normalfont\pkg{xcolor}] \citep{xcolor} gestion des couleurs
%   dans le document, notamment celle des hyperliens;
% \item[\normalfont\pkg{textcomp}] multitude de symboles spéciaux,
%   dont un beau symbole de copyright, \textcopyright.
% \end{description}
% L'\autoref{sec:meo} sur la mise en œuvre de la classe fournit
% plus de détails sur la liste des paquetages chargés et les raisons
% pour lesquelles ils sont requis dans la classe.
%
% \subsection{Paquetage incompatible}
% \label{sec:incompatible}
%
% Le paquetage \pkg{geometry} est incompatible avec la classe à cause
% de sa mauvaise interaction avec \class{memoir}. Son chargement dans
% le préambule du document cause une erreur lors de la compilation.
%
% \section{Thèses et mémoires par articles}
% \changes{5.0}{2018-06-08}{Ajout d'une section spécifique sur les
%   thèses et mémoires par articles.}
%
% La rédaction de la thèse ou du mémoire par articles est une pratique
% qui gagne en popularité. Elle consiste à remplacer le corps
% principal du manuscrit par un ou des articles scientifiques. Chaque
% article fait normalement l'objet d'un chapitre dans le document. Les
% résumés en français (obligatoire) et dans la langue de la thèse ou
% du mémoire (le cas échéant) y apparaissent sous forme de
% sections numérotées.
%
% L'autre caractéristique qui distingue la thèse et le mémoire par
% articles des versions standards, c'est la nécessité d'utiliser une
% bibliographie distincte pour chacun des chapitres. Les options de la
% classe |bibchapitre| et |bibsection| permettent de réaliser cette
% mise en forme particulière. La \autoref{sec:bibliographie:multiples}
% explique comment en tirer profit.
%
% Des gabarits de document maître et de chapitre fournissent la
% structure de base pour la thèse et le mémoire par articles. Vous
% pouvez, au besoin, combiner les caractéristiques de ces gabarits à
% celles de la thèse en cotutelle ou du mémoire en bidiplomation. La
% \autoref{sec:gabarits} fournit les détails sur les gabarits.
%
% \section{Français et autres langues}
% \label{sec:langues}
%
% Une complication additionnelle si vous rédigez dans une langue autre
% que l'anglais consiste à adapter {\LaTeX} à votre langue, qu'il
% s'agisse des mots clés, de la typographie ou de la césure des mots.
% La solution standard à ce problème provient du paquetage
% \pkg{babel}. Celui-ci permet de combiner plusieurs langues dans un
% même document et de passer de l'une à l'autre facilement. Il est
% chargé par défaut par la classe \class{ulthese}.
%
% Aucune langue n'est spécifiée dans la classe. Vous aurez sans doute
% recours à l'anglais et au français, ne serait-ce que pour les deux
% résumés demandés par la FESP. Vous devez spécifier les langues
% utilisées dans le document en options à la classe, tel que mentionné
% à la \autoref{sec:utilisation:options}. La \emph{dernière} langue
% spécifiée devient par défaut la langue active du document.
%
% \begin{DescribeMacro}{\selectlanguage}
%   La commande \cmd{\selectlanguage} de \pkg{babel} permet de passer de
%   la langue courante à la langue spécifiée en argument.
% \end{DescribeMacro}
%
% \begin{DescribeEnv}{otherlanguage}
%   L'environnement |otherlanguage| de \pkg{babel} permet de faire
%   la même chose que la commande \cmd{\selectlanguage}, sauf que le
%   changement de langue est local à l'environnement --- utile pour
%   les brefs changements de langue.
% \end{DescribeEnv}
%
% Si vous n'êtes pas autrement familier avec le paquetage
% \pkg{babel}, consultez sa %
% \doc{http://texdoc.net/pkg/babel/}. %
% Celle-ci est éclatée en un
% document principal, pour le cœur du paquetage et plusieurs
% autres pour les fonctionnalités propres à une langue: %
% \doc[anglais]{http://texdoc.net/pkg/babel-english}, %
% \doc[français]{http://texdoc.net/pkg/babel-french}, %
% etc. Consultez au moins les documents consacrés aux
% langues utilisées dans votre thèse ou mémoire. Le plus simple
% consiste sans doute à consulter en ligne sur CTAN les
% \link{https://mirrors.ctan.org/tex-archive/macros/latex/required/babel/contrib/}{%
% documents spécifiques par langue}.
% \changes{4.4}{2017-05-14}{Correction d'une url vers la documentation
%   de babel.}
%
% \begin{DescribeMacro}{\nombre}
%   Le paquetage \pkg{numprint} étant chargé dans la classe avec
%   \pkg{babel}, vous pouvez utiliser la commande \cmd{\nombre} pour
%   formater automatiquement les nombres. Par exemple, le résultat de
%   |\nombre{123456789}| est \nombre{123456789}.
% \end{DescribeMacro}
%
% Les utilisateurs de {\XeLaTeX} qui souhaiteraient plutôt utiliser le
% plus récent paquetage \pkg{polyglossia} \citep{polyglossia} peuvent
% empêcher le chargement de \pkg{babel} avec l'option |nobabel| de la
% classe. Ils devront toutefois charger et configurer
% \pkg{polyglossia} eux-mêmes dans l'entête de leur document. Ce
% paquetage est moins évolué que \pkg{babel} pour la typographie
% française.
%
% \section{Bibliographie}
% \label{sec:bibliographie}
%
% \begin{DescribeMacro}{\bibliography}
%   Tel qu'expliqué au chapitre~8 de \citet{formation-latex-ul},
%   il est fortement recommandé d'utiliser {\BibTeX} pour la
%   préparation de la bibliographie d'un document. Celle-ci est
%   insérée dans le document à l'endroit où apparait la commande
%   \cmd{\bibliography} dans le code source. Cette commande prend en
%   arguments les noms des bases de données bibliographiques séparés
%   par des virgules.
% \end{DescribeMacro}
%
% \subsection{Mise en forme des citations}
% \label{sec:bibliographie:options}
%
% Dans la classe \class{ulthese}, la mise en forme des citations est
% confiée au paquetage \pkg{natbib} (à moins que l'option |nonatbib|
% ne soit spécifiée). Le paquetage est chargé avec les options par
% défaut, soit |round|, |semicolon| et |authoryear|. Pour spécifier
% d'autres options, vous avez deux possibilités:
% \begin{enumerate}
% \item utiliser l'option |nonatbib| de la classe et ensuite charger
%   explicitement \pkg{natbib} avec ses options;
% \item
%   \begin{DescribeMacro}{\setcitestyle}
%     utiliser la commande \cmd{\setcitestyle} pour passer de
%     nouvelles options à \pkg{natbib}.
%   \end{DescribeMacro}
% \end{enumerate}
%
% Par exemple, pour utiliser un style de citation numérique où le
% numéro de la référence se trouve entre crochets, vous pouvez
% procéder de l'une ou l'autre des deux manières suivantes:
% \begin{quote}
%   |\documentclass[nonatbib]{ulthese}| \\
%   |\usepackage[numbers,square]{natbib}| \\
%   |...| \\
%   |\bibliographystyle{plain-fr}|
% \end{quote}
% ou
% \begin{quote}
%   |\documentclass{ulthese}| \\
%   |...| \\
%   |\setcitestyle{numbers,square}| \\
%   |\bibliographystyle{plain-fr}|
% \end{quote}
%
% Consultez la %
% \doc{http://texdoc.net/pkg/natbib/} %
% de \pkg{natbib} pour les détails.
%
% \subsection{Bibliographies multiples}
% \label{sec:bibliographie:multiples}
%
% L'ajout de |bibchapitre| ou de |bibsection| dans les options de la
% classe permet de générer une bibliographie distincte pour chaque
% chapitre, une chose particulièrement utile pour les thèses et
% mémoire par articles.
%
% La classe fait appel au paquetage \pkg{chapterbib}
% \citep{chapterbib} pour composer des bibliographies multiples. La
% \doc{http://texdoc.net/pkg/chapterbib/} du paquetage
% explique la procédure à suivre pour obtenir une bibliographie par
% chapitre avec {\BibTeX}. En résumé:
% \begin{enumerate}
% \item chaque chapitre comportant sa propre bibliographie doit
%   absolument être inséré dans le document principal avec la commande
%   \cmd{\include} \citep[section~3.5]{formation-latex-ul};
% \item chaque fichier visé doit contenir une commande
%   \cmd{\bibliographystyle} (\autoref{sec:bibliographie:style}) et
%   une commande \cmd{\bibliography};
% \item on obtient les diverses bibliographies en compilant avec
%   {\BibTeX} les fichiers individuels, et non le document maître.
% \end{enumerate}
%
% La plupart des auteurs préféreront sans doute l'effet de l'option
% |bibsection|, où la bibliographie apparait comme une section normale
% à la fin d'un chapitre. Lorsque \pkg{babel} est chargé, la section
% sera numérotée. Pour supprimer la numérotation, insérer dans le
% préambule du document la commande
% \begin{quote}
%   |\addto\extras|\meta{langue}|{\renewcommand{\bibsection}{%| \\
%   |  \section*{\bibname}\prebibhook}}|
% \end{quote}
% où \meta{langue} est la langue par défaut du document
% (\autoref{sec:langues}).
%
% \subsection{Style de la bibliographie}
% \label{sec:bibliographie:style}
%
% \begin{DescribeMacro}{\bibliographystyle}
%   Le format général de la bibliographie est contrôlé par un
%   \emph{style} choisi avec la commande \cmd{\bibliographystyle} dans
%   le préambule du document ou dans les fichiers de chaque chapitre
%   dans une thèse ou un mémoire par articles. Les styles standards de
%   {\LaTeX} sont |plain|, |unsrt|, |alpha| et |abbrv|.
% \end{DescribeMacro}
%
% Le paquetage \pkg{natbib} chargé par défaut par la classe supporte
% le style de citation auteur-année fréquemment employé en sciences
% naturelles, plusieurs commandes de citation, un grand nombre de
% styles de bibliographie ainsi que des entrées spécifiques pour les
% numéros ISBN et les URL. Le paquetage fournit des styles de
% bibliographie |plainnat|, |unsrtnat| et |abbrvnat| similaires aux
% styles standards, mais plus complets. Il existe des
% \link{https://mirrors.ctan.org/biblio/bibtex/contrib/bib-fr/}{versions
% francisées} de ces styles (et de quelques autres) dans CTAN.
%
% Le paquetage \pkg{francais-bst} \citep{francais-bst} fournit une
% feuille de style compatible avec \pkg{natbib} permettant de composer
% des bibliographies auteur-année respectant les normes de typographie
% française proposées dans \cite{Malo:1996}. Pour utiliser ce style,
% spécifier dans le préambule du document LaTeX
% \begin{quote}
%   |\bibliographystyle{francais}|
% \end{quote}
%
% Autrement, la FESP n'a pas d'exigences particulières quant à la
% présentation de la bibliographie (présentation du titre, des auteurs
% et autres informations bibliographiques).
%
% \section{Police de caractères du document}
% \label{sec:police}
%
% Les documents {\LaTeX} sont facilement reconnaissables par leur
% police de caractères par défaut, {\fontfamily{cmr}\selectfont
% Computer Modern}. Avec toute distribution {\LaTeX} moderne, il est
% maintenant simple d'utiliser l'une ou l'autre des polices PostScript
% standards. D'ailleurs la classe \class{ulthese} utilise la police
% sans empattements \textsf{Helvetica} pour composer les pages de
% titre.
%
% La FESP permet l'utilisation des polices Times, Palatino (la police
% du présent document) et Lucida~Bright dans les thèse et mémoires.
% \citet[section~10.2]{formation-latex-ul} explique comment utiliser ces
% polices dans votre document.
%
% \section{Gabarits}
% \label{sec:gabarits}
%
% Afin de vous épauler dans la préparation de votre manuscrit, la
% classe \class{ulthese} est livrée avec un ensemble de gabarits. Tout
% d'abord, pour le document maître de différents types de thèses et de
% mémoires:
% \begin{itemize}[noitemsep]
% \item |gabarit-doctorat.tex|;
% \item |gabarit-doctorat-articles.tex|;
% \item |gabarit-doctorat-cotutelle.tex|;
% \item |gabarit-maitrise.tex|;
% \item |gabarit-maitrise-articles.tex|;
% \item |gabarit-maitrise-bidiplomation.tex|.
% \end{itemize}
% Ensuite, pour quelques parties d'un document:
% \begin{itemize}[noitemsep]
% \item |resume.tex|;
% \item |abstract.tex|;
% \item |remerciements.tex|;
% \item |avantpropos.tex|;
% \item |introduction.tex|;
% \item |chapitre1.tex|;
% \item |chapitre1-articles.tex|;
% \item |chapitre2.tex|;
% \item |chapitre2-articles.tex|;
% \item |conclusion.tex|;
% \item |annexe.tex|.
% \end{itemize}
% Vous pourriez devoir combiner les caractéristiques de deux gabarits
% de document maître comme, par exemple, pour une thèse en cotutelle
% par articles.
%
% Les gabarits comportent des commentaires succincts pour vous guider
% dans leur utilisation. Les sections suivantes fournissent des
% détails additionnels, et ce, dans l'ordre où les commandes
% apparaissent dans les fichiers.
%
% \begin{rem}
%   Il n'y a pas de gabarit spécifique pour un examen de doctorat,
%   pour un essai de maîtrise, pour un projet de recherche ou pour un
%   rapport de stage. Utilisez un gabarit de thèse ou de maîtrise et
%   inscrivez l'option |examen|, |essai|, |projet| ou |stage| dans la
%   commande \cmd{\documentclass}.
% \end{rem}
%
% \subsection{Encodage des fichiers}
%
% Composer en {\LaTeX} de longs textes dans une langue ayant recours
% aux signes diacritiques devient rapidement pénible si l'on utilise
% des commandes telles que |\'e|, |\`a| ou |\"o| pour entrer des lettres
% accentuées. Afin de pouvoir plutôt entrer directement |é|, |à| ou
% |ö|, {\LaTeX} doit être configuré pour reconnaître les lettres
% accentuées. C'est le rôle du paquetage \pkg{inputenc}
% \citep{inputenc}.
%
% Il existe plusieurs manières différentes d'encoder --- ou
% d'enregistrer --- les lettres accentuées et autres caractères
% spéciaux (comme, par exemple, le symbole de l'euro) dans un
% ordinateur. La méthode la plus répandue et celle standard sur les
% versions récentes des systèmes d'exploitation Linux et macOS est
% l'UTF-8 de la norme
% \link{https://fr.wikipedia.org/wiki/Unicode}{Unicode}. Les gabarits
% sont livrés dans ce type d'encodage.
%
% La déclaration
% \begin{quote}
%   |\usepackage[utf8]{inputenc}|
% \end{quote}
% dans le préambule assure que {\LaTeX} traitera correctement des
% fichiers source encodés en UTF-8.
%
% La norme Unicode n'est pas aussi uniformément supportée par Windows.
% Selon l'éditeur de texte employé et la version du système
% d'exploitation, il peut être nécessaire d'utiliser les normes
% d'encodage %
% \link{https://fr.wikipedia.org/wiki/ISO_8859-1}{ISO~8859-1} %
% (ou Latin-1; option |latin1| de \pkg{inputenc}), %
% \link{https://fr.wikipedia.org/wiki/ISO_8859-15}{ISO~8859-15} %
% (ou Latin-9; option |latin9|) ou %
% \link{https://fr.wikipedia.org/wiki/Windows-1252}{Windows-1252} %
% (options |cp1252| ou |ansinew|).
%
% La situation est plus simple avec {\XeLaTeX} puisqu'il gère
% nativement Unicode. Le paquetage \pkg{inputenc} est non seulement
% inutile, mais incompatible avec {\XeLaTeX}. C'est pourquoi, dans les
% gabarits, \pkg{inputenc} est chargé seulement lorsque {\XeTeX}
% n'est pas le moteur employé pour compiler le document.
%
% \subsection{Paquetages additionnels}
%
% Tel qu'expliqué à la \autoref{sec:paquetages}, la classe charge déjà
% quelques paquetages. Cependant, il est fort probable que vous devrez
% en charger d'autres pour composer votre document. Les gabarits
% prévoient un endroit pour le chargement de paquetages additionnels.
%
% Si vous utilisez un paquetage non standard dans les distributions
% courantes (\TeX~Live, Mac\TeX, MiK\TeX), vous devez le fournir avec
% le code source de votre document lors du dépôt final.
%
% \subsection{Changement de police de caractères}
%
% Les gabarits comportent des déclarations types pour utiliser les
% polices Palatino ou Times sous {\LaTeX} ou, sous {\XeLaTeX}, leurs
% équivalents Pagella et Termes du projet
% \link{https://www.gust.org.pl/projects/e-foundry/tex-gyre/}{TeX~Gyre}.
%
% \subsection{Options de \pkg{babel}}
%
% \begin{DescribeMacro}{\frenchbsetup}
%   La commande \cmd{\frenchbsetup} de \pkg{babel} permet de contrôler
%   certains ajustements typographiques apportés par le paquetage en
%   mode français. Consultez la documentation de \pkg{babel} pour la
%   liste des options de configuration disponibles.
% \end{DescribeMacro}
%
% Les concepteurs de la classe \class{ulthese}  proposent trois
% ajustements dans les gabarits:
% \begin{enumerate}
% \item l'option |StandardItemizeEnv=true| évite que le mode français de
%   \pkg{babel} ne diminue l'espacement vertical dans les listes;
% \item l'option |ThinSpaceInFrenchNumbers=true| fait en sorte qu'une espace fine
%   sera utilisée comme séparateur des milliers dans les nombres plutôt
%   qu'une espace pleine;
% \item les options |og=«| et |fg=»| déclarent que les caractères « et
%   » utilisés dans le code source représentent les guillemets ouvrant
%   et fermant, respectivement. Cela évite de devoir utiliser les
%   commandes \cmd{\og} et \cmd{\fg} de \pkg{babel} tout en bénéficiant de
%   l'ajutement automatique des espaces autour des symboles.
% \end{enumerate}
% Vous devez évidemment désactiver ces ajustements si l'option
% |nobabel| est spécifiée au chargement de la classe.
%
% \subsection{Style de la bibliographie}
% \label{sec:bibliographystyle}
%
% Tel qu'expliqué à la \autoref{sec:bibliographie:style}, la commande
% \cmd{\bibliographystyle} permet de définir le mode de citation dans
% le texte et la présentation des notices bibliographiques. Cette
% commande devrait figurer dans les fichiers des chapitres
% individuels dans une thèse ou un mémoire par articles.
%
% \subsection{Déclarations de la page frontispice}
%
% Les gabarits comportent toutes les déclarations nécessaires pour
% composer la page frontispice des divers types de thèse ou de
% mémoires. Vous devez remplacer les éléments se trouvant entre
% crochets |< >| en respectant la forme indiquée. Assurez-vous de
% supprimer les caractères |<| et |>| afin qu'ils n'apparaissent pas
% sur la page frontispice de votre document.
%
% \subsection{Pages liminaires}
%
% \begin{DescribeMacro}{\frontmatter}
%   La commande \cmd{\frontmatter} déclare que {\LaTeX} doit considérer le
%   matériel qui suit comme des pages liminaires. En pratique, cela
%   résulte essentiellement en une numérotation des pages en chiffres
%   romains.
% \end{DescribeMacro}
%
% Les normes de présentation de la FESP édictent que les thèses et
% mémoires devraient comporter les pages liminaires suivantes, dans
% l'ordre:
% \begin{enumerate}[noitemsep]
% \item page frontispice (obligatoire);
% \item résumé en français (obligatoire);
% \item résumé en anglais (recommandé, mais non obligatoire);
% \item table des matières (obligatoire);
% \item liste des tableaux;
% \item liste des figures;
% \item liste des abbréviations et des sigles;
% \item dédicace;
% \item épigraphe;
% \item remerciements;
% \item avant-propos (obligatoire dans le cas d'une thèse ou d'un
%   mémoire par articles).
% \end{enumerate}
%
% Les commandes
% \begin{quote}
%   \cmd{\frontispice} \\
%   \cmd{\tableofcontents} \\
%   \cmd{\listoftables} \\
%   \cmd{\listoffigures} \\
%   \cmd{\dedicace}\marg{texte} \\
%   \cmd{\epigraphe}\marg{texte}\marg{auteur}
% \end{quote}
% permettent de générer les pages correspondantes. Seules les deux
% dernières commandes admettent des arguments.
%
% \begin{rem}
%   L'examen de doctorat, l'essai de maîtrise, le rapport de projet de
%   recherche et le rapport de stage comportent une page de titre.
%   Vous devez ajouter dans les gabarits les commandes nécessaires
%   pour la production de la page de titre (\autoref{sec:pagetitre})
%   et supprimer les autres pages liminaires. Rappelons que l'option
%   de classe |examen| a pour effet de supprimer la page frontispice.
% \end{rem}
%
% \begin{DescribeMacro}{\chapter*}
%   Les résumés, la liste des abbréviations et des sigles, les
%   remerciements et l'avant-propos sont composés comme des chapitres
%   normaux, mais sans être numérotés. Il faut donc définir ces éléments
%   avec la commande \cmd{\chapter*}.
% \end{DescribeMacro}
%
% \begin{DescribeMacro}{\label}
%   La commande \cmd{\label} permet d'utiliser le système de renvois
%   automatiques de {\LaTeX} \citep[section~3.4]{formation-latex-ul}.
%   Les étiquettes des chapitres non numérotés sont déjà configurées
%   pour l'ajout de ces chapitres dans la table des matières (voir
%   ci-dessous). Nous recommandons fortement de définir vos propres
%   étiquettes pour les chapitres normaux.
% \end{DescribeMacro}
%
% \begin{DescribeMacro}{\phantomsection}
%   \begin{DescribeMacro}{\addcontentsline}
%     \begin{DescribeMacro}{\nameref}
%       Les chapitres déclarés avec la commande \cmd{\chapter*}
%       n'apparaissent pas dans la table des matières. Comme les
%       normes de présentation de la FESP exigent que toutes les pages
%       liminaires y figurent, nous faisons suivre les commandes
%       \cmd{\chapter*}\marg{Titre} et \cmd{\label}\marg{nom} des
%       commandes
%     \end{DescribeMacro}
%   \end{DescribeMacro}
% \end{DescribeMacro}
% \begin{quote}
%   |\phantomsection\addcontentsline{toc}{chapter}{%| \\
%   |  \nameref|\marg{nom}|}|
% \end{quote}
% Celles-ci ajoutent à la table des matières (|toc|) une section de
% niveau |chapter| dont le titre est celui du chapitre auquel
% correspond l'étiquette \meta{nom}. La commande \cmd{\phantomsection}
% est rendue nécessaire (ou recommandée) par le paquetage
% \pkg{hyperref}.
%
% \subsection{Corps du document}
%
% \begin{DescribeMacro}{\mainmatter}
%   La commande \cmd{\mainmatter} délimite le début du corps du document. La
%   numérotation des pages passe en chiffres arabes.
% \end{DescribeMacro}
%
% Le corps du document devrait normalement compter une introduction (non
% numérotée), un développement divisé en chapitres (numérotés) et une
% conclusion (non numérotée).
%
% \subsection{Annexes}
%
% \begin{DescribeMacro}{\appendix}
%   Si la thèse ou le mémoire comporte une ou plusieurs annexes,
%   composer celles-ci comme des chapitres normaux insérés dans le
%   document maître après la commande \cmd{\appendix}. Cette commande a
%   pour effet de passer d'un mode de numération numérique (1, 1.1, 2, 2.1,
%   \dots) à un mode alphanumérique (A, A.1, B, B.1, \dots).
% \end{DescribeMacro}
%
% \subsection{Bibliographie}
%
% La bibliographie figure normalement à la toute fin du document, sous
% forme de chapitre non numéroté. Tel qu'expliqué à la
% \autoref{sec:bibliographie:multiples}, la commande
% \cmd{\bibliography} doit plutôt se trouver dans les fichiers des
% chapitres dans une thèse ou un mémoire par articles.
%
% \section{Commentaires, suggestions et aide}
% \changes{5.0}{2018-08-09}{Code source du projet déplacé de GitHub
%   vers GitLab.}
%
% Le \link{https://gitlab.com/vigou3/ulthese/}{dépôt du projet} dans
% GitLab demeure le meilleur endroit pour rapporter des bogues ou pour
% proposer des améliorations à la classe \class{ulthese}.
%
% Pour obtenir de l'aide additionnelle sur l'utilisation de la classe
% --- et non sur celle de {\LaTeX} en général --- consultez d'abord:
% \begin{enumerate}
% \item le \link{http://www.theses.ulaval.ca/wiki/}{WikiThèse} de
%   l'Université Laval, en particulier la
%   \link{http://www.theses.ulaval.ca/wiki/index.php?title=FAQ}{Foire aux questions};
% \item les
%   \link{http://listes.ulaval.ca/listserv/archives/ulthese-aide.html}{archives}
%   de la liste de distribution \texttt{ulthese-aide}.
% \end{enumerate}
% Si la réponse à votre question ne se trouve ni dans le wiki, ni dans
% les archives, alors écrivez à l'adresse
% \href{mailto:ulthese-aide@listes.ulaval.ca}{\url{ulthese-aide@listes.ulaval.ca}}.
%
% \StopEventually{
%   \begin{thebibliography}{15}
%     \addcontentsline{toc}{section}{Références}
%   \bibitem[{Arseneau(2010)}]{chapterbib}
%     Arseneau, D. 2010,
%     {\selectlanguage{english}\emph{chapterbib. Multiple
%     bibliographies in {\LaTeX}}}.
%     URL~\url{https://www.ctan.org/pkg/chapterbib/}.
%
%   \bibitem[{Braams et Bezos(2016)}]{babel}
%     Braams, J. et J.~Bezos. 2016,
%     {\selectlanguage{english}\emph{Babel}}.
%     URL~\url{https://www.ctan.org/pkg/babel/}.
%
%   \bibitem[{Carlisle et {The \LaTeX3\ Project}(2016)}]{graphicx}
%     Carlisle, D. et {The \LaTeX3\ Project}. 2016,
%     {\selectlanguage{english}\emph{Packages in the `graphics'
%     Bundle}}. URL~\url{https://www.ctan.org/pkg/graphics/}.
%
%   \bibitem[{Charette(2015)}]{polyglossia}
%     Charette, F. 2015, {\selectlanguage{english}\emph{Polyglossia:
%     An Alternative to Babel for {\XeLaTeX} and {\LuaLaTeX}}}.
%     URL~\url{https://www.ctan.org/pkg/polyglossia/}, current
%     maintainer Arthur Reutenauer.
%
%   \bibitem[{Daly(2010)}]{natbib}
%     Daly, P.~W. 2010, {\selectlanguage{english}\emph{Natural
%     Sciences Citations and References}}.
%     URL~\url{https://www.ctan.org/pkg/natbib/}.
%
%   \bibitem[{Goulet(2013)}]{francais-bst}
%     Goulet, V. 2013, {\selectlanguage{french}«Paquetage
%     francais-bst»},
%     URL~\url{https://www.ctan.org/pkg/francais-bst/}.
%
%   \bibitem[{Goulet(2020)}]{formation-latex-ul}
%     Goulet, V. 2020, {\selectlanguage{french}\emph{Rédaction avec
%     {\LaTeX}}}, document libre sous contrat Creative Commons.
%     URL~\url{httpss://ctan.org/pkg/formation-latex-ul}.
%
%   \bibitem[{Harders(2012)}]{numprint}
%     Harders, H. 2012, {\selectlanguage{english}\emph{The
%     numprint package}}.
%     URL~\url{https://www.ctan.org/pkg/numprint/}.
%
%   \bibitem[{Jeffrey et Mittelbach(2015)}]{inputenc}
%     Jeffrey, A. et F.~Mittelbach. 2015,
%     {\selectlanguage{english}\emph{inputenc.sty}}.
%     URL~\url{https://www.ctan.org/pkg/inputenc/}.
%
%   \bibitem[{Kern(2016)}]{xcolor}
%     Kern, D.~U. 2016, {\selectlanguage{english}\emph{Extending
%     {\LaTeX}’s color facilities: the xcolor package}}.
%     URL~\url{https://www.ctan.org/pkg/xcolor/}.
%
%   \bibitem[{Lehman et Wright(2018)}]{etoolbox}
%     Lehman, P. et J.~Wright. 2018,
%     {\selectlanguage{english}\emph{The etoolbox package}}.
%     URL~\url{https://www.ctan.org/pkg/etoolbox/}.
%
%   \bibitem[{Malo(1996)}]{Malo:1996}
%     Malo, M. 1996, {\selectlanguage{french}\emph{Guide de la
%     communication écrite au cégep, à l'université et en
%     entreprise}}, Québec Amérique. ISBN~978-2-8903-7875-9.
%
%   \bibitem[{Rahtz et Oberdiek(2017)}]{hyperref}
%     Rahtz, S. et H.~Oberdiek. 2017,
%     {\selectlanguage{english}\emph{Hypertext marks in {\LaTeX}: a
%     manual for hyperref}}.
%     URL~\url{https://www.ctan.org/pkg/hyperref/}.
%
%   \bibitem[{Robertson et Hosny(2017)}]{fontspec}
%     Robertson, W. et K.~Hosny. 2017,
%     {\selectlanguage{english}\emph{The fontspec package: Font
%     selection for {\XeLaTeX} and {\LuaLaTeX}}}.
%     URL~\url{https://www.ctan.org/pkg/fontspec/}.
%
%   \bibitem[{Wilson(2016)}]{memoir}
%     Wilson, P. 2016, {\selectlanguage{english}\emph{The Memoir Class
%     for Configurable Typesetting}}, 8{\ieme} éd., The Herries Press.
%     URL~\url{https://www.ctan.org/pkg/memoir/}, maintained by Lars
%     Madsen.
%   \end{thebibliography}
%   \PrintChanges
% }
%
% \appendix
%
% \section{Création d'une page de titre}
% \label{sec:pagetitre}
%
% La thèse ou le mémoire déposé à la FESP ne doit \emph{pas} comporter
% de page de titre. Les options de classe et les commandes mentionnées
% dans cette section permettent néanmoins de composer une page de
% titre conforme aux normes de présentation matérielles, notamment
% pour les examens de doctorat, les essais de maîtrise ou les rapports
% de projet de recherche ou de stage.
%
% En premier lieu, les \meta{options} suivantes s'ajoutent à celles de
% la \autoref{sec:utilisation:options} lorsque l'on souhaite créer une
% page de titre.
%
% \begin{DescribeMacro}{multifacultaire}
%   Identifie une thèse multifacultaire. Valide uniquement avec un
%   grade de doctorat. Cette déclaration requiert ensuite d'utiliser
%   la commande \cmd{\faculteUL}.
% \end{DescribeMacro}
%
% \begin{DescribeMacro}{extensionUdeS}
%   Identifie une thèse réalisée en extension à l'Université de
%   Sherbrooke. Valide uniquement avec un grade de doctorat. Cette
%   déclaration requiert ensuite d'utiliser les commandes
%   \cmd{\faculteUL} et \cmd{\faculteUdeS}.
% \end{DescribeMacro}
%
% \begin{DescribeMacro}{extensionUQO}
%   Identifie une thèse réalisée en extension à l'Université du
%   Québec en Outaouais (UQO). Valide uniquement avec un grade de doctorat.
%   Cette déclaration requiert ensuite d'utiliser les commandes
%   \cmd{\faculteUL} et \cmd{\faculteUQO}.
% \end{DescribeMacro}
%
% \begin{DescribeMacro}{extensionUQAC}
%   Identifie un mémoire réalisé en extension à l'Université du Québec
%   à Chicoutimi (UQAC). Valide uniquement avec un grade de maîtrise.
%   Cette déclaration requiert ensuite d'utiliser les commandes
%   \cmd{\faculteUL} et \cmd{\faculteUQAC}.
% \end{DescribeMacro}
%
% Ensuite, la composition d'une page de titre peut nécessiter
% d'utiliser les commandes ci-dessous en plus de celles de la
% \autoref{sec:utilisation:commandes}.
%
% \begin{DescribeMacro}{\annee}
%   Année du dépôt final de la thèse ou du mémoire.
% \end{DescribeMacro}
%
% \begin{DescribeMacro}{\univcotutelle}
%   Nom, ville et pays de l'université de cotutelle, déclarés sous la
%   forme
%   \begin{quote}
%     |\univcotutelle{Nom de l'université \\ Ville, Pays}|
%   \end{quote}
%   Cette commande prend effet seulement lorsque la classe est chargée
%   avec l'option |cotutelle|.
% \end{DescribeMacro}
%
% \begin{DescribeMacro}{\gradecotutelle}
%   Grade conféré par l'université de cotutelle, déclaré
%   sous la forme
%   \begin{quote}
%     |\gradecotutelle{Nom du grade (sigle)}|
%   \end{quote}
%   Cette commande prend effet seulement lorsque la classe est chargée
%   avec l'option |cotutelle|.
% \end{DescribeMacro}
%
% \begin{DescribeMacro}{\univbidiplomation}
%   Nom, ville et pays de l'université de bidiplomation, déclarés sous
%   la forme
%   \begin{quote}
%     |\univbidiplomation{Nom de l'université \\ Ville, Pays}|
%   \end{quote}
%   Cette commande prend effet seulement lorsque la classe est chargée
%   avec l'option |bidiplomation|.
% \end{DescribeMacro}
%
% \begin{DescribeMacro}{\gradebidiplomation}
%   Grade conféré par l'université de bidiplomation, déclaré sous la
%   forme
%   \begin{quote}
%     |\gradebidiplomation{Nom du grade (sigle)}|
%   \end{quote}
%   Cette commande prend effet seulement lorsque la classe est chargée
%   avec l'option |bidiplomation|.
% \end{DescribeMacro}
%
% \begin{DescribeMacro}{\faculteUL}
%   La commande a deux usages:
%   \begin{enumerate}
%   \item noms des facultés pour les thèses et mémoires
%     multifacultaires, séparés par des commandes |\\|;
%   \item nom de la faculté de l'Université Laval où sont réalisés les
%     thèses et mémoires en extension à l'Université de Sherbrooke, à
%     l'UQO ou à l'UQAC.
%   \end{enumerate}
%   Cette commande prend effet seulement lorsque la classe est chargée
%   avec l'une ou l'autre des options |multifacultaire|,
%   |extensionUdeS|, |extensionUQO| ou |extensionUQAC|.
% \end{DescribeMacro}
%
% \begin{DescribeMacro}{\faculteUdeS}
%   Nom de la faculté de l'Université de Sherbrooke hébergeant la
%   thèse en extension. Cette commande prend effet seulement lorsque
%   la classe est chargée avec l'option |extensionUdeS|.
% \end{DescribeMacro}
%
% \begin{DescribeMacro}{\faculteUQO}
%   Nom de la faculté de l'Université du Québec en Outaouais
%   hébergeant la thèse en extension. Cette commande prend effet
%   seulement lorsque la classe est chargée avec l'option
%   |extensionUQO|.
% \end{DescribeMacro}
%
% \begin{DescribeMacro}{\faculteUQAC}
%   Nom de la faculté de l'Université du Québec à Chicoutimi
%   hébergeant le mémoire en extension. Cette commande prend effet
%   seulement lorsque la classe est chargée avec l'option
%   |extensionUQAC|.
% \end{DescribeMacro}
%
% \begin{DescribeMacro}{\pagetitre}
%   Création de la page de titre. Ne pas utiliser la commande
%   \cmd{\pagetitle} de {\LaTeX} pour ce faire. De toutes les
%   commandes ci-dessus, c'est la seule qui doit se trouver dans le
%   corps du document plutôt que dans le préambule.
% \end{DescribeMacro}
%
%
% ^^A Début du code de la classe
%
% \section{Mise en œuvre}
% \label{sec:meo}
%
% Cette annexe passe en revue le code {\TeX} et {\LaTeX} de la classe.
% Elle n'est susceptible d'intéresser que les personnes qui souhaitent
% explorer comment la classe est programmée.
%
% \subsection{Tests et valeurs booléennes}
% \changes{4.4}{2017-06-01}{Tests et valeurs booléennes réalisés sans ifthen.}
%
% Le paquetage \pkg{ifxetex} permet de tester si un document est
% compilé avec \XeTeX.
%    \begin{macrocode}
%<*class>
\RequirePackage{ifxetex}
%    \end{macrocode}
% Nous définissons ici toutes les valeurs booléennes requises par la
% classe.
%    \begin{macrocode}
\newif\ifUL@babel        \UL@babeltrue         % charger babel?
\newif\ifUL@natbib       \UL@natbibtrue        % charger natbib?
\newif\ifUL@hyperref     \UL@hyperreftrue      % charger hyperref?
\newif\ifUL@chapterbib   \UL@chapterbibfalse   % charger chapterbib?
\newif\ifUL@sectionbib   \UL@sectionbibfalse   % option sectionbib de chapterbib?
\newif\ifUL@isthesis                           % programme est une thèse?
\newif\ifUL@iscotutelle  \UL@iscotutellefalse  % thèse en cotutelle?
\newif\ifUL@isexam       \UL@isexamfalse       % examen de doctorat?
\newif\ifUL@isessay      \UL@isessayfalse      % essai de maîtrise?
\newif\ifUL@isresearch   \UL@isresearchfalse   % projet de recherche?
\newif\ifUL@isworkterm   \UL@isworktermfalse   % rapport de stage?
\newif\ifUL@hassubtitle  \UL@hassubtitlefalse  % document a un sous-titre?
\newif\ifUL@hastitlepage \UL@hastitlepagefalse % page titre utilisée?
%    \end{macrocode}
%
% \subsection{Options de la classe}
%
% Les grandes catégories d'options propres à la classe sont: la
% possibilité d'empêcher le chargement du paquetage \pkg{natbib};
% l'activation des bibliographies par chapitre ou par section; la
% possibilité d'empêcher le chargement du paquetage \pkg{babel}; la
% possibilité d'empêcher le chargement du paquetage \pkg{hyperref}; la
% taille de la police de caractères en points; le type de grade; la
% déclaration qu'il s'agit d'un examen de doctorat ou d'un essai de
% maîtrise.
%
% \begin{macro}{nonatbib}
%   L'option |nonatbib| permet d'empêcher la classe de charger le
%   paquetage \pkg{natbib} en cas d'incompatibilité avec d'autres
%   paquetages spécialisés de mise en forme de la bibliographie.
%    \begin{macrocode}
\DeclareOption{nonatbib}{\UL@natbibfalse}
%    \end{macrocode}
% \end{macro}
%
% \begin{macro}{bibchapitre}
%   \begin{macro}{bibsection}
%     L'option |bibchapitre| entraine le chargement du paquetage
%     \pkg{chapterbib} qui permet de composer des bibliographies
%     multiples dans un document, habituellement une par
%     chapitre. Avec l'option |bibsection|, la bibliographie se
%     présente sous la forme d'une section plutôt que d'un chapitre.
%     La seconde option implique la première.
%    \begin{macrocode}
\DeclareOption{bibchapitre}{\UL@chapterbibtrue}
\DeclareOption{bibsection}{\UL@chapterbibtrue\UL@sectionbibtrue}
%    \end{macrocode}
%   \end{macro}
% \end{macro}
%
% \begin{macro}{nobabel}
%   L'option |nobabel| permet d'empêcher la classe de charger le
%   paquetage \pkg{babel}. Cette option peut s'avérer utile pour
%   les utilisateurs de {\XeLaTeX} qui souhaitent plutôt utiliser
%   \pkg{poyglossia} pour le traitement des langues dans leur
%   document.
%    \begin{macrocode}
\DeclareOption{nobabel}{\UL@babelfalse}
%    \end{macrocode}
% \end{macro}
%
% \begin{macro}{nohyperref}
%   L'option |nohyperref| permet d'empêcher la classe de charger le
%   paquetage \pkg{hyperref}. Cette option peut s'avérer utile pour
%   charger d'autres paquetages comme \pkg{cleveref} ou
%   \pkg{glossaries} après \pkg{hyperref}.
%    \begin{macrocode}
\DeclareOption{nohyperref}{\UL@hyperreffalse}
%    \end{macrocode}
% \end{macro}
%
% \begin{macro}{10pt,11pt,12pt}
%   Les valeurs possibles pour la taille de la police de caractères
%   sont |10pt|, |11pt| et |12pt|. Cette option est gérée au niveau de
%   la classe afin de s'assurer que les divers éléments sur les pages
%   de titre sont toujours de la même taille. La taille de la police
%   par défaut permet de déterminer si, par exemple, le titre du
%   document doit être dans la taille \cmd{\Huge}, \cmd{\huge} ou \cmd{\LARGE} de
%   \class{memoir}.
%
%   La taille de la police est passée à \class{memoir} et la macro
%   \cmd{\UL@ptsize} stocke la taille des caractères pour usage futur.
%    \begin{macrocode}
\newcommand*{\UL@ptsize}{}
\DeclareOption{10pt}{%
  \PassOptionsToClass{10pt}{memoir}
  \renewcommand*{\UL@ptsize}{10}}
\DeclareOption{11pt}{%
  \PassOptionsToClass{11pt}{memoir}
  \renewcommand*{\UL@ptsize}{11}}
\DeclareOption{12pt}{%
  \PassOptionsToClass{12pt}{memoir}
  \renewcommand*{\UL@ptsize}{12}}
%    \end{macrocode}
% \end{macro}
%
% \begin{macro}{PhD,MSc,MA,...}
%   Définition du type de grade et si la thèse ou le mémoire est
%   en cotutelle, en bidiplomation, multifacultaire ou en
%   extension. Seuls les doctorats sont offerts en cotutelle et seules
%   les maîtrises offrent l'option d'une bidiplomation. Les options
%   «multifacultaire» et «en extension» ne sont pertinentes que pour créer
%   une page de titre. La page de titre des programmes en extension
%   comporte une mention «offert en extension» ou «offerte en
%   extension» selon qu'il s'agit d'un doctorat ou d'une maîtrise; le
%   bon terme est défini avec l'option correspondante.
%    \begin{macrocode}
\newcommand*{\UL@typenum}{}
\DeclareOption{LLD}{%
  \UL@isthesistrue
  \renewcommand*{\UL@typenum}{0}
  \newcommand*{\UL@degree}{Docteur en droit (LL.~D.)}}
\DeclareOption{DMus}{%
  \UL@isthesistrue
  \renewcommand*{\UL@typenum}{0}
  \newcommand*{\UL@degree}{Docteur en musique (D.~Mus.)}}
\DeclareOption{DPsy}{%
  \UL@isthesistrue
  \renewcommand*{\UL@typenum}{0}
  \newcommand*{\UL@degree}{Docteur en psychologie (D.~Psy.)}}
\DeclareOption{DThP}{%
  \UL@isthesistrue
  \renewcommand*{\UL@typenum}{0}
  \newcommand*{\UL@degree}{Docteur en th\'eologie pratique (D.~Th.~P.)}}
\DeclareOption{PhD}{%
  \UL@isthesistrue
  \renewcommand*{\UL@typenum}{0}
  \newcommand*{\UL@degree}{Philosophi{\ae} doctor (Ph.~D.)}}
\DeclareOption{MATDR}{%
  \UL@isthesisfalse
  \renewcommand*{\UL@typenum}{0}
  \newcommand*{\UL@degree}{Ma\^itre en am\'enagement du territoire %
    et d\'eveloppement r\'egional (M.ATDR)}}
\DeclareOption{MArch}{%
  \UL@isthesisfalse
  \renewcommand*{\UL@typenum}{0}
  \newcommand*{\UL@degree}{Ma\^itre en architecture (M.~Arch.)}}
\DeclareOption{MA}{%
  \UL@isthesisfalse
  \renewcommand*{\UL@typenum}{0}
  \newcommand*{\UL@degree}{Ma\^itre \`es arts (M.A.)}}
\DeclareOption{LLM}{%
  \UL@isthesisfalse
  \renewcommand*{\UL@typenum}{0}
  \newcommand*{\UL@degree}{Ma\^itre en droit (LL.~M.)}}
\DeclareOption{MErg}{%
  \UL@isthesisfalse
  \renewcommand*{\UL@typenum}{0}
  \newcommand*{\UL@degree}{Ma\^itre en ergoth\'erapie (M.~Erg.)}}
\DeclareOption{MMus}{%
  \UL@isthesisfalse
  \renewcommand*{\UL@typenum}{0}
  \newcommand*{\UL@degree}{Ma\^itre en musique (M.~Mus.)}}
\DeclareOption{MPht}{%
  \UL@isthesisfalse
  \renewcommand*{\UL@typenum}{0}
  \newcommand*{\UL@degree}{Ma\^itre en physioth\'erapie (M.~Pht.)}}
\DeclareOption{MSc}{%
  \UL@isthesisfalse
  \renewcommand*{\UL@typenum}{0}
  \newcommand*{\UL@degree}{Ma\^itre \`es sciences (M.~Sc.)}}
\DeclareOption{MScGeogr}{%
  \UL@isthesisfalse
  \renewcommand*{\UL@typenum}{0}
  \newcommand*{\UL@degree}{Ma\^itre en sciences g\'eographiques (M.~Sc.~g\'eogr.)}}
\DeclareOption{MServSoc}{%
  \UL@isthesisfalse
  \renewcommand*{\UL@typenum}{0}
  \newcommand*{\UL@degree}{Ma\^itre en service social (M.~Serv.~soc.)}}
\DeclareOption{MPsEd}{%
  \UL@isthesisfalse
  \renewcommand*{\UL@typenum}{0}
  \newcommand*{\UL@degree}{Ma\^itre en psycho\'education (M.~Ps.~\'ed.)}}
\DeclareOption{multifacultaire}{%
  \ifUL@isthesis
    \renewcommand*{\UL@typenum}{1}
  \else
    \ClassError{ulthese}{%
      Incompatible option multifacultaire}
    {Use this option with a doctorate degree only.}
  \fi}
\DeclareOption{cotutelle}{%
  \ifUL@isthesis
    \renewcommand*{\UL@typenum}{2}
    \UL@iscotutelletrue
  \else
    \ClassError{ulthese}{%
      Incompatible option cotutelle}
    {Use this option with a doctorate degree only.}
  \fi}
\DeclareOption{bidiplomation}{%
  \ifUL@isthesis
    \ClassError{ulthese}{%
      Incompatible option bidiplomation}
    {Use this option with a master degree only.}
  \else
    \renewcommand*{\UL@typenum}{2}
  \fi}
\DeclareOption{extensionUdeS}{%
  \ifUL@isthesis
    \renewcommand*{\UL@typenum}{3}
    \newcommand*{\UL@offered}{offert}
    \newcommand*{\UL@extensionat}{Universit\'e de Sherbrooke}
    \newcommand*{\UL@extensionloc}{Sherbrooke, Canada}
  \else
    \ClassError{ulthese}{%
      Incompatible option extensionUdeS}
    {Use this option with a doctorate degree only.}
  \fi}
\DeclareOption{extensionUQO}{%
  \ifUL@isthesis
    \renewcommand*{\UL@typenum}{3}
    \newcommand*{\UL@offered}{offert}
    \newcommand*{\UL@extensionat}{Universit\'e du Qu\'ebec en Outaouais}
    \newcommand*{\UL@extensionloc}{Gatineau, Canada}
  \else
    \ClassError{ulthese}{%
      Incompatible option extensionUQO}
    {Use this option with a doctorate degree only.}
  \fi}
\DeclareOption{extensionUQAC}{%
  \ifUL@isthesis
    \ClassError{ulthese}{%
      Incompatible option extensionUQAC}
    {Use this option with a master degree only.}
  \else
    \renewcommand*{\UL@typenum}{3}
    \newcommand*{\UL@offered}{offerte}
    \newcommand*{\UL@extensionat}{Universit\'e du Qu\'ebec \`a Chicoutimi}
    \newcommand*{\UL@extensionloc}{Chicoutimi, Canada}
  \fi}
%    \end{macrocode}
% \end{macro}
%
% \begin{macro}{examen,essai}
%   L'option |examen| change l'appellation «Thèse» sur la couverture
%   pour «Examen de doctorat». L'option |essai|, quant à elle, change
%   l'appellation «Maîtrise» pour «Essai». L'option |examen|
%   n'est compatible qu'avec l'une des options de thèse, alors que
%   |essai| n'est compatible qu'avec une option de maîtrise. Un
%   message d'erreur est émis en cas d'incompatibilité.
%    \begin{macrocode}
\DeclareOption{examen}{%
  \ifUL@isthesis
    \UL@isexamtrue
  \else
    \ClassError{ulthese}{%
      Incompatible option examen}
    {Use this option with a thesis type only.}
  \fi}
\DeclareOption{essai}{%
  \ifUL@isthesis
    \ClassError{ulthese}{%
      Incompatible option essai}
    {Use this option with a Master's type only.}
  \else
    \UL@isessaytrue
  \fi}
%    \end{macrocode}
% \end{macro}
%
% \begin{macro}{projet,stage}
%   L'option |projet| inscrit l'appellation «Projet de recherche» sur
%   la page de titre et aucun type de diplôme, alors que l'option
%   |stage| inscrit «Rapport de stage». La page de titre est
%   autrement traitée comme celle d'une thèse ou d'un mémoire
%   standard. Ces options ont préséance sur toute déclaration de type de
%   doctorat ou de maîtrise. Il faut donc les utiliser avec prudence.
%    \begin{macrocode}
\DeclareOption{projet}{%
  \UL@isresearchtrue
  \renewcommand*{\UL@typenum}{0}
  \def\UL@degree{\mbox{}}}
\DeclareOption{stage}{%
  \UL@isworktermtrue
  \renewcommand*{\UL@typenum}{0}
  \def\UL@degree{\mbox{}}}
%    \end{macrocode}
% \end{macro}
%
% \subsection{Chargement de la classe \class{memoir}}
%
% Toutes les options de la classe sont passées à \class{memoir}. Le
% format de papier et la taille de police par défaut sont, dans
% l'ordre, |letterpaper| et |11pt|. On vérifie qu'un type de grade a
% bien été déclaré (à moins que l'option |projet| ou l'option |stage|
% ne soit présente). L'option de \class{memoir} |oneside| est
% explicitement déclarée afin d'éviter toute tentative de passer outre
% à cette exigence de la FESP.
%    \begin{macrocode}
\DeclareOption*{\PassOptionsToClass{\CurrentOption}{memoir}}
\ExecuteOptions{11pt,letterpaper}
\ProcessOptions\relax
\ifx\UL@typenum\empty
  \ifUL@isresearch\else
    \ifUL@isworkterm\else
      \ClassError{ulthese}{%
        No thesis type specified}
        {Declare the thesis type as a class option.}
    \fi
  \fi
\fi
\LoadClass[oneside]{memoir}
%    \end{macrocode}
%
% \subsection{Paquetages requis}
% \changes{5.0}{2018-06-10}{Le paquetage unicode-math n'est plus
%   chargé par défaut par la classe.}
% \changes{5.0}{2018-08-11}{Le paquetage hyperref est désormais chargé
%   dans la classe plutôt que dans les gabarits. Requiert le paquetage
%   etoolbox.}
%
% La classe s'efforce de charger un minimum de paquetages afin d'éviter
% les conflits potentiels.
%
% {\XeLaTeX} requiert le paquetage \pkg{fontspec} pour le traitement
% des polices. Sous {\LaTeX}, il est aujourd'hui préférable d'utiliser
% les polices T1.
%    \begin{macrocode}
\ifxetex
  \RequirePackage{fontspec}
  \defaultfontfeatures{Ligatures=TeX}
\else
  \RequirePackage[T1]{fontenc}
\fi
%    \end{macrocode}
%
% Les paquetages \pkg{natbib} et \pkg{chapterbib} doivent être chargés
% avant \pkg{babel} pour bien fonctionner, le cas échéant. Tel que
% précisé dans la documentation de \pkg{natbib}, l'option |sectionbib|
% est passée à ce paquetage lorsqu'il est chargé, ou à
% \pkg{chapterbib} autrement.
%    \begin{macrocode}
\ifUL@natbib
  \ifUL@sectionbib
    \PassOptionsToPackage{sectionbib}{natbib}
  \fi
  \RequirePackage[round,semicolon,authoryear]{natbib}
\fi
\ifUL@chapterbib
  \ifUL@sectionbib
     \ifUL@natbib\else
       \PassOptionsToPackage{sectionbib}{chapterbib}
     \fi
  \fi
  \RequirePackage{chapterbib}
\fi
%    \end{macrocode}
%
% Le support pour les langues autres que l'anglais est offert par le
% paquetage \pkg{babel} --- à moins que l'option |nobabel| n'ait
% été spécifiée au chargement de la classe. Les langues sont passées
% en option de la classe, et non du paquetage. Le paquetage
% \pkg{numprint} est requis par \pkg{babel} pour la définition
% de la commande de mise en forme des nombres \cmd{\nombre}.
%    \begin{macrocode}
\ifUL@babel
  \RequirePackage{babel}
  \RequirePackage[autolanguage]{numprint}
\fi
%    \end{macrocode}
%
% Le paquetage obligatoire \pkg{hyperref} est chargé en tout dernier
% avec la commande \cmd{\AtEndPreamble} du paquetage \pkg{etoolbox}.
% Si l'option |nohyperref| a été spécifiée, la classe vérifie que le
% paquetage a été chargé. La configuration de base de la classe fait
% en sorte que tous les liens sont signalés par la couleur
% |ULlinkcolor| définie à la \autoref{sec:couleurs}.
%    \begin{macrocode}
\ifUL@hyperref
  \RequirePackage{etoolbox}
  \AtEndPreamble{\RequirePackage{hyperref}}
\else
  \AtBeginDocument{%
    \@ifpackageloaded{hyperref}{\relax}{%
      \ClassError{ulthese}{%
        Package hyperref is required}
      {Delete option nohyperref or load the package in the preamble.}}}
\fi
\PassOptionsToPackage{colorlinks,allcolors=ULlinkcolor}{hyperref}
%    \end{macrocode}
%
% L'insertion du logo de l'Université sur la page de titre requiert
% \pkg{graphicx}. Les coloration des hyperliens requiert \pkg{xcolor}.
%    \begin{macrocode}
\RequirePackage{graphicx}
\RequirePackage{xcolor}
%    \end{macrocode}
%
% La commande \cmd{\textcopyright} utilisée sur la page de titre requiert le
% paquetage \pkg{textcomp} pour obtenir un beau signe de copyright.
%    \begin{macrocode}
\RequirePackage{textcomp}
%    \end{macrocode}
%
% \subsection{Paquetage incompatible}
%
% Le chargement du paquetage \pkg{geometry} avec la classe
% \class{memoir} modifie les marges du document. Pour cette raison,
% \pkg{geometry} est déclaré incompatible avec la classe.
%    \begin{macrocode}
\AtBeginDocument{%
  \@ifpackageloaded{geometry}{%
    \ClassError{ulthese}{%
      Package geometry is incompatible with this class}
    {Use the memoir class facilities to change the page layout.}}{\relax}}
%    \end{macrocode}
%
% \subsection{Couleur des hyperliens}
% \label{sec:couleurs}
%
% La classe définit une couleur standard pour les hyperliens, une
% teinte de bleu assez foncée pour être à la fois visible en couleur
% et peu contrastante si le document est imprimé en noir et blanc.
%    \begin{macrocode}
\definecolor{ULlinkcolor}{rgb}{0,0,0.3}
%    \end{macrocode}
%
% \subsection{Marges}
%
% Les marges exigées par les normes de présentation de la FESP sont de
% 30~mm pour les marges gauche et droite et 25~mm pour les marges
% supérieure et inférieure. Le pied de page est placé de sorte que le
% folio de page se retrouve à 10~mm du bas de la page.
%    \begin{macrocode}
\setlrmarginsandblock{30mm}{30mm}{*}
\setulmarginsandblock{25mm}{25mm}{*}
\checkandfixthelayout[nearest]
\setlength{\footskip}{\lowermargin}
\addtolength{\footskip}{-10mm}
%    \end{macrocode}
%
% Comme les thèses et mémoires comportent normalement plusieurs pages
% liminaires, il arrive que des folios (en chiffres romains) dépassent
% dans la marge de droite dans la table des matières. Pour régler ce
% problème, nous augmentont la largeur de la boîte prévue pour les
% imprimer.
%    \begin{macrocode}
\renewcommand{\@pnumwidth}{3em}
\renewcommand{\@tocrmarg}{4em}
%    \end{macrocode}
%
% \subsection{Interligne}
%
% L'espacement entre les lignes est d'un interligne et demi.
% L'espacement «double» entre les paragraphes est fixé à
% |0.5\baselineskip| afin d'en arriver à une disposition agréable à
% l'œil. Le retrait de première ligne est supprimé puisque plus
% nécessaire suite à l'ajout de l'espacement entre les paragraphes.
%    \begin{macrocode}
\OnehalfSpacing
\setlength{\parskip}{0.5\baselineskip}
\setlength{\parindent}{0em}
%    \end{macrocode}
%
% La table des matières, la liste des tableaux et la liste des figures
% sont composées à interligne simple.
%    \begin{macrocode}
\renewcommand{\tocheadstart}{\SingleSpacing\chapterheadstart}
\renewcommand{\lotheadstart}{\SingleSpacing\chapterheadstart}
\renewcommand{\lofheadstart}{\SingleSpacing\chapterheadstart}
%    \end{macrocode}
%
% \subsection{Entêtes et pieds de page}
%
% Les règles pour les entêtes et pieds de page sont uniformes pour
% tout le document: aucun entête et folio au centre du pied de page.
% Ceci correspond au style standard |plain|.
%    \begin{macrocode}
\pagestyle{plain}
%    \end{macrocode}
%
% \subsection{Pages de titre}
%
% Le code pour traiter et composer la page de titre (le cas échéant)
% et la page frontispice constitue l'essentiel de la classe.
%
% \subsubsection{Famille et style de la police de caractères}
%
% Les pages de titre sont composées avec la police \textsf{Helvetica}
% (famille \texttt{phv} dans la classification NFSS) dans les
% tailles\footnote{%
%   La police Helvetica produite par {\LaTeX} est plus grande que
%   celle utilisée par Microsoft Word. Pour cette raison, les tailles
%   utilisées dans la classe sont toutes quelques points inférieures à
%   celles des gabarits Word.}%
% et les graisses présentées au \autoref{tab:polices}. La
% déclaration |\fontencoding{T1}| est nécessaire avec {\XeLaTeX} pour
% explicitement charger la même police que sous {\LaTeX}.
%
% \begin{table}
%   \centering
%   \begin{tabular}{ll}
%     \toprule
%     Élément          & Police \\
%     \midrule
%     Titre            & 17~points gras \\
%     Sous-titre       & 14~points gras \\
%     Auteur           & 12~points gras \\
%     Nom du programme & 12~points gras \\
%     Autres éléments  & 12~points normal \\
%     \bottomrule
%   \end{tabular}
%   \caption{Tailles et graisses de la police Helvetica des éléments
%     de la page de titre}
%   \label{tab:polices}
% \end{table}
%
%    \begin{macrocode}
\newcommand*{\UL@phvfamily}{\fontencoding{T1}\fontfamily{phv}\selectfont}
%    \end{macrocode}
% Les commandes sélectionnant ces polices sont adaptées selon la
% taille de police choisie pour le document afin d'être toujours
% identiques. Nous utilisons les déclarations de taille de police de
% la classe \class{memoir}, présentées au tableau~3.9 de sa
% documentation.
%    \begin{macrocode}
\ifnum\UL@ptsize=10\relax
  \newcommand*{\UL@fonttitle}{\normalfont\huge\bfseries\UL@phvfamily}
  \newcommand*{\UL@fontsubtitle}{\normalfont\LARGE\bfseries\UL@phvfamily}
  \newcommand*{\UL@fontauthor}{\normalfont\Large\bfseries\UL@phvfamily}
  \newcommand*{\UL@fontprogram}{\UL@fontauthor}
  \newcommand*{\UL@fontbase}{\normalfont\Large\UL@phvfamily}
\fi
\ifnum\UL@ptsize=11\relax
  \newcommand*{\UL@fonttitle}{\normalfont\LARGE\bfseries\UL@phvfamily}
  \newcommand*{\UL@fontsubtitle}{\normalfont\Large\bfseries\UL@phvfamily}
  \newcommand*{\UL@fontauthor}{\normalfont\large\bfseries\UL@phvfamily}
  \newcommand*{\UL@fontprogram}{\UL@fontauthor}
  \newcommand*{\UL@fontbase}{\normalfont\large\UL@phvfamily}
\fi
\ifnum\UL@ptsize=12\relax
  \newcommand*{\UL@fonttitle}{\normalfont\Large\bfseries\UL@phvfamily}
  \newcommand*{\UL@fontsubtitle}{\normalfont\large\bfseries\UL@phvfamily}
  \newcommand*{\UL@fontauthor}{\normalfont\normalsize\bfseries\UL@phvfamily}
  \newcommand*{\UL@fontprogram}{\UL@fontauthor}
  \newcommand*{\UL@fontbase}{\normalfont\normalsize\UL@phvfamily}
\fi
%    \end{macrocode}
%
% \subsubsection{Interfaces interne et externe}
%
% Définition des commandes permettant de construire les pages de titre.
% L'interface utilisateur est basée sur un ensemble de commandes
% internes. On commence par celles-ci.
%    \begin{macrocode}
\newcommand{\UL@maintitle}{}
\newcommand{\UL@subtitle}{}
\newcommand*{\UL@author}{}
\newcommand*{\UL@year}{}
\newcommand*{\UL@program}{}
\newcommand*{\UL@director}{}
\newcommand*{\UL@codirector}{}
\newcommand*{\UL@nameother}{}
\newcommand*{\UL@degreeother}{}
\newcommand*{\UL@facUL}{}
\newcommand*{\UL@facother}{}
%    \end{macrocode}
% Puis les commandes visibles pour les utilisateurs, qui redéfinissent
% les commandes internes. Voir la \autoref{sec:utilisation:commandes}
% pour leur signification.
%    \begin{macrocode}
\newcommand{\titre}[1]{\renewcommand{\UL@maintitle}{#1}}
\newcommand{\soustitre}[1]{%
  \UL@hassubtitletrue
  \renewcommand{\UL@subtitle}{#1}}
\newcommand*{\auteur}[1]{\renewcommand*{\UL@author}{#1}}
\newcommand*{\annee}[1]{\renewcommand*{\UL@year}{#1}}
\newcommand*{\programme}[1]{\renewcommand*{\UL@program}{#1}}
\newcommand*{\direction}[1]{\renewcommand*{\UL@director}{#1}}
\newcommand*{\codirection}[1]{\renewcommand*{\UL@codirector}{#1}}
\newcommand*{\univcotutelle}[1]{\renewcommand*{\UL@nameother}{#1}}
\newcommand*{\gradecotutelle}[1]{\renewcommand*{\UL@degreeother}{#1}}
\newcommand*{\univbidiplomation}[1]{\renewcommand*{\UL@nameother}{#1}}
\newcommand*{\gradebidiplomation}[1]{\renewcommand*{\UL@degreeother}{#1}}
\newcommand{\faculteUL}[1]{\renewcommand*{\UL@facUL}{#1}}
\newcommand*{\faculteUdeS}[1]{\renewcommand*{\UL@facother}{#1}}
\newcommand*{\faculteUQO}[1]{\renewcommand*{\UL@facother}{#1}}
\newcommand*{\faculteUQAC}[1]{\renewcommand*{\UL@facother}{#1}}
%    \end{macrocode}
%
% \subsubsection{Titre et sous-titre}
%
% Le titre et le sous-titre peuvent s'étendre sur plus d'une ligne.
% Sans traitement spécial, un long titre ou sous-titre aurait pour
% impact de décaler vers le bas tous les autres éléments de la page
% de titre. Pour contrer ce phénomène, nous devrons mesurer la hauteur du
% titre et du sous-titre pour ensuite ajuster en conséquence la
% distance entre ce bloc et les éléments qui suivent.
%
% \begin{macro}{\UL@measuretitle}
%   On place le titre et le sous-titre centrés dans des boîtes
%   \cmd{\UL@titlebox} et \cmd{\UL@subtitlebox}. La commande
%   \cmd{\UL@measuretitle} permettra de mesurer leur hauteur lorsque le
%   titre sera créé avec \cmd{\pagestitre}, plus loin. Un espacement
%   vertical d'un demi interligne est ajouté entre le titre et le
%   sous-titre, le cas échéant.
%    \begin{macrocode}
\newsavebox{\UL@titlebox}
\newsavebox{\UL@subtitlebox}
\newlength{\UL@titleboxtotht}
\newlength{\UL@subtitleboxtotht}
\newcommand{\UL@measuretitle}{%
  \setbox\UL@titlebox=\vbox{%
    \centering\UL@fonttitle\UL@maintitle}
  \setlength{\UL@titleboxtotht}{%
    \dimexpr\ht\UL@titlebox+\dp\UL@titlebox}
  \ifUL@hassubtitle
    \setbox\UL@subtitlebox=\vbox{%
      \centering\vspace*{0.5\baselineskip}\UL@fontsubtitle\UL@subtitle}
    \setlength{\UL@subtitleboxtotht}{%
      \dimexpr\ht\UL@subtitlebox+\dp\UL@subtitlebox}
  \fi}
%    \end{macrocode}
% \end{macro}
%
% \subsubsection{Type de document}
%
% \begin{macro}{\UL@typeofdoc}
%   La commande |\UL@typeofdoc| contient le type de document qui est
%   produit: thèse, thèse en cotutelle, maîtrise, examen de doctorat,
%   essai, projet de recherche ou rapport de stage.
%    \begin{macrocode}
\ifUL@isthesis
  \ifUL@iscotutelle
    \newcommand*{\UL@typeofdoc}{Th\`ese en cotutelle}
  \else
    \newcommand*{\UL@typeofdoc}{Th\`ese}
  \fi
\else
  \newcommand*{\UL@typeofdoc}{M\'emoire}
\fi
\ifUL@isexam
  \renewcommand*{\UL@typeofdoc}{Examen de doctorat}
\fi
\ifUL@isessay
  \renewcommand*{\UL@typeofdoc}{Essai}
\fi
\ifUL@isresearch
  \renewcommand*{\UL@typeofdoc}{Projet de recherche}
\fi
\ifUL@isworkterm
  \renewcommand*{\UL@typeofdoc}{Rapport de stage}
\fi
%    \end{macrocode}
% \end{macro}
%
% \begin{macro}{\UL@docid}
%   La commande |\UL@docid| prépare ensuite la mention mise en forme
%   du type de document pour les pages de titre. La thèse ou le
%   mémoire en cotutelle ou en bidiplomation requiert un traitement
%   différent puisque le programme d'étude apparaît immédiatement sous
%   la mention.
%    \begin{macrocode}
\newcommand{\UL@docid}{%
  {\UL@fontprogram\UL@typeofdoc\par
  \ifnum\UL@typenum=2 \UL@program\par \fi}}
%    \end{macrocode}
% \end{macro}
%
% \subsubsection{Détails sur les facultés et universités d'attache}
%
% \begin{macro}{\Ul@details}
%   La commande |\Ul@details| est la plus complexe puisque la
%   disposition des informations additionnelles sur le document varie
%   beaucoup selon le type de thèse ou de mémoire. Il existe quatre
%   grandes catégories de disposition des éléments sur la page de titre:
%   standard; multifacultaire; en cotutelle ou en bidiplomation
%   (disposition identique); en extension.
%
%   Tel qu'expliqué à l'\autoref{sec:pagetitre}, certains
%   types de grade requièrent expressément que certaines informations
%   soient fournies pour la composition d'une page de titre. Si un
%   élément d'information manque, un avertissement est émis.
%    \begin{macrocode}
\newcommand{\UL@details}{%
  \ifcase\UL@typenum\relax% 0 standard
    \vspace{96pt}
    {\UL@fontprogram\UL@program}\par
    \ifnum\ifUL@isresearch 1\else\ifUL@isworkterm 1\else 0\fi\fi
      =1 %
      \vspace{\baselineskip}
    \else
      \UL@degree
    \fi\par
    \vspace{112pt}
    Qu\'ebec, Canada\par
  \or%                      1 multifacultaire
    \vspace{96pt}
    {\UL@fontprogram\UL@program}\par
    \UL@degree\par
    \vspace{36pt}
    \ifx\UL@facUL\empty
      \ClassWarningNoLine{ulthese}{UL faculty names missing.}
    \else
      \UL@facUL\par
    \fi
    \vspace{48pt}
    Qu\'ebec, Canada\par
  \or%                      2 cotutelle et bidiplomation
    \vspace{72pt}
    Universit\'e Laval\par Qu\'ebec, Canada\par
    \UL@degree\par
    \vspace{\baselineskip} et\par \vspace{\baselineskip}
    \ifx\UL@nameother\empty
      \ClassWarningNoLine{ulthese}{Other university name and location missing}
    \else
      \UL@nameother\par
    \fi
    \ifx\UL@degreeother\empty
      \ClassWarningNoLine{ulthese}{Other university degree missing}
    \else
      \UL@degreeother\par
    \fi
  \or%                      3 extension
    \vspace{48pt}
    {\UL@fontprogram\UL@program\ de l'Universit\'e Laval\par
      \UL@offered\ en extension \`a l'\UL@extensionat}\par
    \vspace{36pt}
    \UL@degree\par
    \vspace{36pt}
    \ifx\UL@facother\empty
      \ClassWarningNoLine{ulthese}{Other university faculty name missing}
    \else
      \UL@facother\par
    \fi
    \UL@extensionat\par
    \UL@extensionloc\par
    \vspace{\baselineskip}
    \ifx\UL@facUL\empty
      \ClassWarningNoLine{ulthese}{UL faculty name missing}
    \else
      \UL@facUL\par
    \fi
    Universit\'e Laval\par Qu\'ebec, Canada\par
  \fi}
%    \end{macrocode}
% \end{macro}
%
% \subsubsection{Conception de la page frontispice}
%
% \begin{macro}{\frontispice}
%   Les thèses et mémoires doivent obligatoirement débuter par une
%   page frontispice. Celle-ci comporte, de haut en bas:
%   \begin{enumerate}
%   \item un espace vide (créé avec une boîte de réglure invisible, ou
%     \emph{strut}) de la même hauteur que le logo de l'Université
%     Laval;
%   \item le titre et le sous-titre, le cas échéant;
%   \item le type de document (thèse, thèse en cotutelle,
%     mémoire, etc.);
%   \item le nom complet de l'auteur;
%   \item les noms des directeur et codirecteurs de recherche.
%   \end{enumerate}
%
%   Nous devons rétablir pour les pages de titre l'interligne simple et
%   l'espacement nul entre les paragraphes (\cmd{\parskip}). Ensuite, nous
%   devons ajuster la distance entre le bloc de titre et le type de
%   document (\cmd{\UL@docidspacing}) et celle entre ce dernier et le nom
%   de l'auteur (\cmd{\UL@authorspacing}). Cela fait en sorte que les
%   éléments des pages de titre se retrouvent (presque) toujours au même
%   endroit sur la page. Une distance minimale d'un interligne est
%   conservée entre le bloc de titre et le type de document
%   (précaution nécessaire pour l'éventuel cas d'un bloc de titre
%   s'étendant sur plusieurs lignes).
%
%   Le nom de l'auteur et ceux des directeurs et codirecteurs de
%   recherche sont insérées directement dans le code de la commande
%   \cmd{\frontispice}.
%    \begin{macrocode}
\newlength{\UL@docidspacing}
\setlength{\UL@docidspacing}{82pt}
\newlength{\UL@authorspacing}
\setlength{\UL@authorspacing}{72pt}
\newcommand{\UL@settitledims}{%
  \UL@measuretitle
  \addtolength{\UL@docidspacing}{-\UL@titleboxtotht}
  \addtolength{\UL@docidspacing}{-\UL@subtitleboxtotht}
  \ifdim\UL@docidspacing<\baselineskip\relax
  \setlength{\UL@docidspacing}{\baselineskip}
  \addtolength{\UL@authorspacing}{-\baselineskip}
  \fi}
\newcommand{\frontispice}{{%
    \ifUL@isexam\else\ifUL@isworkterm\else
      \clearpage
      \pagestyle{empty}
      \SingleSpacing\setlength{\parskip}{0pt}
      \centering
      \UL@fontbase
      \ifUL@hastitlepage\else
        \UL@settitledims
      \fi
      \ifnum\UL@typenum>1\relax
        \vspace*{0pt}\par
      \else
        \rule{0mm}{15mm}\par    % strut
      \fi
      \vspace{82pt}
      \box\UL@titlebox
      \box\UL@subtitlebox
      \vspace{\UL@docidspacing}
      \UL@docid
      \vspace{\UL@authorspacing}
      {\UL@fontauthor\UL@author}\par
      \vspace{72pt}
      Sous la direction de:\par
      \vspace{\baselineskip}
      \UL@director\par
      \UL@codirector
    \fi\fi
    \clearpage}}
%    \end{macrocode}
% \end{macro}
%
% \begin{macro}{\pagetitre}
%   La composition d'une page de titre est optionnelle pour les thèses
%   et mémoires de l'Université Laval. La commande \cmd{\pagetitre}
%   permet néanmoins d'en créer une conforme aux normes de
%   présentation matérielle de la FESP. La page de titre comporte, de
%   haut en bas:
%   \begin{enumerate}
%   \item le logo de l'Université Laval (sauf pour les thèses ou
%     mémoires réalisés en cotutelle, en bidiplomation ou en
%     extension);
%   \item le titre et le sous-titre, le cas échéant;
%   \item le type de document (thèse, thèse en cotutelle,
%     mémoire, etc.);
%   \item le nom complet de l'auteur;
%   \item une description du programme, du grade obtenu et des
%     facultés ou universités d'attache, le cas échéant;
%   \item la mention «Québec, Canada» si le logo de l'Université
%     Laval est présent;
%   \item la notice de copyright.
%   \end{enumerate}
%    \begin{macrocode}
\newcommand{\pagetitre}{{%
    \UL@hastitlepagetrue
    \clearpage
    \pagestyle{empty}
    \SingleSpacing\setlength{\parskip}{0pt}
    \centering
    \UL@fontbase
    \UL@settitledims
    \ifnum\UL@typenum>1\relax
      \vspace*{0pt}\par
    \else
      \includegraphics[height=15mm,keepaspectratio=true]{ul_p}\par
    \fi
    \vspace{82pt}
    \copy\UL@titlebox
    \copy\UL@subtitlebox
    \vspace{\UL@docidspacing}
    \UL@docid
    \vspace{\UL@authorspacing}
    {\UL@fontauthor\UL@author}\par
    \UL@details
    \vfill
    {\textcopyright} \UL@author, \UL@year\par
    \clearpage}}
%    \end{macrocode}
% \end{macro}
%
% \subsection{Listes des figures et des tableaux}
%
% \begin{macro}{\listfigurename}
%   Le paquetage \pkg{babel} définit comme titre pour la liste des
%   figures «Table des figures», alors que la liste des tableaux est
%   «Liste des tableaux». Pour une plus grande symétrie, la classe
%   redéfinit le titre correspondant à \cmd{\listoffigures}. La commande
%   \cmd{\addto} est nécessaire pour éviter que \pkg{babel} redéfinisse
%   le titre à |\begin{document}|.
%    \begin{macrocode}
\ifUL@babel
  \addto\captionsfrench{\renewcommand{\listfigurename}{Liste des figures}}
\fi
%    \end{macrocode}
%   Si \pkg{babel} n'est pas chargé, ce sera à l'utilisateur de faire
%   une correction équivalente. Avec \pkg{polyglossia}, la commande à
%   insérer dans l'entête du document est la même que ci-dessus.
% \end{macro}
%
% \subsection{Dédicace et épigraphe}
%
% La dédicace et l'épigraphe sont mises en forme avec la commande
% \cmd{\epigraph} de \class{memoir}.
% \begin{macro}{\dedicace}
%   La dédicace est une épigraphe simplifiée placée seule sur une
%   page, alignée à droite à une dizaine de lignes de la marge
%   supérieure, sans auteur ou source et sans ligne de démarcation.
%    \begin{macrocode}
\newcommand{\dedicace}[1]{{%
    \clearpage
    \pagestyle{empty}
    \setlength{\beforeepigraphskip}{10\baselineskip}
    \setlength{\epigraphrule}{0pt}
    \epigraphtextposition{flushright}
    \mbox{}\epigraph{\itshape #1}{}}}
%    \end{macrocode}
% \end{macro}
% \begin{macro}{\epigraphe}
%   L'épigraphe de début de document est placée seule sur une page à
%   une dizaine de lignes de la marge supérieure. Pour le reste, on
%   s'en remet à la commande \cmd{\epigraph} de \class{memoir}.
%    \begin{macrocode}
\newcommand{\epigraphe}[2]{{%
    \clearpage
    \pagestyle{empty}
    \setlength{\beforeepigraphskip}{10\baselineskip}
    \mbox{}\epigraph{#1}{#2}}}
%    \end{macrocode}
% \end{macro}
%
% \subsection{Citations}
%
% \begin{environment}{quote}
%   La classe redéfinit l'environnement |quote| de \class{memoir}
%   afin que le texte des citations se trouve en retrait de 10~mm à
%   gauche et à droite, conformément aux règles de présentation de la
%   FESP.
%    \begin{macrocode}
\renewenvironment{quote}{%
  \list{}{\rightmargin 10mm \leftmargin 10mm}%
  \item[]}{\endlist}
%    \end{macrocode}
%   \end{environment}
%
% \begin{environment}{quotation}
%   Il en va de même de l'environnement |quotation|. Cependant, cet
%   environnement passe également à l'interligne simple et la classe
%   ajuste l'espacement vertical entre les paragraphes afin que
%   ceux-ci soient bien distincts les uns des autres tout en demeurant
%   raisonnablement compacts. Cet espacement est ici fixé à 6~points.
%    \begin{macrocode}
\renewenvironment{quotation}{%
  \list{}{%
    \SingleSpacing
    \listparindent 0em
    \itemindent    \listparindent
    \leftmargin    10mm
    \rightmargin   \leftmargin
    \parsep        6\p@ \@plus\p@}%
  \item[]}{\endlist}
%    \end{macrocode}
% \end{environment}
%
% \subsection{Numérotation des divisions du document}
%
% Par défaut, \class{memoir} numérote les divisions du document
% seulement jusqu'au niveau des sections. La classe étend la
% numérotation aux sous-sections.
%    \begin{macrocode}
\setsecnumdepth{subsection}
%</class>
%    \end{macrocode}
% ^^A Fin du code de la classe
%
% ^^A Il faut placer ces notes de changement relatives aux gabarits
% ^^A avant \Finale
% \changes{5.0}{2018-06-08}{Gabarits: meilleure indication du texte à
%   remplacer.}
% \changes{5.0}{2018-06-08}{Gabarits: bonification de la structure
%   d'un chapitre de thèse ou de mémoire par articles.}
% \changes{5.0}{2018-04-25}{Gabarits: ajout d'étiquettes et de renvois
%   dans les gabarits des parties du document.}
% \changes{5.0}{2018-04-25}{Gabarits: automatisation de l'ajout
%   d'entrées dans la table des matières à l'aide d'étiquettes et de
%   renvois.}
% \changes{4.4}{2017-05-16}{Gabarits: nouveau gabarit pour une thèse
%   par articles.}
% \changes{4.2}{2016-03-29}{Gabarits: nouvelles options de babel.}
%
% \Finale
%
% \iffalse
% ^^A Gabarits du document maître
%<*gabarit>
%<phd&standard>%% GABARIT POUR THÈSE STANDARD
%<phd&mesure>%% GABARIT POUR THÈSE SUR MESURE
%<phd&articles>%% GABARIT POUR THÈSE PAR ARTICLES
%<phd&multifac>%% GABARIT POUR THÈSE MULTIFACULTAIRE
%<phd&cotutelle>%% GABARIT POUR THÈSE EN COTUTELLE
%<phd&UdeS>%% GABARIT POUR THÈSE EN EXTENSION À L'UNIVERSITÉ DE SHERBROOKE
%<phd&UQO>%% GABARIT POUR THÈSE EN EXTENSION À L'UQO
%<m&standard>%% GABARIT POUR MÉMOIRE STANDARD
%<m&mesure>%% GABARIT POUR MÉMOIRE SUR MESURE
%<m&bidiplomation>%% GABARIT POUR MÉMOIRE EN BIDIPLOMATION
%<m&UQAC>%% GABARIT POUR MÉMOIRE EN EXTENSION À L'UQAC
%%
%% Consulter la documentation de la classe ulthese pour une
%% description détaillée de la classe, de ce gabarit et des options
%% disponibles.
%%
%% [Ne pas hésiter à supprimer les commentaires après les avoir lus.]
%%
%% Déclaration de la classe avec le type de grade
%<phd>%%   [l'un de LLD, DMus, DPsy, DThP, PhD]
%<m>%%   [l'un de MATDR, MArch, MA, LLM, MErg, MMus, MPht, MSc, MScGeogr,
%<m>%%    MServSoc, MPsEd]
%% et les langues les plus courantes. Le français sera la langue par
%<!articles>%% défaut du document.
%<articles>%% défaut du document. L'option 'bibsection' permet de créer des
%<articles>%% bibliographies par chapitre présentées sous forme de section
%<articles>%% numérotée.
%<phd&standard>\documentclass[PhD,english,french]{ulthese}
%<phd&articles>\documentclass[PhD,bibsection,english,french]{ulthese}
%<phd&cotutelle>\documentclass[PhD,cotutelle,english,french]{ulthese}
%<m&standard>\documentclass[MSc,english,french]{ulthese}
%<m&articles>\documentclass[MSc,bibsection,english,french]{ulthese}
%<m&bidiplomation>\documentclass[MA,bidiplomation,english,french]{ulthese}
  %% Encodage utilisé pour les caractères accentués dans les fichiers
  %% source du document. Les gabarits sont encodés en UTF-8. Inutile
  %% avec XeLaTeX, qui gère Unicode nativement.
  \ifxetex\else \usepackage[utf8]{inputenc} \fi

  %% Charger ici les autres paquetages nécessaires pour le document.
  %% Quelques exemples; décommenter au besoin.
  %\usepackage{amsmath}       % recommandé pour les mathématiques
  %\usepackage{icomma}        % gestion de la virgule dans les nombres

  %% Utilisation d'une autre police de caractères pour le document.
  %% - Sous LaTeX
  %\usepackage{mathpazo}      % texte et mathématiques en Palatino
  %\usepackage{mathptmx}      % texte et mathématiques en Times
  %% - Sous XeLaTeX
  %\setmainfont{TeX Gyre Pagella}      % texte en Pagella (Palatino)
  %\setmathfont{TeX Gyre Pagella Math} % mathématiques en Pagella (Palatino)
  %\setmainfont{TeX Gyre Termes}       % texte en Termes (Times)
  %\setmathfont{TeX Gyre Termes Math}  % mathématiques en Termes (Times)

  %% Options de mise en forme du mode français de babel. Consulter la
  %% documentation du paquetage babel pour les options disponibles.
  %% Désactiver (effacer ou mettre en commentaire) si l'option
  %% 'nobabel' est spécifiée au chargement de la classe.
  \frenchbsetup{%
    StandardItemizeEnv=true,       % format standard des listes
    ThinSpaceInFrenchNumbers=true, % espace fine dans les nombres
    og=«, fg=»                     % caractères « et » sont les guillemets
  }

%<!articles>  %% Style de la bibliographie.
%<!articles> \bibliographystyle{}
%<articles>  %% Suppression du numéro de section de la bibliographie. Utilisation
%<articles>  %% de \extrasfrench parce que c'est la dernière langue déclarée dans
%<articles>  %% \documentclass, ci-dessus.
%<articles>  %\addto\extrasfrench{%
%<articles>  %  \renewcommand{\bibsection}{\section*{\bibname}\prebibhook}}

  %% Composition de la page frontispice. Remplacer les éléments entre < >.
  %% Supprimer les caractères < >. Couper un long titre ou un long
  %% sous-titre manuellement avec \\.
  \titre{<Titre principal>}
  % \titre{Ceci est un exemple de long titre \\
  %   avec saut de ligne manuel}
  % \soustitre{Sous-titre le cas échéant}
  % \soustitre{Ceci est un exemple de long sous-titre \\
  %   avec saut de ligne manuel}
  \auteur{<Prénom Nom>}
%<phd&(standard|cotutelle)>  \programme{Doctorat en <discipline> <-- majeure, s'il y a lieu>}
%<m&standard>  \programme{Maîtrise en <discipline> <-- majeure, s'il y a lieu>}
%<m&bidiplomation>  \programme{Maîtrise en <discipline>}
%<!(cotutelle|bidiplomation)>  \direction{<Prénom Nom>, <directeur ou directrice> de recherche}
%<cotutelle>  \direction{<Prénom Nom>, <directeur ou directrice> de recherche \\
%<cotutelle>             <Prénom Nom>, <directeur ou directrice> de cotutelle}
%<bidiplomation>  \direction{<Prénom Nom>, <directeur ou directrice> de recherche \\
%<bidiplomation>             <Prénom Nom>, <directeur ou directrice> de recherche}
  % \codirection{<Prénom Nom>, <codirecteur ou codirectrice> de recherche}
  % \codirection{<Prénom Nom>, <codirecteur ou codirectrice> de recherche \\
  %              <Prénom Nom>, <codirecteur ou codirectrice> de recherche}

  %% Les commandes ci-dessous servent uniquement pour la création
  %% d'une page de titre (interdite lors du dépôt à la FESP).
  % \annee{<20xx>}
%<cotutelle>  % \univcotutelle{<Université de cotutelle> \\ <Ville>, <Pays>}
%<bidiplomation>  % \univcotutelle{<Université de bidiplomation> \\ <Ville>, <Pays>}
%<cotutelle|bidiplomation>  % \gradecotutelle{<Nom du grade> (<sigle du grade>)}

\begin{document}

\frontmatter                    % pages liminaires

\frontispice                    % production de la page frontispice

\chapter*{Résumé}               % ne pas numéroter
\label{chap-resume}             % étiquette pour renvois
\phantomsection\addcontentsline{toc}{chapter}{\nameref{chap-resume}} % inclure dans TdM

Les environnements forestiers jouent un rôle majeur à l’échelle planétaire à travers leurs fonctions écosystémiques et leur apport économique. 
Toutefois, la croissance démographique humaine et l’augmentation des besoins en produits ligneux entraînent une intensification de l'exploitation des forêts. 
Ces changements peuvent amener  une perte d'habitats essentiels à la biodiversité et perturber les dynamiques écologiques entre les espèces qui peuplent ces écosystèmes. 
Cette étude vise à quantifier les effets des traitements sylvicoles tels que les coupes totales et les coupes partielles sur la dynamique de la faune du sol en forêt mixte tempérée. 
Les objectifs étaient premièrement de mesurer l’impact des coupes forestières sur des variables environnementales influençant l’utilisation de l’habitat par la faune du sol 
et deuxièmement d’évaluer les effets directs et indirects de ces traitements sur les amphibiens et les arthropodes du sol à travers le réseau trophique. 
Nous avons émis une hypothèse que les variables environnementales associées à l’habitat de la petite faune (volume de débris ligneux, ouverture de la canopée et profondeur de litière) varient en fonction de l’intensité des coupes forestières. 
Ainsi, les traitements de coupes constitueraient une variable englobant les changements de conditions environnementales. 
Selon notre deuxième hypothèse, les traitements de coupe forestière entraîneraient une modification de l'utilisation de l'habitat 
par les salamandres cendrées (\textit{Plethodon cinereus}), les carabes (\textit{Carabidae}), et les collemboles (\textit{Collembola}), de manière directe ou via des interactions trophiques, 
allant des prédateurs vers les proies. 
Pour tester ces hypothèses, nous avons développé un modèle d’équations structurelles (SEM) intégrant des modèles d’occupation et des modèles linéaires mixtes. 
Cette approche nous a permis de mesurer simultanément, sous forme de réseau, l’effet des coupes forestières sur les variables environnementales, d’une part, et sur la cooccurrence entre les différents taxons, d’autre part. 
Nos résultats montrent que les traitements de coupes forestières influencent significativement les variables environnementales. 
Les coupes totales ont entraîné une ouverture plus grande de la canopée, une diminution de la profondeur de la litière et une réduction du volume de débris ligneux, 
tandis que les coupes partielles permettent une meilleure rétention de ces attributs environnementaux. 
Cependant, les traitements sylvicoles n’ont, dans l’ensemble, pas eu d’effet direct significatif sur la probabilité d’occupation des salamandres et des carabes, ni sur la biomasse des collemboles. 
Nous avons toutefois mesuré une probabilité d’occupation légèrement plus faible pour les salamandres dans les traitements de coupe totale par rapport aux coupes partielles, bien que cet effet soit marginal (IC à 90\%). 
De plus, aucun changement de cooccurrence entre ces groupes d’espèces n’a été observé, ce qui suggère que les effets des coupes forestières ne se propagent pas à travers le réseau trophique des prédateurs vers les proies du sol forestier. 
Notre étude met en évidence la pertinence des modèles d’équations structurelles pour analyser les réseaux trophiques complexes des écosystèmes forestiers. 
Cette recherche a permis d'approfondir la compréhension des relations existantes entre les pratiques sylvicoles et la dynamique de la faune du sol forestier. 
L'apport de ces connaissances permettra de fournir de meilleurs outils pour orienter les plans de gestion et adapter les pratiques sylvicoles. 


\begin{otherlanguage*}{french}
\textbf{Mots-clés:} \textit{Plethodon cinereus}, carabe, collembole, modèle d'équations structurelles, cooccurrence, coupe forestière.
\end{otherlanguage*}
                % résumé français
\chapter*{Abstract}             % ne pas numéroter
\label{chap-abstract}           % étiquette pour renvois
\phantomsection\addcontentsline{toc}{chapter}{\nameref{chap-abstract}} % inclure dans TdM

Forest environments play a major role on a global scale through their ecosystem functions and the economic benefits they provide. 
However, human population growth and increasing demand for wood products are leading to increased logging. 
These changes can lead to habitat loss essential for biodiversity and disrupt ecological dynamics between species inhabiting these ecosystems. 
This study aims to quantify the effects of silvicultural treatments, such as clearcutting and partial cutting, on soil fauna dynamics in temperate mixed forests. 
The objectives were, first, to assess the impact of logging on environmental variables influencing soil fauna habitat use and, second, 
to evaluate the direct and indirect effects of these treatments on soil amphibians and arthropods through trophic networks. 
We hypothesized that environmental variables associated with small fauna habitats (volume of coarse woody debris, canopy openness, and litter depth) vary depending on the intensity of logging. 
Thus, harvesting treatments would act as an overarching variable encompassing changes in environmental conditions. 
According to our second hypothesis, forest harvesting treatments would modify habitat use by Eastern red-backed salamanders (\textit{Plethodon cinereus}), ground beetles (\textit{Carabidae}), and springtails (\textit{Collembola}), 
either directly or through trophic interactions from predators to prey. 
To test these hypotheses, we developed a structural equation model (SEM) integrating occupancy models and linear mixed models. 
This approach allowed us to simultaneously measure, in a network framework, the effects of forest logging on environmental variables and the cooccurrence between different taxa. 
Our results show that harvesting treatments significantly influence environmental variables. 
Clearcutting led to greater canopy openness, reduced litter depth, and decreased coarse woody debris volume,  
whereas partial cutting allowed for better retention of these environmental attributes. 
Overall, silvicultural treatments had no significant direct effects on the occupancy probabilities of salamanders and ground beetles or the biomass of springtails. 
However, we measured a slightly lower occupancy probability for salamanders in clearcuts compared to partial cuts, although this effect was marginal (90\% CI). 
Moreover, no changes in cooccurrence between these groups were observed, suggesting that the effects of logging do not propagate through trophic networks from predators to prey. 
Our study highlights the relevance of structural equation models for analyzing complex trophic networks of forest ecosystems. 
This research deepens our understanding of the relationships between silvicultural practices and the dynamics of forest soil fauna. 
\hl{The environmental variables influencing the habitat of the studied species (coarse woody debris volume, litter depth, and canopy openness) vary depending on the intensity of logging. 
We therefore recommend promoting harvesting methods that minimize habitat disturbance, such as partial cuts. 
These practices help preserve key environmental attributes, benefiting soil fauna.} 
The knowledge gained from this study provides valuable tools for guiding management plans and adapting silvicultural practices. 

\begin{otherlanguage*}{english}
  \textbf{Keywords:} \textit{Plethodon cinereus}, ground beetle, springtail, structural equation modelling, cooccurrence, logging.
  
\end{otherlanguage*}
              % résumé anglais
\cleardoublepage

\tableofcontents                % production de la TdM
\cleardoublepage

\listoftables                   % production de la liste des tableaux
\cleardoublepage

\listoffigures                  % production de la liste des figures
\cleardoublepage

\dedicace{<Dédicace si désiré>}
\cleardoublepage

\epigraphe{<Texte de l'épigraphe>}{<Source ou auteur>}
\cleardoublepage

\chapter*{Remerciements}        % ne pas numéroter
\label{chap-remerciements}      % étiquette pour renvois
\phantomsection\addcontentsline{toc}{chapter}{\nameref{chap-remerciements}} % inclure dans TdM


 Ce travail n'aurait pas pris sa forme actuelle sans le soutien que j'ai reçu tout au long de mon parcours de maîtrise.

 Mes remerciements vont d'abord à mon directeur, Marc J. Mazerolle, pour m'avoir acceuilli au sein de son laboratoire de recherche et pour son soutien constant tout au long de ce projet. 
 Sa confiance, ses précieux conseils en écologie, en statistique et en gestion de projet, ainsi que sa disponibilité ont été des atouts inestimables.
 Je souhaite également exprimer ma gratitude à mon co-directeur, Mathieu Bouchard, dont le soutien et les conseils dans un domaine que je maîtrisais peu, la foresterie, ont été d'une grande aide. 
 Sa vision du projet et ses recommandations m'ont permis de prendre des décisions plus judicieuses et éclairées.

 Je suis également reconnaissant envers Patricia Raymond et Émilie Champagne du Ministère des Ressources Naturelles et de Forêts, Direction de la recherche forestière, pour leur aide et 
 pour m'avoir permis de participer au projet enrichissant qu'est DREAM-Qc.
 Un grand merci également aux techniciennes Karine Thériault et Marie-Claude Martin pour leur soutien logistique et leurs conseils avisés.
 Je suis reconnaissant pour l'aide que j'ai reçu de la part de mon auxiliaire de terrain, Rebecca Dubé Messier. Sa motivation et son travail assidu ont largement contribué à la qualité de ce mémoire.

 Le parcours d'un projet de maîtrise et toujour plus facile quand on est bien entouré. 
 Pour cela je remercie les membres de mon laboratoire : 
 Aurore Fayard, Laura Millard, Mariano Feldman, Anais Baillet, Lucas Voirin, Félicia Beaulieu, Jeanne Dudemaine, Laurianne Plante, Naomie Herpin Saunier. 
 Votre soutien, nos activités extérieures ainsi que nos discussions ont été des moments précieux pour moi.
 Un merci spécial à Joëlle Spooner, qui a été une amie, une colocataire et une collègue de laboratoire exceptionnelle tout au long de ce parcours. 
 Sa présence a été un véritable soutien.

 Je souhaite également remercier mes parents, Isabelle et Patrick, mes beaux-parents, Michel et Anne-Chantal, ainsi que mes frères, Pierre et Guillaume, pour leur soutien depuis 
 mon plus jeune age. Votre support dans ma passion pour la biologie est inestimable.
 Je conclus ces remerciements en exprimant ma reconnaissance envers la personne qui m'accompagne au quotidien, Léonie St-Onge. 
 Merci pour ton écoute, ta patience, tes conseils et pour être simplement toi.

         % remerciements
\chapter*{Avant-propos}         % ne pas numéroter
\label{chap-avantpropos}        % étiquette pour renvois
\phantomsection\addcontentsline{toc}{chapter}{\nameref{chap-avantpropos}} % inclure dans TdM

%\usepackage[french]{babel}

Ce mémoire « par article » en Sciences forestières présenté à l'Université Laval est divisé en plusieurs parties : 
l'introduction, écrite en français, aborde en premier lieu la problématique liée à l'impact potentiel des traitements sylvicoles sur 
l'utilisation de l'habitat par la petite faune du sol dans le cadre d'un projet de migration assistée des arbres.
Le corps de ce document comprend le chapitre principal, intitulé 
"Direct and indirect effects on soil fauna of silvicultural treatments in the context of forest assisted migration", rédigé en anglais sous la forme d'un article scientifique. 
Suite au dépôt de ce mémoire, ce chapitre sera soumis à la revue scientifique \textit{Forest Ecology and Management}. 
La conclusion, rédigée en français, constitue la dernière partie de ce mémoire.
William Devos est l'auteur principal de cet article, suivi de Mathieu Bouchard, co-directeur de recherche, et Marc J. Mazerolle, directeur de recherche et dernier auteur. 
Les trois auteurs ont participé à l'élaboration du projet et à la rédaction. 
William Devos a pris en charge la planification du terrain, ainsi que la collecte et le traitement des données.
Marc J. Mazerolle et William Devos ont élaboré l'intégralité des modèles statistiques.
Marc J. Mazerolle et Mathieu Bouchard ont supervisé l'ensemble du projet de recherche.
           % avant-propos

\mainmatter                     % corps du document

\chapter*{Introduction}         % ne pas numéroter
\label{chap-introduction}       % étiquette pour renvois
\phantomsection\addcontentsline{toc}{chapter}{\nameref{chap-introduction}} % inclure dans TdM

% \usepackage[french]{babel}

\section*{Mise en contexte}
\label{sec:contexte}
\phantomsection\addcontentsline{toc}{section}{\nameref{sec:contexte}}

Les environnements forestiers jouent un rôle majeur à l’échelle planétaire à travers leurs fonctions écosystémiques et l'apport \hl{économique qu’ils représentent} en fournissant de nombreux biens et services \citep{Balvanera2006Quantifyingevidence}. 
Ces milieux permettent de réguler les flux de nutriments et d’énergie en participant notamment à la séquestration du carbone, à la stabilisation du climat, à la rétention de l’eau ainsi qu’à la conservation de la biodiversité \citep{Balvanera2006Quantifyingevidence,Diaz2006BiodiversityLoss,Canadell2008Managingforests,Pawson2013Plantationforests}. 
Toutefois, la croissance démographique et l’augmentation des besoins en produits ligneux amènent une hausse de l’exploitation forestière depuis maintenant plusieurs décennies \citep{Foley2005GlobalConsequences}. 
Par conséquent, le changement d’affectation des sols, causé par la récolte de bois, l’agriculture et l’urbanisation, occasionne une diminution globale du territoire forestier, entraînant une perte de biodiversité \citep{Bengtsson2000Biodiversitydisturbances,Sala2000Globalbiodiversity,Naeem2012functionsbiological,Bichet2016Maintaininganimal}. 
La biodiversité est cependant essentielle pour le bon fonctionnement des écosystèmes terrestres, et plusieurs engagements internationaux ont souligné l'urgence de freiner ce déclin et encouragent une gestion durable des forêts \citep{Newbold2015Globaleffects}. 

De nombreuses recherches ont démontré l'impact de la gestion forestière sur différents groupes fauniques, tels que les oiseaux, les mammifères, les reptiles, les amphibiens et les insectes \citep{Drapeau2000LandscapeScaleDisturbances,Suffice2015Shorttermeffects,bouchardBeetleCommunityResponse2016b,Hughes2019Impactnatural,Mazerolle2021Woodlandsalamander}. 
Les pratiques sylvicoles peuvent provoquer une hausse de mortalité chez les espèces, perturber ou restreindre leurs déplacements, intensifier les interactions biotiques, modifier leurs cycles de vie et altérer leur morphologie ainsi que leur physiologie \citep{Sergio2018Animalresponses}. 
De plus, l'exploitation des peuplements entraîne des pertes d'habitat et la fragmentation des milieux naturels, diminuant ainsi l'accès à la nourriture, aux refuges et aux sites de reproduction, tout en réduisant la taille et la diversité génétique des populations \citep{Bouderbala2023Longtermeffect}.
Ces pratiques causent également des problèmes de connectivité entre les habitats, les groupes d'espèces et les processus écologiques \citep{Lindenmayer2006Generalmanagement}. 
Une baisse de la connectivité peut compromettre la capacité des populations à se rétablir et à se maintenir après une perturbation, ce qui affecte négativement la conservation de la faune \citep{Lamberson1994ReserveDesign}. 
Cela restreint de la même façon le déplacement des individus au sein de leur habitat et le long des corridors écologiques, limitant ainsi le flux génétique entre populations et augmentant le risque d'extinction locale \citep{Saccheri1998Inbreedingextinction}. 

De plus, à l’échelle du paysage, l'exploitation forestière peut amener une homogénéisation du paysage en ramenant les peuplements à des stades précoces de succession, ce qui favorise une sous-représentation des forêts des stades de succession plus avancés \citep{Cyr2009Forestmanagement,Boucher2017Cumulativepatterns}. 
Cette pratique peut réduire la complexité structurelle des forêts, affectant parfois des aspects tels que la composition des espèces d'arbres, la stratification verticale, la structure d'âges, la dynamique de succession et la fréquence des perturbations \citep{Bergeron2000Speciesstand,Commarmot2005Structurevirgin,Varga2005Treesizediversity}. 
Ces modifications sont particulièrement préoccupantes pour la biodiversité, notamment pour les espèces qui vivent en forêts mixtes et conifériennes matures \citep{Tremblay2018Harvestinginteracts,Cadieux2020Projectedeffects}.
Le maintien de l'hétérogénéité structurelle joue un rôle important, car cela amène une variété d'habitats et permet de renforcer la connectivité en améliorant la dispersion de certaines espèces. 
C’est pourquoi l’aménagement forestier peut représenter un risque particulièrement important pour les espèces dépendantes des attributs de forêts complexes ainsi que pour celles ayant une faible capacité de dispersion  \citep{Norden2001Conceptualproblems,Martin2021indicatorspecies}. 
De plus, un peuplement avec une structure plus hétérogène est plus résilient lors d'une perturbation, car il facilite la régénération des espèces d'arbres présentes et le retour de celles qui ont disparu \citep{Kuuluvainen2009Forestmanagement}. 

Les modifications des caractéristiques forestières peuvent participer à la perte d'une partie de la diversité spécifique et fonctionnelle associée au milieu forestier. 
Certaines pratiques de gestion peuvent ainsi réduire la résilience des forêts à l'échelle locale et entraîner une diminution des services écosystémiques \citep{Hooper2012globalsynthesis,Edwards2014Maintainingecosystem}. 


\section*{Traitements sylvicoles}
\label{sec:sylvicole}
\phantomsection\addcontentsline{toc}{section}{\nameref{sec:sylvicole}}

Au cours des dernières décennies, les pratiques sylvicoles qui favorisent la sauvegarde de la biodiversité ont été encouragées afin de réduire les effets néfastes de la gestion intensive \citep{Gustafsson2012Retentionforestry}. 
\hl{L'aménagement écosystémique} a ainsi été suggérée comme une stratégie pour une gestion forestière durable \citep{Perry1998scientificbasis,Kuuluvainen2002Naturalvariabilitya}. 
Cette approche a pour but d'imiter les perturbations naturelles, les structures des peuplements et les dynamiques de succession qui en découlent. 
Elle s'inspire des processus naturels de la forêt, comme les ouvertures créées par les chablis, la mortalité partielle ou les perturbations à plus grande échelle comme les feux de forêt ou les épisodes cycliques de ravageurs \citep{Bergeron1999Forestmanagementa,Bergeron2007Usingknowledge}.
Elle permet également de préserver la biodiversité et d'entretenir la résilience des écosystèmes tout en assurant une variété de services écosystémiques \citep{Szaro1998emergenceecosystem,MacDicken2015Globalprogress}. 
\hl{Plus précisément, l'aménagement écosystémique utilise les traitements sylvicoles à l'échelle des peuplements, comme la coupe totale et la coupe partielle, pour reproduire les impacts écologiques associés aux perturbations naturelles. }
Toutefois, le niveau de perturbation varie selon le type de traitement appliqué, allant de perturbations plus légères à des perturbations plus intensives dans le milieu \citep{Thiffault2006Harvestingintensity}. 

Historiquement, les coupes totales sont parmi les pratiques sylvicoles les plus répandues dans les forêts tempérées et boréales \citep{Fedrowitz2014Canretention,Chaudhary2016Impactforest}. 
Elles font partie d’une gestion forestière permettant d’accroître la rentabilité des opérations forestières à relativement court terme, afin de répondre aux besoins des industries de transformation du bois \citep{Irland2011TimberProductivitya}. 
Sur le plan technique, elles sont relativement simples à réaliser, car l'intégralité du couvert forestier est retirée en une seule récolte, occasionnant une déforestation temporaire de la parcelle. 
Les coupes totales entraînent généralement une homogénéisation de la structure des peuplements \citep{Brashears2004AssessmentCanopy,Martin2020Forestmanagement}. 
\hl{Cette pratique peut ainsi être intégrée dans une approche de gestion écosystémique visant à imiter des perturbations naturelles de grande ampleur, telles que les feux de forêt, et à favoriser le maintien des stades précoces de la succession forestière} \citep{Bergeron1999Forestmanagementa}. 
\hl{De plus, certaines espèces d’arbres se régénèrent plus efficacement en coupe totale, en fonction de leur tolérance à l’ombre et de leur dépendance aux perturbations du sol} \citep{Bergeron1999Forestmanagementa}. 
Toutefois la simplification engendré par ce type de traitement peu parfois bouleverser des processus écologiques et évolutifs en amenant un changement drastique au niveau de la diversité biologique \citep{Holling2001UnderstandingComplexity}. 
Ainsi, l’utilisation de méthodes qui modifient radicalement la structure des milieux par rapport à leur état d’origine peut accroître le risque de déclin de certaines communautés biologiques \citep{Hanski2000Extinctiondebt}. 

Les coupes partielles sont un autre exemple d'\hl{aménagement écosystémique} visant à maintenir la composition, la structure et les fonctions de l'écosystème dans les limites de leur variabilité naturelle \citep{Bergeron1999Forestmanagementa,Raymond2009irregularshelterwood}.
Ce type de coupe implique une récolte sélective d'arbres, souvent orientée par des critères écologiques ou économiques spécifiques. 
Cela permet de conserver une couverture forestière partielle qui protège les semis et favorise une régénération naturelle \citep{Raymond2009irregularshelterwood}.
Les coupes partielles favorisent une structure forestière irrégulière en maintenant plusieurs classes d'âge ou de tailles d'arbres. 
Cette hétérogénéité peut inclure une diversité verticale et horizontale à l'intérieur des peuplements forestiers \citep{Raymond2009irregularshelterwood}.
\hl{Ainsi, les coupes partielles peuvent être appliquées afin de maintenir les stades de succession plus près de la dynamique naturelle des écosystèmes en forêt mixte} \citep{Bergeron1999Forestmanagementa}. 
Contrairement aux coupes totales, les coupes partielles ont une période prolongée de régénération, 
permettant une gestion à long terme et l'intégration d'objectifs de conservation \citep{Raymond2009irregularshelterwood}. 
Les coupes partielles \hl{peuvent permettre} une meilleure conservation des espèces associées au milieu forestier en maintenant des attributs importants comme les débris ligneux, 
plusieurs strates de canopée, différentes tailles d'arbres et un microclimat favorable, 
ce qui est bénéfique pour les espèces sensibles aux variations climatiques ou à la perte d'habitat \citep{Hansen1991Conservingbiodiversity}. 


\section*{Faune du sol}
\label{sec:soilfauna}
\phantomsection\addcontentsline{toc}{section}{\nameref{sec:soilfauna}}

Au sein de la biodiversité forestière, la faune du sol joue un rôle écologique prédominant en contribuant à la circulation de la matière et de l'énergie à travers la chaine alimentaire et en participant au recyclage des nutriments \citep{Seibold2021contributioninsects,Kudrin2023metaanalysiseffects}. 
Toutefois, ces organismes ont souvent une capacité de dispersion limitée en raison de leur petite taille, ce qui les rend vulnérables face aux pertes locales d’habitats \citep{Kudrin2023metaanalysiseffects}. 
Cette communauté est ainsi identifiée comme l'une des plus touchées par les perturbations environnementales \citep{Marshall2000Impactsforest,Coyle2017Soilfauna}. 

Les coupes forestières amènent généralement des changements brusques et importants dans les habitats de la faune du sol \citep{Lindo2003Microbialbiomass,Paillet2010Biodiversitydifferences,Fedrowitz2014Canretention,Chaudhary2016Impactforest}. 
Par exemple, le retrait de la canopée augmente l'exposition de la surface du sol au rayonnement solaire, entraînant une hausse des températures ainsi qu'une modification de l'humidité du sol \citep{Lindo2003Microbialbiomass,Brook2008Synergiesextinction,Zhang2022Intensiveforest}. 
Ces effets sont amplifiés par une augmentation de la vitesse du vent et une intensification des précipitations atteignant le sol \citep{Keenan1993ecologicaleffects,Heithecker2007Edgerelatedgradients}. 
Les modifications de la structure du sol, telles qu'une compaction accrue due au passage des machines forestières, réduisent également la porosité du sol et affectent négativement la faune qui y vit \citep{Battigelli2004Shorttermimpact}. 
De plus, les opérations forestières perturbent la disponibilité des nutriments en modifiant la qualité et la quantité de la litière, en altérant les sécrétions racinaires, en favorisant le lessivage et en modifiant les propriétés chimiques du sol \citep{Covington1981Changesforest,Marshall2000Impactsforest,Lindo2003Microbialbiomass,Battigelli2004Shorttermimpact}. 
Ces bouleversements des conditions microclimatiques et physico-chimiques ont des répercussions directes sur la biodiversité, 
notamment sur les organismes dépendant de microhabitats tels que le bois mort, les cavités des arbres matures ou les plaques racinaires \citep{Spies1999Dynamicforest,Christensen2005Deadwood,Brassard2008EffectsForest}. 
Ces éléments structurels, souvent réduits ou éliminés par les opérations de coupe, sont généralement remplacés par des pistes de débardage et des chemins d'exploitation \citep{Hansen1991Conservingbiodiversity}. 

La présence de microhabitats joue pourtant un rôle essentiel pour une grande partie de la faune du sol, notamment les arthropodes, les amphibiens et les microorganismes \citep{Paillet2010Biodiversitydifferences,Fedrowitz2014Canretention,Chaudhary2016Impactforest}. 
La disparition de ces structures augmente la vulnérabilité de nombreuses espèces, tout en créant des conditions environnementales défavorables à leur survie. 
Par conséquent, les espèces forestières qui dépendent de la fraîcheur et de l'humidité caractéristiques des sols non perturbés peuvent voir leurs populations décliner, voire disparaître localement \citep{Kudrin2023metaanalysiseffects}. 
À l'inverse, les espèces adaptées aux milieux plus ouverts et secs peuvent coloniser les parcelles récemment coupées, entraînant ainsi une modification de la composition spécifique. 

Plusieurs études se sont intéressées aux changements des caractéristiques forestières comme le débris ligneux grossier, la profondeur de la litière, l'ouverture de la canopée, 
ainsi que sur leur impact sur la faune du sol après les opérations de récolte, dans le but d'orienter la gestion forestière \citep{Semlitsch2002CriticalElements,McKenny2006Effectsstructural}. 
Cependant, la faune du sol englobe une grande diversité de taxons, présentant des différences biologiques et écologiques significatives \citep{Kudrin2023metaanalysiseffects}. 
En conséquence, leur réponse à l'exploitation forestière varie selon le type de traitement appliqué et le groupe d'espèces étudiées \citep{Malmstrom2009Dynamicssoil,Paillet2010Biodiversitydifferences}. 

Les amphibiens et les arthropodes jouent un rôle important dans l'écologie forestière, principalement au niveau de la chaine trophique et dans la circulation du carbone au niveau des sols \citep{Burton1975Energyflow,loveiEcologyBehaviorGround1996,Handa2014Consequencesbiodiversity}. 
Malgré le fait qu'ils soient abondants dans les forêts nord-américaines, ces groupes ne sont pas souvent pris en compte dans les plans de gestion forestière \citep{deMaynadier1995relationshipforest}. 
Pourtant, les amphibiens et les arthropodes ont connu un fort déclin au cours des dernières décennies, principalement en raison des pratiques d'aménagement forestier \citep{Houlahan2000Quantitativeevidence,Stuart2004Statustrends,Wagner2021Insectdecline}. 
Leur sensibilité aux changements environnementaux et leur capacité de dispersion limitée en font des modèles d'étude pertinents pour mieux comprendre les effets des pratiques forestières 
sur la biodiversité et l'intégrité écologique des forêts \citep{pongeVerticalDistributionCollembola2000,Ojala2001Dispersalmicroarthropods,birdChangesSoilLitter2004,Maleque2009Arthropodsbioindicators}. 
Parmi les groupes d'espèces fréquemment étudiés chez ces taxons, on retrouve la salamandre cendrée de l'Est (\textit{Plethodon cinereus} (Green, 1818)), les carabes (Carabidae) et les collemboles (Collembola). 

\subsection*{Salamandre cendrée}

La salamandre cendrée, de la famille des Plethodontidae, représente l’une des plus grandes biomasses de vertébrés dans de nombreuses forêts de l’est de l’Amérique du Nord \citep{Burton1975Salamanderpopulations,Petranka1993Effectstimber,semlitschAbundanceBiomassProduction2014a}. 
En tant qu'espèce exclusivement terrestre et dépourvue de poumons, elle dépend de la respiration cutanée et donc des conditions d'humidité pour assurer ses échanges gazeux \citep{Heatwole1961Relationsubstrate}. 
Ce mode de respiration la contraint à occuper des microhabitats spécifiques, en surface lorsque la température et l'humidité sont favorables, ou en profondeur dans le sol durant des périodes moins propices \citep{Grizzell1949HibernationSite,FraserEmpiricalEvaluation1976,Jaeger1980MicrohabitatsTerrestrial}. 
Le retrait des refuges à la surface, tels que les débris ligneux, après la récolte peut réduire la qualité de l'habitat pour les salamandres et diminuer le temps qu'elles passent à la surface du sol forestier \citep{Achat2015Quantifyingconsequences,Peele2017Effectswoody}. 
Étant donné que les salamandres cendrées se nourrissent et se reproduisent à la surface du sol forestier, des conditions de surface défavorables, un compactage des sols et de faibles niveaux de débris ligneux peuvent affecter la dynamique des populations \citep{Peterman2014Spatialvariation}. 
Prédateur généraliste, la salamandre cendrée contribue de manière significative à la régulation des invertébrés détritivores, influençant directement les processus de décomposition, la circulation des nutriments dans les sols et la dynamique du carbone \citep{Burton1975Energyflow,Wyman1998Experimentalassessment,Walton2013Topdownregulation,Hickerson2017Easternredbacked}. 
Par ailleurs, elle constitue une proie de haute valeur nutritive pour divers prédateurs tels que les oiseaux, les mammifères et les reptiles, renforçant ainsi son rôle dans les dynamiques trophiques forestières \citep{Burton1975Energyflow,Pough1987abundancesalamanders}. 
Compte tenu de sa sensibilité aux perturbations environnementales, notamment en raison de sa respiration cutanée, la salamandre cendrée est souvent utilisée comme indicateur de la qualité des sols forestiers \citep{Welsh2001caseusing,Fisher-Reid2024Easternredbacked}. 
Son abondance et sa présence sont des paramètres couramment évalués pour mesurer l'impact des pratiques sylvicoles sur la biodiversité \citep{Harpole1999Effectsseven,Grialou2000effectsforest,Homyack2009Longtermeffects,Hocking2013Effectsexperimental,Mazerolle2021Woodlandsalamander}. 

\subsection*{Carabes}

Avec environ 40 000 espèces répertoriées, les carabes constituent l'une des familles les plus diversifiées parmi les coléoptères et représentent l'une des plus fortes abondances d'arthropodes du sol \citep{Erwin1985taxonpulse,loveiEcologyBehaviorGround1996,Rochefort2006GroundBeetle}. 
Ils sont largement distribués et abondants dans presque tous les écosystèmes terrestres, bien que les espèces varient dans leur sélection d'habitats \citep{loveiEcologyBehaviorGround1996,kotzeFortyYearsCarabid2011a,Larochelle2003naturalhistory}. 
Les changements environnementaux et les modifications d'habitat favorisent certaines espèces au détriment d'autres, ce qui rend leur diversité et leur sensibilité particulièrement intéressantes pour étudier l'impact des perturbations environnementales \citep{Rainio2003Groundbeetles}. 
Ces insectes ont une courte durée de vie, occupent une position élevée dans la chaine alimentaire et réagissent rapidement et de manière complexe aux modifications de leur environnement \citep{loveiEcologyBehaviorGround1996}.
Les carabes jouent un rôle écologique dans la régulation des populations d'invertébrés, se nourrissant principalement d'aphides, de collemboles et de gastéropodes, tout en étant à leur tour des proies pour divers amphibiens, reptiles, oiseaux et mammifères \citep{loveiEcologyBehaviorGround1996}. 
Les carabes sont particulièrement sensibles à la complexité structurelle des peuplements forestiers, et ce, à différentes échelles temporelles et spatiales \citep{Butterfield1995Carabidbeetle,loveiEcologyBehaviorGround1996,Niemela2007effectsforestry}.

\subsection*{Collemboles}

Parmi les taxons de la mésofaune du sol, les collemboles représentent un des groupes les plus abondants et les plus diversifiés de petits arthropodes \citep{rusekBiodiversityCollembolaTheir1998}. 
Différentes communautés de collemboles occupent un ensemble de niches écologiques allant de la litière aux différents horizons du sol \citep{pongeVerticalDistributionCollembola2000}.
La répartition verticale de ces communautés dépend essentiellement des conditions abiotiques du sol telles que la luminosité, le taux d’humidité ou encore la porosité.
Ces invertébrés exercent une influence sur l'écologie forestière à travers leur influence sur les micro-organismes, sur la décomposition et le cycle des nutriments.
Étant principalement fongivores et détritivores, les différentes communautés de collemboles influencent les taux de décomposition du bois, affectent la structure physique et les taux de minéralisation de la litière, 
influencent l'absorption des nutriments et régulent la communauté microbienne, en plus de participer à la formation de microstructure du sol \citep{Petersen1982comparativeanalysis,Neher2012Linkinginvertebrate,Maass2015Functionalrole,Potapov2016Connectingtaxonomy}. 
Les collemboles jouent également un rôle dans chaine trophique en constituant une ressource alimentaire abondante pour plusieurs organismes comme les amphibiens, les coléoptères, les arachnides, les oiseaux et les reptiles \citep{Burton1975Energyflow,Bauer1982Predationcarabid,rusekBiodiversityCollembolaTheir1998}.

\section*{Objectifs et hypothèses}
\label{sec:objectifs}
\phantomsection\addcontentsline{toc}{section}{\nameref{sec:objectifs}}

Grâce à leurs relations trophiques, leur sensibilité aux variations des conditions environnementales et leur dépendance aux caractéristiques forestières, ces trois groupes d'espèces constituent des modèles pertinents pour étudier l'impact des pratiques forestières sur la faune du sol. 
La plupart des travaux qui se sont intéressés aux impacts des traitements sylvicoles sur la faune discutent généralement des effets directs des perturbations sur un ou plusieurs groupes d'espèces, 
sans tenir compte des relations existantes entre les variables environnementales et les différents groupes d'espèces, 
négligeant ainsi les effets indirects des coupes sur la faune du sol \citep{josephIntegratingOccupancyModels2016,Pollierer2021Diversityfunctional,Kudrin2023metaanalysiseffects}. 
\hl{Mon projet a pour but de combler cette lacune en essayant de comprendre comment les effets à court terme de différents types de récolte forestière se propagent à l’intérieur du réseau écologique forestier et influencent la dynamique de la faune du sol. }
Ultimement, ces nouvelles connaissances permettront de mieux informer la gestion durable des forêts.
Les objectifs liés à cette étude sont :

\begin{enumerate}
    \item De quantifier l'effet des traitements de coupes forestières sur les variables environnementales qui influencent l'utilisation de l'habitat par la faune du sol. 
    \item De mesurer l'impact des coupes forestières sur l'utilisation de l'habitat par la faune du sol.
\end{enumerate}

L’hypothèse 1.1 attachée à notre premier objectif soutient que les variables environnementales (volume de débris ligneux, ouverture de la canopée, profondeur de litière) favorables à l’utilisation de l’habitat par les espèces varient en fonction de l’intensité des coupes forestières. 
En ce sens, l’intensité du traitement de coupe est utilisée comme indicateur global représentant les changements dans les conditions environnementales.

L'hypothèse 2.1 liée à notre deuxième objectif stipule que les traitements de coupe forestière entraînent une modification de l'utilisation de l'habitat 
par la faune du sol et que ces effets se propagent à travers le réseau trophique, allant des prédateurs vers les proies. 
Spécifiquement, les coupes affectent l'utilisation de l'habitat par les salamandres et les grands carabes (compétiteurs de salamandres), 
ce qui modifie ensuite la sélection d'habitats des petits carabes (proies de salamandres), puis enfin des collemboles (proies de grands carabes et salamandres). 
Cette hypothèse repose sur le fait que la salamandre cendrée, en tant que prédateur sensible aux perturbations, est reconnue pour son rôle important dans la régulation des communautés d'invertébrés du sol forestier \citep{Wyman1998Experimentalassessment,MichaelWalton2005Salamandersforestfloor,Walton2006Salamandersforestfloor,Walton2013Topdownregulation,Hickerson2017Easternredbacked}. 
Parallèlement, les carabes ont une relation complexe avec les salamandres, se positionnant à la fois comme compétiteurs et proies de celles-ci \citep{Jaeger1980MicrohabitatsTerrestrial,loveiEcologyBehaviorGround1996,Gall2003BehavioralInteractions}. 


\cleardoublepage

\bibliography{References.bib}
\bibliographystyle{ecologyNewFR.bst}
          % introduction
%<!articles>\include{chapitre1}             % chapitre 1
%<articles>\chapter{Direct and indirect effects on soil fauna of silvicultural treatments in the context of forest assisted migration}     % numéroté
\label{chapitre1-articles}    

William Devos$^1$, Mathieu Bouchard$^1$, Marc Mazerolle$^1$

%href{mailto:william.devos.1@ulaval.ca}
$^1$ Centre d'étude de la forêt, Département des sciences du bois \\ 
et de la forêt, Université Laval, Québec, QC G1V 0A6, Canada. \\ 

\clearpage

\section*{Résumé}
\label{sec:resume1}
\phantomsection\addcontentsline{toc}{section}{\nameref{sec:resume1}}

\begin{otherlanguage*}{french}
  <Résumé de l'article en français. Obligatoire.>

  \textbf{Mots-clés} : <ajouter des mots clés>
\end{otherlanguage*}

\clearpage

\section*{Abstract}
\label{sec:abstract1}
\phantomsection\addcontentsline{toc}{section}{\nameref{sec:abstract1}}

\begin{otherlanguage*}{english}
  <English abstract of the paper. Optional, but recommended.>

\textbf{Keywords}: <add some keywords> 
\end{otherlanguage*}

\cleardoublepage

\section*{Introduction}
\label{sec:intro1}
\phantomsection\addcontentsline{toc}{section}{\nameref{sec:intro1}}

%\defcitealias{keylist}{alias}

Due to their ecosystem services and economic values, forests play a predominant role on a global scale \citep{Balvanera2006Quantifyingevidence}. 
Within terrestrial ecosystems, they maintain significant biodiversity and act as regulators of biogeochemical factors \citep{Pawson2013Plantationforests}. 
However, the reality of climate change poses urgent challenges for the sustainability of current forests \citep{McKenney2009Climatechange,Messier2022Warningnatural,Seidl2017Forestdisturbances,Trumbore2015Foresthealth}. 
Despite international commitments to reduce greenhouse gas emissions, climate projections predict a global temperature exceeding 1.5 \up{o}C above pre-industrial temperatures \citep{Matthews2022Currentglobal}. 
Canada is particularly vulnerable to this warming, due to its northern latitudes \citep{Alo2008Potentialfuture,Bush2019Canadachanging} and forests in eastern North America will be particularly affected by this disruption \citep{Park2014Canboreal,Mahony2017closerlook,Messier2022Warningnatural,Sittaro2017Treerange}. 
Various studies have predicted lengthening and intensification of drought periods, an increase in wildfires and a higher presence of biotic disturbances \citep{Gatti2021Amazoniacarbon,Heidari2021Effectsclimate,Joyce2013Climatechange,Parmesan2007Influencesspecies}. 
Additionally, shifts in phenology and plant distribution are expected \citep{Aitken2008Adaptationmigration,Chuine2010Whydoes,Gray2013Trackingsuitable,Zhu2012Failuremigrate}. 
However, climate changes are occurring more rapidly than trees ability to adapt or migrate \citep{Aitken2008Adaptationmigration,Harrison2020Plantcommunity,Loarie2009velocityclimate,Messier2022Warningnatural,Williams2013Preparingclimate,Vitt2010Assistedmigration}, 
consequently threatening the growth and survival of these species \citep{Sittaro2017Treerange,Woodall2018Decadalchanges,Zhu2012Failuremigrate}.
This would ultimately result in a shift in forest composition, impacting forest management and conservation efforts \citep{Chmura2011Forestresponses,Lo2011Linkingclimate,McKenney2009Climatechange}.

Several calls for adaptation in forest management have been made to preserve forest ecosystems and their benefits \citep{Messier2021sakeresilience,Nagel2017Adaptivesilviculture}. 
Among the proposed solutions, assisted tree migration is suggested as a mitigation measure involving the movement of individuals or genetic material from their original climatic range to a more suitable area for species survival and growth in the future \citep{Dumroese2015Considerationsrestoring,Palik2022Operationalizingforestassisted,Park2023Provenancetrials,Park2018Informationunderload,Pedlar2011implementationassisted,Vitt2010Assistedmigration,Williams2013Preparingclimate}. 
This movement would rapidly change stand composition, meeting conservation needs, maintaining ecosystem services and preserving economic value \citep{Pedlar2011implementationassisted,Ste-Marie2011Assistedmigration,Winder2011Ecologicalimplications}. 
However, there remains a lack of knowledge and uncertainty surrounding assisted tree migration \citep{Park2018Informationunderload,Klenk2015assistedmigration}, particularly regarding the trade-offs between preserving one species and the risks to the ecosystem of the host territory \citep{Hewitt2011Takingstock,McLachlan2007frameworkdebate,Vitt2010Assistedmigration}.

To address this knowledge gap and reduce uncertainty, various forestry scenarios are currently under examination to mitigate risks \citep{royoDesiredREgenerationAssisted2023}. 
Therefore, the study of different silvicultural interventions such as overstory treatments is considered since they influence stand growth, health and composition \citep{Ameray2021Forestcarbon,Chaudhary2016Impactforest,Man2008Elevenyearresponses,MontoroGirona2018ConiferRegeneration,PamerleauCouture2015Effectthree}. 
For example, clear-cut treatments involve the removal of all trees within a designated area and is commonly used in intensive forest management plans focused on increasing wood productivity and quality over a short period to meet industry needs \citep{Ameray2021Forestcarbon}. 
On the other hand, partial-cut treatments entail selective tree removal, maintaining a portion of the stand and are usually applied in extensive management plans that favor natural regeneration, mimic natural disturbances and preserve the ecological value habitats \citep{Ameray2021Forestcarbon,Barg1999Influencepartial,Irland2011Timberproductivity,Tong2020Forestmanagement}. 
However, forest harvesting leads to changes in environmental conditions, such as soil compaction, increased solar exposure, higher winds and increased precipitation reaching the forest floor \citep{Keenan1993ecologicaleffects,Lindo2003Microbialbiomass,Heithecker2007Edgerelatedgradients}. 
Ultimately, these changes can affect nutrient availability and impact soil biodiversity \citep{Battigelli2004Shorttermimpact,Chaudhary2016Impactforest,Covington1981Changesforest,Fedrowitz2014Canretention,Kudrin2023metaanalysiseffects,Lindo2003Microbialbiomass,Paillet2010Biodiversitydifferences,rousseauLongtermEffectsBiomass2018}.
Soil fauna plays a crucial role in forest ecosystems by contributing to the circulation of matter and energy \citep{Kudrin2023metaanalysiseffects,Seibold2021contributioninsects}. 
Amphibians and arthropods are among the groups most often affected by environmental disturbances such as forestry practices \citep{Hartshorn2021reviewforest,Semlitsch2009Effectstimber,Stuart2004Statustrends} or climate change \citep{Alford1999Globalamphibian,Houlahan2000Quantitativeevidence,Milanovich2010Projectedloss,Parmesan2006EcologicalEvolutionary,Pounds2006Widespreadamphibian,Warren2018projectedeffect}.
Consequently, to better understand the impact of  treatments on soil fauna, we chose to study the Eastern Red-backed Salamander (\textit{Plethodon cinereus} (Green, 1818)), ground beetles (Carabidae) and springtails (Collembola).

The Eastern Red-backed Salamander is one of the most significant biomass among vertebrates in North American forests \citep{Burton1975Salamanderpopulations,Petranka1993Effectstimber,semlitschAbundanceBiomassProduction2014a}. 
Like other Plethodontidae, this salamander is strictly terrestrial and relies on skin respiration due to the absence of lungs \citep{Heatwole1961Relationsubstrate}. 
It occupies forest soils when temperature and humidity levels are optimal for cutaneous respiration. 
Outside these periods, it will move vertically in the soil to maintain favorable conditions to its survival \citep{Grizzell1949HibernationSite}. 
This species has a small home range and typically shows philopatric behavior \citep{Yurewicz2004ResourceAvailability}. 
Its role is significant in forest ecosystems as it acts as a generalist predator and regulates detritivore invertebrate populations \citep{Burton1975Energyflow,Hickerson2017Easternredbacked,Walton2013Topdownregulation}. 
It also serves as prey in trophic networks and constitutes a rich energy food source \citep{Burton1975Energyflow,Pough1987abundancesalamanders}. 
Due to its cutaneous respiration, this salamander is highly sensitive to environmental disturbances \citep{Welsh2001caseusing} and is commonly used as a bioindicator \citep{Baecher2018Environmentalgradients,gibbsDistributionWoodlandAmphibians1998,Heatwole1962EnvironmentalFactors,Harpole1999Effectsseven,Hocking2013Effectsexperimental,Mazerolle2021Woodlandsalamander}.

On their part, ground beetles gather the highest specific diversity among beetles with 40,000 identified species and represent one of the most abundant groups among soil arthropods \citep{Erwin1985taxonpulse,loveiEcologyBehaviorGround1996,Rochefort2006GroundBeetle}. 
As voracious carnivores and polyphagous predators, they act as regulators of invertebrate populations \citep{loveiEcologyBehaviorGround1996}. 
Ground beetles are also prey for several species of amphibians, reptiles, birds and mammals \citep{loveiEcologyBehaviorGround1996}. 
While widely distributed in terrestrial ecosystems, habitat selection varies among species \citep{Larochelle2003naturalhistory}. 
Ground beetles are often classified into three species communities : mature and closed forest species, open habitat species and generalist species \citep{Niemela2007effectsforestry}. 
This variation in habitat selection makes ground beetles an interesting taxon to study during environmental disturbances \citep{bouchardBeetleCommunityResponse2016b,Halme1993Carabidbeetles,Heliola2001Distributioncarabid,koivulaBorealCarabidbeetleColeoptera2002a,Luff1992Classificationprediction,Niemela2001Carabidbeetles,Rainio2003Groundbeetles,Work2008Evaluationcarabid}.

As for springtails, they are polyphyletic species of arthropods belonging to the mesofauna established in forest soils. 
These invertebrates have a high species richness and represent a significant abundance \citep{rusekBiodiversityCollembolaTheir1998}. 
Different springtail communities occupy a range of ecological niches from litter to various soil horizons \citep{pongeVerticalDistributionCollembola2000}.
The vertical distribution of these communities depends mainly on abiotic conditions such as light, humidity, or porosity. 
Therefore, springtails can be used to characterize a substrate based on the community found there \citep{rusekBiodiversityCollembolaTheir1998}. 
Primarily fungivore and detritivore, these organisms play a predominant ecological role by feeding largely on fungi, bacteria, actinomycetes and algae. 
They contribute to the decomposition of organic matter, nutrient transformation and energy transfer in terrestrial ecosystems \citep{Cuchta2019importantrole,Hattenschwiler2005Biodiversitylitter,Marsden2020Howagroforestry,Petersen2000Collembolapopulations,rusekBiodiversityCollembolaTheir1998,Wolters1991SoilInvertebrates}. 
Springtails also represent a food source for several species of arachnids, beetles, amphibians, reptiles and birds. 
This group of species is commonly used in studies focusing on the effects of environmental changes on forest soils and mesofauna \citep{farskaManagementIntensityAffects2014,rousseauWoodyBiomassRemoval2019,Salmon2008Relationshipssoil}.

These three taxa are relevant to study the impact of silvicultural treatments as their sensitivity to environmental changes and trophic relationships enable the analysis of disturbance effects on soil fauna dynamics.
Several studies have already studied the effects of those treatments on soil fauna.
However, most of them have only focused on the direct impacts of disturbances for one or more species group, neglecting the existing relationships between environmental variables and species groups \citep{josephIntegratingOccupancyModels2016,Kudrin2023metaanalysiseffects,Pollierer2021Diversityfunctional}. 
It is essential to have a better comprehension of those treatments effects, their propagation within the ecological network and their repercussions on soil organism communities.
Ultimately, this knowledge acquisition will provide useful tools to facilitate sustainable forest management.

Our study aims to understand how silvicultural practices, conducted in an assisted tree migration context, affect the dynamics of forest soil ecosystems. 
The objectives associated with this aim are firstly to measure the impact of overstory treatments on habitat use by soil fauna, and secondly to compare overstory 
treatments and environmental variables that influence habitat use by soil fauna.
Regarding our first objective, we hypothesized that overstory treatments result in a modification of habitat use by soil fauna and propagate 
through the trophic network. Specifically, the harvests initially affect habitat use by salamanders and large ground beetles, 
followed by modifications in habitat selection for small ground beetles, and ultimately affecting springtail biomass.
The hypothesis associated with our second objective suggested that environmental variables, favorable to habitat use by taxa, fluctuate according to 
the intensity of forest harvests. Therefore, overstory treatments can be used as a variable encompassing changes in environmental conditions.
As a result, more intense treatments would decrease litter depth and CWD, as a consequence of reduced accumulation of leaves and woody debris on the forest floor, 
and increase canopy openness. 


\section*{Material and methods}
\label{sec:matmet1}
\phantomsection\addcontentsline{toc}{section}{\nameref{sec:matmet1}}

\subsection*{Study area}
\label{subsec:area}
\phantomsection\addcontentsline{toc}{subsection}{\nameref{subsec:area}}

\begin{otherlanguage*}{english}
  Our study was conducted within the Portneuf Wildlife Reserve in the Captiale-Nationale administrative region near Lac des Amanites and Rivière-à-Pierre (47°07'N, 72°24'W, Figure \ref{fig:area}). 
  This area is located within the balsam fir (\textit{Abies balsamea})-yellow birch (\textit{Betula alleghaniensis}) bioclimatic domain \citep{saucierChapitreEcologieForestiere2009}.
  Other tree species include sugar maple (\textit{Acer saccharum}), red maple (\textit{Acer rubrum}), white spruce (\textit{Picea glauca}), black spruce (\textit{Picea mariana}), red spruce (\textit{Picea rubens}), white birch (\textit{Betula papyrifera}) and quaking aspen (\textit{Populus tremuloides}) \citep{olaBelowgroundCarbonStocks2024}. 
  Study sites rest on a deep glacial till as surface deposit with a moderately well-drained sandy loams soil \citep{CanadianSystemSoil1998}.
  The mean daily temperature is 4\up{o}C based on the 1981-2010 period at the nearest weather station (Lac aux sables, \citealp{environmentcanadaCanadianClimateNormals2019}). 
  Based on the same report, the mean annual precipitation and snowfall are 1133.2 mm and 230.3 cm, respectively.
  We used the assisted migration experimental system established in 2018 by the Ministère des Ressources naturelles et des Forêts to collect our data (\citealp{royoDesiredREgenerationAssisted2023}).
  This system is a factorial experimental design with a split-plots replicated in four blocks. 
  Each whole block (200 m x 140 m) is split in two overstory treatment : clear-cut and partial-cut. 
  See \cite{royoDesiredREgenerationAssisted2023} for more details about the assisted migration experiment.

\end{otherlanguage*}

\begin{figure}[ht!]
	\centering
	\includegraphics[scale=0.60]{fig_area4.png}
	\caption[Localization of the Captiale-Nationale administrative region in Quebec, Canada and position of the study area near Lac des Amanites in Portneuf Wildlife Reserve, Quebec, Canada.]
  {Localization of the Captial-Nationale administrative region in Quebec, Canada (A) and position of the study area near Lac des Amanites in Portneuf Wildlife Reserve, Quebec, Canada (B) where the assisted migration experimental system was implemented in 2018 (47°07'N, 72°24'W).}
	\label{fig:area}
	\end{figure}  



\subsection*{Sampling design}
\label{subsec:sampling}
\phantomsection\addcontentsline{toc}{subsection}{\nameref{subsec:sampling}}


We selected a total of 60 sampling units measuring 10 meters by 7.5 meters each to collect our data: we used four blocks serving as replicates, 
with each block containing six sampling units for both the clear-cut and partial-cut overstory treatments, 
while three additional sampling units were positioned outside the block, serving as controls (Figure \ref*{fig:blockSU}).
Sampling units serving as controls were separated from blocks by at least 10 meters to remove treatments effects.
In each sampling unit, we used three sampling methods to collect species data: artificial coverboards, pitfall traps, and soil cores. 

Artificial coverboards is commonly employed to count salamanders \citep{hesedUncoveringSalamanderEcology2012,Mazerolle2021Woodlandsalamander,mooreComparisonPopulationEastern2009c}. 
This method helps to standardize the number and size of sampled ground objects under which animals will hide, thereby reducing variability \citep{hydeSamplingPlethodontidSalamanders2001}. 
Coverboards were made of untreated spruce wood, measured 25 cm x 30 cm x 5 cm and were placed directly on the ground without litter underneath \citep{Mazerolle2021Woodlandsalamander}. 
Six coverboards were set per sample unit and spaced by at least 2.5 m, resulting in a total of 360 coverboards.
Two rows of three boards each were aligned along the length of the sampling units and centrally positioned (Figure \ref{fig:blockSU}).
All coverboards were placed outdoors in March 2022 to allow for natural aging, increasing the likelihood of salamander usage \citep{hedrickEffectsCoverboardAge2021,smithEffectsCoverBoard2014a}.
Throughout the surveys, coverboards were inspected once on the same day, and salamanders were counted without any manipulation.

Pitfall traps were used to capture ground beetles.
This method is commonly employed to assess the abundance and species richness of soil invertebrates, such as ground beetles \citep{baarsCatchesPitfallTraps1979,knappEffectPitfallTrap2012,kotzeFortyYearsCarabid2011a,loveiEcologyBehaviorGround1996,spenceSamplingCarabidAssemblages1994a}. 
Pitfall traps were Multipher\up{\textregistered{}} traps and included a container with a diameter of 12.5 cm, a depth of 25 cm and a cover raised 4.5 cm above the trap 
to prevent debris and rain from filling the container \citep{bouchardBeetleCommunityResponse2016b,Jobin1988MultiPherinsect,mooreEffectsTwoSilvicultural2004}.
They were equipped with a protective grid with a mesh size of 15 mm, limiting trap access to carabid-sized individuals and reducing the chances of predation by small mammals. 
Typically, a preserving liquid (such as propylene glycol or alcohol) is placed in the bottom of the container to preserve captured individuals. 
In this study, to avoid killing salamanders that are small enough to pass through the grid, we decided to use dry pitfall traps without any preserving liquid \citep{luffFeaturesInfluencingEfficiency1975}. 
We added wet sponges to the bottom of each container to maintain a suitable level of humidity for salamanders.
We centered one pitfall in each sampling and control units, resulting in a total of 60 pitfall traps (Figure \ref{fig:blockSU}). 
Traps were inserted in the soil at a depth allowing the container's opening to be juxtaposed with the soil surface. 
During periods when the traps were not visited daily, all pitfall traps were closed with adhesive tape around the opening to prevent individuals capture. 
On the first day of each survey, the traps were opened and captured ground beetles were collected for each remaining day for a five day period. 
Individuals were preserved in 70\% alcohol and identified at the species level afterwards.
Identification was conducted with a ZEISS SteREO Discovery.V12 bionocular microscope using \cite{larochelleManuelIdentificationCarabidae1976} identification keys.
Ground beetles were categorized in two groups as salamander prey or competitors based on the salamander gape size (Table \ref{tab:carabid}, \citealp{jaegerFoodLimitedResource1972,magliaModulationPreycaptureBehavior1995,magliaOntogenyFeedingEcology1996}).

Soil cores was used to sample springtails and are commonly employed to assess mesofauna in litter and different soil horizons \citep{chauvatChangesSoilFaunal2011a,farskaManagementIntensityAffects2014,pongeVerticalDistributionCollembola2000,salamonEffectsPlantDiversity2004,wuCompositionSpatiotemporalVariation2014}. 
Two soil cores with litter collection were harvested per sampling unit per survey using a soil sampling pedological probe. 
Cores had a diameter of 5 cm, a depth of 5 cm and a 15 cm x 15 cm litter quadrat was collected above each soil sample \citep{raymond-leonardSpringtailCommunityStructure2018a,rousseauForestFloorMesofauna2018}.
Both substrates were used to target mesofauna and obtain springtail communities directly related to the ecology of salamanders and ground beetles \citep{chauvatChangesSoilFaunal2011a,edwardsAssessmentPopulationsSoilinhabiting1991,raymond-leonardSpringtailCommunityStructure2018a,rousseauForestFloorMesofauna2018}.
Soil and litter from the same unit were pooled in Ziploc\up{\texttrademark{}} bags and stored in a cooler at $\pm$ 4 °C \citep{chauvatChangesSoilFaunal2011a,rousseauForestFloorMesofauna2018}, providing 60 samples per survey.
Each sample was placed in an individual Tullgren dry-funnel for springtail extraction within 48 h after they were collected \citep{rousseauForestFloorMesofauna2018,rusekBiodiversityCollembolaTheir1998,wuCompositionSpatiotemporalVariation2014}. 
The extraction process lasted six days with a gradual temperature increase (25 °C to 50 °C) \citep{raymond-leonardSpringtailCommunityStructure2018a}.
Springtails were preserved in 75\% alcohol \citep{wuCompositionSpatiotemporalVariation2014} prior to isolation from surrounding organisms, with subsequent identification at the family level.
Identification was done with a ZEISS SteREO Discovery.V12 binocular microscope and a Leitz orthoplan phase-contrast fluorescent trinocular microscope using \cite{bellingerChecklistCollembolaWorld1996} identification keys.
Springtail dry biomass of each sampled was measured with micro balance (Sartorius Cubis\up{\texttrademark{}} MSA3.6P-000-DM, city, state, country) after being lyophilized (Labconco FreeZone Bulk tray dryer 78060 series, city, state, country).

We conducted four surveys of five consecutive days each, namely in mid-May, mid-June, mid-July, and mid-August, during the summer 2022.
Blocks were visited in a random sequence to reduce effects of time of day and observer fatigue.

\pagebreak

\begin{figure}[ht]
	\centering
	\includegraphics[scale=0.50]{fig_blockSU2.png}
	\caption[Design of one block and one sampling unit with three sampling methods.]{
  Design of a block (left) and a sampling unit (right). 
  The block contains two overstory treatments : clear-cut (grey background), partial-cut (white background). 
  Fifteen sampling units were used per block : six per overstory treatment and three controls (\textbf{c}) outside each block.
  Each sampling unit contained six artificial coverboards (rhombus) and one pitfall trap (triangles). Two soil cores (circles) were collected per survey.
  }
	\label{fig:blockSU}
	\end{figure}  

\vspace{0.5cm}


\subsection*{Environmental variables}
\label{subsec:EnvVar}
\phantomsection\addcontentsline{toc}{subsection}{\nameref{subsec:EnvVar}}

In each sampling unit, we measured several environmental variables that could affect occupancy probability.
CWD and litter depth play a crucial role in habitat use for salamanders, ground beetles and springtails as
they serve for feeding and protection  \citep{birdChangesSoilLitter2004,groverInfluenceCoverMoisture1998a,harmonEcologyCoarseWoody1986,koivula.LeafLitterSmallscale1999,mckennyEffectsStructuralComplexity2006,patrickEffectsExperimentalForestry2006a}. \\
Salamanders also utilize CWD as shelter to maintain suitable temperature and moisture levels during dry periods \citep{groverInfluenceCoverMoisture1998a,Jaeger1980MicrohabitatsTerrestrial,patrickEffectsExperimentalForestry2006a}
We used 400 m\up2 plots centered inside every sampling unit to estimate CWD (20 m $\times$  20 m)(\citealp{methotGuideInventaireEchantillonnage2014}). 
We only selected CWD with a basal diameter greater than or equal to 9 cm and a length greater than or equal to 1 m.
Subsequently, we measured the basal diameter, the apical diameter, and length of CWD with a tree caliper.
Segments of CWD outside the plot boundaries were not considered.
We employed the conic–paraboloid formula to estimate the volume of each CWD \citep{fraverRefiningVolumeEstimates2007} :

\begin{equation}
  \text{Volume} = L/12 \times (5A_b + 5A_u + 2\sqrt{A_b \times A_u})
\end{equation}

\vspace{0.5cm}

Where $L$ is the length of log (cm), $A_b$ the basal area (cm\up{2}), and $A_u$ the apical area (cm\up{2}).
We measured litter depth next to each coverboard and estimated the mean depth per sampling unit \citep{Mazerolle2021Woodlandsalamander}. \\
We also assessed canopy openness, as it may influence species occupancy \citep{henneronForestPlantCommunity2017,koivulaBorealCarabidbeetleColeoptera2002a,kotzeFortyYearsCarabid2011a,messereForestFloorDistribution1998,tilghmanMetaanalysisEffectsCanopy2012}.
Measurements were conducted at the center of all sampling units using a spherical densiometer \citep{lemmonSphericalDensiometerEstimating1956} and were taken 130 cm above the ground. 
We took four measurements per sampling unit, oriented toward each of the four cardinal points, and computed the mean as an estimate of canopy openness in each sampling unit.

We collected data for air temperature, air humidity and precipitation levels during the summer 2022.
Two compact weather stations (Em50 Digital Decagon Data Logger, Part \#40800, Meter Group Inc., USA), were used inside both overstory treatments.
Each weather station was equipped with a probe measuring temperature, air humidity, and atmospheric pressure, 1.30 m above the ground (VP-4 Sensor (Temp/RH/Barometer), Part \#40023). 
Rain gauges were installed in the clear-cut treatments to monitor precipitation levels.
The temperature and humidity sensors were programmed for record data every 15 minutes. 
We used the means of both weather stations to get daily average measurements.
These variables fluctuate on a daily basis, affecting species activity and, consequently, the probability of detecting individuals. 
\citep{butterfieldCarabidLifeCycle1996,kotzeFortyYearsCarabid2011a,loveiEcologyBehaviorGround1996,odonnellPredictingVariationMicrohabitat2014a,spotilaRoleTemperatureWater1972}.

\subsection*{Statistical analyses}
\label{subsec:analyses}
\phantomsection\addcontentsline{toc}{subsection}{\nameref{subsec:analyses}} 

% Expliquer pourquoi d'un coté on regarde les variables emvironnementales et pourquoi de l'utre on regarde les espèces
% POurquoi ne pas intégré les varable du sol avec les espèces
% Expliquer comment les les modèles d'occupation et les GLM sont inclus dans le SEM comment il sont intgrer
% expliquer dans les methodes les analyses de comparaison entre coupe forestière et variables environnementales
% préciser que l'enseble de test sont fait avec le SEM

\subsubsection{Structural equations models} 

To assess the effects of overstory treatments on the habitat use (hypothesis 1.1) and the relationship between overstory treatments and environmental variables (hypothesis 2.1), 
we employed a structural equation models (SEM) combining occupancy models and linear mixed models (LMM) \citep{graceSpecificationStructuralEquation2010,josephIntegratingOccupancyModels2016,mackenzieOccupancyEstimationModeling2006a}.
This approach enables us to test both hypotheses within the same analysis using the same data and method, thereby reducing estimation variations.
One part of the SEM was designed to evaluate the direct and indirect effects of overstory treatments on taxa (Figure \ref*{fig:SEM}), 
while another section focused on studying the variations in environmental variables across different overstory treatments. 
SEM are employed to analyze complex relationships among observed and latent variables (unobserved but inferred from observed data) to investigate multicausal ecological processes \citep{graceStructuralEquationModeling2008}.

We combined occupancy models with SEM to measure the impact of overstory treatments on the occupancy probabilities of salamanders and ground beetles. 
Occupancy models allow accounting for imperfect detection since salamanders and ground beetles are cryptic species \citep{baileyEstimatingSiteOccupancy2004,spiersEstimatingSpeciesMisclassification2022}.
This method use repeated surveys to estimate the presence probabilities of species and the probability of detecting them if they are present, 
leading to an estimation of the species true occupancy \citep{mackenzieEstimatingSiteOccupancy2002,mazerolleMakingGreatLeaps2007}.

By combining SEM and occupancy models, the occupancy probability estimated from the occupancy models becomes a latent variable in the SEM. 
The SEM then explores how the explanatory variables directly and indirectly affect this latent occupancy variable.
This approach enables us to investigate complex relationships influencing the presence of salamanders and ground beetles 
while addressing the issue of imperfect detection.

LMM were incorporated into the SEM to estimate parameters of springtail biomass and environmental variables.
LMM allows us to account for random variations of blocks, enhancing model accuracy. 
Those models can handle data with hierarchical structures and can estimate the amount of variation due to fixed effects while also accounting for 
random effects.




\begin{figure}[ht]
	\centering
	\includegraphics[scale=0.55]{fig_sem.png}
	\caption[Theoretical model illustrating the anticipated relationships between overstory treatments, environmental variables and species groups.]
  {Theoretical model illustrating the anticipated relationships between overstory treatments, coarse woody debris volume, canopy openness, litter depth,
   salamander occupancy, ground beetle occupancy, and springtail biomass in the Portneuf Wildlife Reserve, Quebec, Canada. 
   Each arrow indicates the direction of a potential effect, from an explanatory variable to a response variable.}
	\label{fig:SEM}
\end{figure}  

\subsubsection{Occupancy models} 


To estimate the occupancy probabilities of salamanders and ground beetles, we used the observed detection at site $i$ during survey $j$, 
represented by $\text{Salamander}_{ij}$, $\text{Carabid.comp}_{ij}$ and $\text{Carabid.prey}_{ij}$. Observed detection followed a Bernoulli distribution with $z_{i} \times p_{ij}$ as parameter, 
where $z_{i}$ represents the latent occupancy state of species groups at a site $i$ and $p_{ij}$ represents the probability of detecting species groups at a site $i$ during a survey $j$ :


\begin{align}
  \text{Salamander}_{ij} &\sim \text{Bernoulli}(z_{\text{Salamander}_i} \times p_{\text{Salamander}_{ij}}) \nonumber \\
  \text{Carabid.comp}_{ij} &\sim \text{Bernoulli}(z_{\text{Carabid.comp}_i} \times p_{\text{Carabid.comp}_{ij}})  \\
  \text{Carabid.prey}_{ij} &\sim \text{Bernoulli}(z_{\text{Carabid.prey}_i} \times p_{\text{Carabid.prey}_{ij}}) \nonumber
\end{align}


$z_{i}$ follows a Bernoulli distribution with $\psi$ as parameter, 
which represent the occupancy probability of each species group. 
The latent occupancy state $z$ indicates whether a site is occupied by a species group ($z = 1$) or not ($z = 0$) :


\begin{align}
  z_{\text{Salamander}_i} &\sim \text{Bernoulli}(\psi_{\text{Salamander}_{\text{Group}_i}}) \nonumber \\
  z_{\text{Carabid.comp}_i} &\sim \text{Bernoulli}(\psi_{\text{Carabid.comp}_{\text{Group}_i}}) \\
  z_{\text{Carabid.prey}_i} &\sim \text{Bernoulli}(\psi_{\text{Carabid.prey}_{i}}) \nonumber
\end{align}


Salamanders and large ground beetle occupancy probabilities ($\psi_{\text{Salamander}_{\text{Group}_i}}$, $\psi_{\text{Carabid.comp}_{\text{Group}_i}}$ ) are drawn from uniform distribution $\text{U}(0, 1)$ for each overstory group $i$, 
whereas small ground beetle occupancy probabilities ($\psi_{\text{Carabid.prey}_{i}}$) were estimated from a linear predictor with salamander latent occupancy state ($z_{\text{Salamander}_i}$) and overstory treatments 
($\text{Cutpartial}_i$, $\text{Cutclear}_i$) as explanatory variables on a logit scale. 
$\beta$s represent the coefficients associated with each parameter affecting occupancy in the linear predictor. 
We assumed vague normal priors, $\text{N}(0, \sigma = 10)$ for those $\beta$ parameters :


\begin{align}
  \text{logit}(\psi_{\text{Carabid.prey}_i}) &= \beta_{0[\text{Carabid.prey}]} + \beta_{z_{\text{Salamander}}[\text{Carabid.prey}]} \times z_{\text{Salamander}_i} + \nonumber \\
  &\beta_{\text{Cutpartial}[\text{Carabid.prey}]} \times \text{Cutpartial}_i + \\
  &\beta_{\text{Cutclear}[\text{Carabid.prey}]} \times \text{Cutclear}_i \nonumber
\end{align}

We applied a logit scale on parameters influencing the probability of detecting species groups ($\text{logit}(p_{ij})$), allowing us to use 
the coarse woody debris ($\text{CWD}_i$), the precipitation levels ($\text{Precipitation}_{ij}$) as explanatory variables with a block random effect ($\alpha_{Block}$). 
$\alpha$s represent the coefficients associated with each explanatory variable affecting detection in the linear predictor. 
We assumed vague normal priors for CWD, precipitation levels $\text{N}(0, \sigma = 10)$ and block random effects $\text{N}(0, \sigma_{Block})$ 
where $\sigma_{Block}$ is drawn form uniform distribution $\text{U}(1, 10)$. 
Air temperature and relative humidity were excluded from detection variables due to a high correlation with the precipitation level.


\begin{align}
  \text{logit}(p_{\text{Salamander}_{ij}}) &= \alpha_{0[\text{Salamander}]} + \alpha_{\text{CWD}[\text{Salamander}]} \times \text{CWD}_i + \nonumber \\
  &\alpha_{\text{Precipitation}[\text{Salamander}]} \times \text{Precipitation}_{ij} + \alpha_{\text{Block}[\text{Salamander}]_{\text{Block}_i}} \nonumber
\end{align}

\begin{align}
  \text{logit}(p_{\text{Carabid.comp}_{ij}}) &= \alpha_{0[\text{Carabid.comp}]} + \alpha_{\text{CWD}[\text{Carabid.comp}]} \times \text{CWD}_i + \\
  &\alpha_{\text{Precipitation}[\text{Carabid.comp}]} \times \text{Precipitation}_{ij} + \alpha_{\text{Block}[\text{Carabid.comp}]_{\text{Block}_i}} \nonumber 
\end{align}

\begin{align}
  \text{logit}(p_{\text{Carabid.prey}_{ij}}) &= \alpha_{0[\text{Carabid.prey}]} + \alpha_{\text{CWD}[\text{Carabid.prey}]} \times \text{CWD}_i + \nonumber \\
  &\alpha_{\text{Precipitation}[\text{Carabid.prey}]} \times \text{Precipitation}_{ij} + \alpha_{\text{Block}[\text{Carabid.prey}]_{\text{Block}_i}} \nonumber 
\end{align}

%%%%%%%%%%%%%%%%%%%
%%%%%%%%%%%%%%%%%%% fin


% \begin{align}
%   z_{\text{Salamander}_i} &\sim \text{Bernoulli}(\psi_{\text{Salamander}_{\text{Group}_i}}) \nonumber \\
%   \text{logit}(p_{\text{Salamander}_{ij}}) &= 
%   \alpha_{0[\text{Salamander}]} + \alpha_{\text{CWD}[\text{Salamander}]} \times \text{CWD}_i + \\
%   &\alpha_{\text{Precipitation}[\text{Salamander}]} \times \text{Precipitation}_{ij} + \alpha_{\text{Block}[\text{Salamander}]_{\text{Block}_i}} \nonumber \\
%   \text{Salamander}_{ij} &\sim \text{Bernoulli}(z_{\text{Salamander}_i} \times p_{\text{Salamander}_{ij}}) \nonumber
% \end{align} 

% \begin{align}
% z_{\text{Salamander}_i} &\sim 
% \text{Bernoulli}(\psi_{\text{Salamander}_{\text{Group}_i}}) \nonumber \\
% \text{logit}(p_{\text{Salamander}_{ij}}) &= 
% \alpha_{0[\text{Salamander}]} + \alpha_{\text{CWD}[\text{Salamander}]} \times \text{CWD}_i + \\
% &\alpha_{\text{Precipitation}[\text{Salamander}]} \times \text{Precipitation}_{ij} + \alpha_{\text{Block}[\text{Salamander}]_{\text{Block}_i}} \nonumber \\
% \text{Salamander}_{ij} &\sim 
% \text{Bernoulli}(z_{\text{Salamander}_i} \times p_{\text{Salamander}_{ij}}) \nonumber
% \end{align} 

% \begin{align}
%   z_{\text{Carabid.comp}_i} &\sim 
%   \text{Bernoulli}(\psi_{\text{Carabid.comp}_{\text{Group}_i}}) \nonumber \\
%   \text{logit}(p_{\text{Carabid.comp}_{ij}}) &= 
%   \alpha_{0[\text{Carabid.comp}]} + \alpha_{\text{CWD}[\text{Carabid.comp}]} \times \text{CWD}_i + \\
%   &\alpha_{\text{Precipitation}[\text{Carabid.comp}]} \times \text{Precipitation}_{ij} + \alpha_{\text{Block}[\text{Carabid.comp}]_{\text{Block}_i}} \nonumber \\
%   \text{Carabid.comp}_{ij} &\sim \text{Bernoulli}(z_{\text{Carabid.comp}_i} \times p_{\text{Carabid.comp}_{ij}}) \nonumber
%   \end{align}

% \begin{align}
%   \text{logit}(\psi_{\text{Carabid.prey}_i}) &= 
%   \beta_{0[\text{Carabid.prey}]} + \beta_{z_{\text{Salamander}}[\text{Carabid.prey}]} \times z_{\text{Salamander}_i} + \nonumber \\
%   &\beta_{\text{Cutpartial}[\text{Carabid.prey}]} \times \text{Cutpartial}_i + \beta_{\text{Cutclear}[\text{Carabid.prey}]} \times \text{Cutclear}_i \nonumber\\
%   z_{\text{Carabid.prey}_i} &\sim 
%   \text{Bernoulli}(\psi_{\text{Carabid.prey}_i}) \nonumber \\
%   \text{logit}(p_{\text{Carabid.prey}_{ij}}) &= \alpha_{0[\text{Carabid.prey}]} + \alpha_{\text{CWD}[\text{Carabid.prey}]} \times \text{CWD}_i +  \\
%   &\alpha_{\text{Precipitation}[\text{Carabid.prey}]} \times \text{Precipitation}_{ij} + \alpha_{\text{Block}[\text{Carabid.prey}]_{\text{Block}_i}} \nonumber \\
%   \text{Carabid.prey}_{ij} &\sim \text{Bernoulli}(z_{\text{Carabid.prey}_i} \times p_{\text{Carabid.prey}_{ij}}) \nonumber
% \end{align}

\subsubsection{Linear mixed models} 

We used linear mixed models to assess how overstory treatments affect springtail biomass ($\text{Springtail}_{i}$) and 
environmental variables ($\text{CWD}_{i}$, $\text{Canopy}_{i}$, $\text{Litter}_{i}$) at site $i$ :

\begin{align}
  \text{Springtail}_{i} &\sim \text{N} (\mu_{\text{Springtail}_i}, \sigma_{\text{Springtail}_{\text{Group}_i}}) \nonumber \\
  \text{CWD}_{i} &\sim \text{N} (\mu_{\text{CWD}_i}, \sigma_{\text{CWD}_{\text{Group}_i}}) \\
  \text{Canopy}_{i} &\sim \text{N} (\mu_{\text{Canopy}_i}, \sigma_{\text{Canopy}_{\text{Group}_i}}) \nonumber \\
  \text{Litter}_{i} &\sim \text{N} (\mu_{\text{Litter}_i}, \sigma_{\text{Litter}_{i}}) \nonumber 
\end{align}

Springtail biomass and environmental variables estimations are drawn from Normal distributions where means estimations at site $i$ ($\mu_{i}$) included overstory treatments ($\text{Cutpartial}_i$, $\text{Cutclear}_i$) and block random effect ($\alpha_{\text{Block}}$) as linear predictors. 
Estimations of $\mu_{\text{Springtail}_i}$ also accounted for latent occupancy state of salamanders and both ground beetles groups ($z_{Salamander}$, $z_{SCarabid.comp}$, $z_{Carabid.prey}$ ). 
We assumed vague normal priors for overstory treatments, latent occupancy states, $\text{N}(0, \sigma = 10)$ and block random effects $\text{N}(0, \sigma_{Block})$ 
where $\sigma_{Block}$ is drawn form uniform distribution $\text{U}(0, 50)$. 
The $\beta$ coefficients represent the weights assigned to each explanatory in the linear predictor.
Due to heteroscedasticity with $\text{Springtail}$, $\text{CWD}$ and $\text{Canopy}$, we allowed each overstory groups $j$ to have their own variances ($\sigma_j \sim \text{U}(0,150)$) :

\begin{align}
  \mu_{\text{Springtail}_i} &= \beta_{0[\text{Springtail}]} + \beta_{\text{Cutpartial}[\text{Springtail}]} \times \text{Cutpartial}_i + \nonumber\\
  &\beta_{\text{Cutclear}[Springtail]} \times \text{Cutclear}_i + \beta_{z_{\text{Salamander}}[\text{Springtail}]} \times z_{Salamander} +  \nonumber\\
  &\beta_{z_{\text{Carabid.prey}}[\text{Springtail}]} \times z_{Carabid.prey} + \beta_{z_{\text{Carabid.comp}}[\text{Springtail}]} \times z_{Carabid.comp} + \nonumber\\
  &\alpha_{\text{Block}[\text{Springtail}]_{\text{Block}_i}} \nonumber
\end{align}

\begin{align}
  \mu_{\text{CWD}_i} &= \beta_{0[\text{CWD}]} + \beta_{\text{Cutpartial}[\text{CWD}]} \times \text{Cutpartial}_{i} + \nonumber\\
  & \beta_{\text{Cutclear}[\text{CWD}]} \times \text{Cutclear}_{i} + \alpha_{\text{Block}[\text{CWD}]_{\text{block}_i}} 
\end{align}


\begin{align}
  \mu_{\text{Canopy}_i} &= \beta_{0[\text{Canopy}]} + \beta_{\text{Cutpartial}[\text{Canopy}]} \times \text{Cutpartial}_{i} + \nonumber \\
  & \beta_{\text{Cutclear}[\text{Canopy}]} \times \text{Cutclear}_{i} + \alpha_{\text{Block}[\text{Canopy}]_{\text{block}_i}} \nonumber
\end{align}

\begin{align}
  \mu_{\text{Litter}_i} &= \beta_{0[\text{Litter}]} + \beta_{\text{Cutpartial}[\text{Litter}]} \times \text{Cutpartial}_{i} + \nonumber\\
  & \beta_{\text{Cutclear}[\text{Litter}]} \times \text{Cutclear}_{i} + \alpha_{\text{Block}[\text{Litter}]_{\text{block}_i}} \nonumber
\end{align}


Analysis was performed with a Bayesian approach and parameters were estimated using Markov chain Monte Carlo (MCMC) with JAGS 4.3.0 include in the jagsUI package in R 4.3.1 \citep{lunnBUGSProjectEvolution2009,kellnerJagsUIWrapperRjags2024,rcoreteamLanguageEnvironmentStatistical2020}. 
We ran the model with five chains, 200,000 iterations each \citep{gelmanUnderstandingPredictiveInformation2014}. A burn-in period of 75,000 iterations was used and a thinning rate of 5 was applied. 
We verified the convergence of MCMC chains by examining trace plots, posterior density plots, and applying the Brooks-Gelman-Rubin statistic. 
The JAGS model code is available in Table

\clearpage

\section*{Results}
\label{sec:results1}
\phantomsection\addcontentsline{toc}{section}{\nameref{sec:results1}}


\subsection*{Soil fauna}
\label{subsec:taxa}
\phantomsection\addcontentsline{toc}{subsection}{\nameref{subsec:taxa}} 

Model diagnostics indicated that the chain lengths were sufficient, as the Brooks-Gelman-Rubin statistic was below 1.1. 
Trace plot analysis revealed that all chains had converged towards similar values, and none of the ratios of MCMC error to posterior standard deviation exceeded 5\%.

\vspace{10pt}

\begin{figure}[ht]
	\centering
	\includegraphics[scale=0.60]{fig_sem_res3.png}
	\caption[Results from structural equation modeling analysis revealing effects of overstory treatments on coarse woody debris volume,
  canopy openness, litter depth, salamander occupancy, carabid occupancy, and springtail biomass.]
  {Results from SEM analysis showing effects of overstory treatments on CWD, 
  canopy openness, litter depth, salamander occupancy, carabid occupancy, and springtail biomass in the Portneuf Wildlife Reserve, 
  Quebec, Canada. Bold arrows represent significant effects, while gray arrows indicate no discernible effects. 
  Values above bold arrows represent average differences between posterior distributions of two overstory groups: 
  partial-cut (PC), clear-cut (CC), and control (C). Estimates marked with one asterisk (*) 
  indicate a 90\% credible interval (CI) excluding 0, while estimates marked with two asterisks (**) indicate a 95\% CI excluding 0.}
	\label{fig:SEMres}
\end{figure}  

\vspace{10pt}

The average salamander observation per survey was 8 individuals with a range of 17 salamanders, during the summer 2022.
We captured 189 ground beetles belonging to 30 species, with 49 ground beetles found in the partial-cut treatments, 105 in the clear-cuts, and 35 in the control sites (Table \ref{tab:carabid}). 
We collected 468 springtails representing 12 families, with 219 springtails collected from partial-cut treatments, 131 from clear-cuts, and 118 from the control areas (Table \ref{tab:springtail}). 
The average springtail biomass collected per overstory treatments was 24.3 $\mu$g (SD = 18.2 $\mu$g), 56.8 $\mu$g (SD = 78.0 $\mu$g), and 31.1 $\mu$g (SD = 52.8 $\mu$g) in the partial-cut treatments, clear-cuts and control sites, respectively.

The effects of overstory treatments on habitat selection by species groups were generally not significant. 
Salamander occupancy probability was marginally lower in sites subjected to clear-cutting compared to those with partial-cutting (90\% CI : [-0.74, -0.07], Figure \ref{fig:pcin}, Table \ref{tab:overstorysp}). 
However, these two groups did not differ from the control sites. 
Occupancy probability for both carabid groups and the springtail biomass, did not vary significantly between the overstory treatments and control sites (Table \ref{tab:overstorysp}). 
Overall, the presence of salamanders had no significant impact on the occupancy probabilities of ground beetles representing prey for salamanders, and we did not observe significant effects of either salamanders or ground beetles on springtail biomass (Table \ref{tab:overstorysp}).

\vspace{10pt}

\begin{table}[ht]
  \centering
  \caption[Differences between overstory treatments on the forest soil fauna.]
  {Differences between overstory treatments on salamander occupancy, both ground beetle groups occupancy and springtail biomass. 
  This table also show estimated effect of salamanders presence on ground beetles form the salamanders prey group occupancy 
  and the effects of salamanders and both ground beetle groups presence on springtail biomass, during the summer 2022 in the Portneuf Wildlife Reserve, Quebec, Canada.}
  \label{tab:overstorysp}
  \begin{tabular}{lllll} 
      \hline
      &&&&95\% Bayesian \\
      Variable&Unit& Comparison & Estimate &  credible interval \\ [0.5ex] 
      \hline     
      Salamander           &probability& Partial vs control & \hspace{1mm}0.07 & [-0.29, 0.45] \\ 
      occupancy       && Clear vs control  & -0.38 & [-0.75, 0.11] \\ 
                          && Clear vs partial  & -0.45 & [-0.74, -0.07]$^{a}$ \\       
      Carabid$_{competitor}$ &probability& Partial vs control & -0.12 & [-0.35, 0.15] \\
      occupancy       && Clear vs control  & -0.06 & [-0.29, 0.20] \\ 
                          && Clear vs partial  & \hspace{1mm}0.06 & [-0.19, 0.30] \\ 
      Carabid$_{prey}$    &logit& Partial vs control & \hspace{1mm}3.31 & [-10.12, 17.72] \\
      occupancy             && Clear vs control  & \hspace{1mm}10.19 & [-4.15, 24.45] \\ 
                          && Clear vs partial  & \hspace{1mm}6.88 & [-12.81, 23.42] \\  
                          && Salamander        & -2.20 & [-17.15, 16.59] \\  
      Springtail          &$\mu$g& Partial vs control & \hspace{1mm}8.11 & [-9.38, 25.40] \\
      biomass             && Clear vs control  & \hspace{1mm}2.11 & [-13.98, 18.11] \\ 
                          && Clear vs partial  & -6.00 & [-29.09, 17.26] \\  
                          && Salamander        & \hspace{1mm}6.80 & [-10.43, 23.26] \\ 
                          && Carabid$_{competitor}$      & \hspace{1mm}0.56 & [-16.75, 17.66] \\ 
                          && Carabid$_{prey}$      & \hspace{1mm}7.62 & [-8.93, 24.09] \\ 
      \hline
      \multicolumn{5}{l}{\textbf{Note:} Estimates from Bayesian SEM are presented in terms of mean number with 95\%} \\
      \multicolumn{5}{l}{credible intervals, where an interval excluding 0 indicates a difference between groups.} \\
      \multicolumn{5}{l}{$^{a}$Marginal difference based on 90\% Bayesian credible interval excluding 0}
  \end{tabular}
\end{table}


\clearpage

\begin{figure}[ht]
  \centering
  \includegraphics[scale=0.55]{fig_pcin.png}
  \caption[Occupancy probability of salamanders under overstory treatments]
  {Occupancy probability of salamanders within two overstory treatments and controls during the summer 2022 in the Portneuf Wildlife Reserve device, Quebec, Canada. 
  Error bars denote 95\% credible intervals around estimates.}
  \label{fig:pcin}
\end{figure}

\vspace{10pt}

We did not observe significant impacts of CWD and precipitation level on salamander detection probabilities. 
However, the precipitation level had a positive effect on detection probability for both small ground beetles (95\% CI : [0.59, 1.77]) and large carabid (95\% CI : [0.70, 3.23]) (Table \ref{tab:detection}). 
CWD had no significant impact on ground beetle detection probabilities.


\begin{table}[ht]
  \centering
  \caption[Estimated effects of coarse woody debris and precipitation level on detection probabilities of salamanders and both ground beetle groups.]
  {Estimated effects of coarse woody debris and precipitation level on detection probabilities of salamanders and both ground beetle groups, during the summer 2022 in the Portneuf Wildlife Reserve,  Quebec, Canada.}
  \label{tab:detection}
  \begin{tabular}{lllll} 
      \hline
      &&&95\% Bayesian \\
      Variable & Taxa & Estimate &  credible interval \\ [0.5ex] 
      \hline      
      Precipitation       & Salamander              & \hspace{1mm}0.11 & [-0.83, 1.06] \\ 
                          & Carabid$_{competitor}$  & \hspace{1mm}1.17 & [0.59, 1.77] \\ 
                          & Carabid$_{prey}$        & \hspace{1mm}1.87 & [0.70, 3.23] \\  
      \hline      
      Coarse woody debris & Salamander              & -0.59 & [-1.39, 0.12] \\ 
                          & Carabid$_{competitor}$  & \hspace{1mm}0.06 & [-0.26, 0.38] \\ 
                          & Carabid$_{prey}$        & \hspace{1mm}0.27 & [-0.74, 1.37] \\   

      \hline
      \multicolumn{4}{l}{\textbf{Note:} Estimates from Bayesian SEM are presented in terms of mean} \\
      \multicolumn{4}{l}{number with 95\% credible intervals, where an interval excluding 0 indicates} \\
      \multicolumn{4}{l}{a difference between groups.} \\
  \end{tabular}
\end{table}


\subsection*{Environmental variables}
\label{subsec:ResEnv}
\phantomsection\addcontentsline{toc}{subsection}{\nameref{subsec:ResEnv}} 

Environmental variables usually differed between forest cutting treatments and control conditions. 
We found that clear-cutting treatments had significantly less CWD compared to partial-cutting (95\% CI : [-1.15, -0.43]) (Figure \ref{fig:envar} A, Table \ref{tab:overstoryenvar}). 
However, these both treatments did not differ from the control sites. 
Canopy openness was significantly higher in both the partial-cut (95\% CI : [1.97, 11.02]) and clear-cut treatments (95\% CI : [51.39, 77.06]) when compared 
to the control sites, with the clear-cuts having a greater percentage of openness than the partial-cuts (95\% CI : [44.61, 70.76], Figure \ref{fig:envar} B, Table \ref{tab:overstoryenvar}). 
Conversely, litter depth was lower in both the partial-cut (95\% CI : [-2.44, -0.65]) and clear-cut treatments (95\% CI : [-4.28, -2.50]) compared to the controls, 
with the depth being lower in the clear-cuts than in the partial-cuts (95\% CI : [-2.57, -1.12], Figure \ref{fig:envar} C, Table \ref{tab:overstoryenvar}).

\vspace{10pt}

\begin{figure}[ht]
  \centering
  \includegraphics[scale=0.23]{fig_envar2.png}
  \caption[Environmental variables estimations with a potential effect soil species within two different overstory treatments and control.]
  {Environmental variables with a potential effect soil species estimations within two different overstory treatments and control 
  during the summer 2022 in the Portneuf Wildlife Reserve, Quebec, Canada. Error bars denote 95\% credible intervals around estimates.}
  \label{fig:envar}
\end{figure}

\begin{table}[ht]
  \centering
  \caption[Differences between overstory treatments on environmental variables that could affect habitat selection of fauna on the forest soil.]
  {Differences between overstory treatments on environmental variables that could affect habitat selection of fauna on the forest soil during the summer 2022 in the Portneuf Wildlife Reserve,
  Quebec, Canada.}
  \label{tab:overstoryenvar}
  \begin{tabular}{lllll} 
      \hline
      &&&&95\% Bayesian \\
      Variable&Unit& Comparison & Estimate &  credible interval \\ [0.5ex] 
      \hline
      Coarse woody debris &m\up{3}& Partial vs control & \hspace{1mm}0.02 & [-1.01, 1.06] \\ 
                 && Clear vs control  & -0.77 & [-1.79, 0.23] \\ 
                          && Clear vs partial  & -0.79 & [-1.15, -0.43] \\
      Canopy openness     &\%& Partial vs control & \hspace{1mm}6.49 & [1.97, 11.02] \\ 
                      && Clear vs control  & \hspace{1mm}64.19 & [51.39, 77.06] \\ 
                          && Clear vs partial  & \hspace{1mm}57.69 & [44.61, 70.76] \\ 
      Litter depth        &cm& Partial vs control & -1.54 & [-2.44, -0.65] \\ 
                      && Clear vs control  & -3.39 & [-4.28, -2.50] \\ 
                          && Clear vs partial  & -1.85 & [-2.57, -1.12] \\       
      \hline
      \multicolumn{5}{l}{\textbf{Note:} Estimates from Bayesian SEM are presented in terms of mean number with 95\%} \\
      \multicolumn{5}{l}{credible intervals, where an interval excluding 0 indicates a difference between groups.} \\
  \end{tabular}
\end{table}




\clearpage

\section*{Discussion}
\label{sec:discu1}
\phantomsection\addcontentsline{toc}{section}{\nameref{sec:discu1}}

\section*{Conclusion}
\label{sec:conclu1}
\phantomsection\addcontentsline{toc}{section}{\nameref{sec:conclu1}}

\section*{Acknowledgements}
\label{sec:acknowl1}
\phantomsection\addcontentsline{toc}{section}{\nameref{sec:acknowl1}}

\section*{Conflict of interest}
\label{sec:conflict1}
\phantomsection\addcontentsline{toc}{section}{\nameref{sec:conflict1}}

None declared
\section*{Author contributions}
\label{sec:author1}
\phantomsection\addcontentsline{toc}{section}{\nameref{sec:author1}}

\cleardoublepage

\begin{otherlanguage}{english}
\bibliographystyle{ecologyNewEN} % Style de citation en français
\bibliography{References}
\addcontentsline{toc}{section}{References}
\end{otherlanguage}
    % chapitre 1
%<!articles>\include{chapitre2}             % chapitre 2, etc.
%<articles>\chapter{<Titre du chapitre ou de l'article>}     % numéroté
\label{chap-}                   % étiquette pour renvois (à compléter!)

\section{Résumé}

\begin{otherlanguage*}{french}
  <Résumé de l'article en français. Obligatoire.>
\end{otherlanguage*}

\section{Abstract}

\begin{otherlanguage*}{english}
  <English abstract of the paper. Optional, but recommended.>
\end{otherlanguage*}

<Texte du chapitre ou de l'article.>

\bibliographystyle{}              % style de la bibliographie
\bibliography{}                   % production de la bibliographie
    % chapitre 2, etc.
\chapter*{Conclusion générale}           % ne pas numéroter
\label{chap-conclusion}         % étiquette pour renvois
\phantomsection\addcontentsline{toc}{chapter}{\nameref{chap-conclusion}} % inclure dans TDM

Le but de cette étude était de comprendre comment les traitements sylvicoles, effectués dans un contexte de migration assistée, affectent la dynamique des écosystèmes du sol forestier.
Les objectifs qui s’y rattachaient étaient, d'une part, de quantifier l'effet des traitements de coupes forestières sur les variables environnementales qui influencent 
l'utilisation de l'habitat par la faune du sol et, d'autre part, de mesurer l'impact des coupes forestières sur l'occupation de cet habitat par la faune du sol.
Pour répondre à ces deux objectifs, nous avons développé un modèle d'équations structurelles pour mesurer l'effet des coupes sur le volume de débris ligneux, 
la profondeur de litière et l'ouverture de la canopée, ainsi que sur la probabilité d'occupation des salamandres cendrées, des carabes et sur la biomasse de collemboles.
Notre première hypothèse stipulait que les variables environnementales favorables à l'utilisation de l'habitat par les espèces fluctuent 
en fonction de l'intensité des coupes forestières et qu'ainsi, les traitements de coupes forestières constituent une variable englobant les changements de conditions environnementales. 
Notre seconde hypothèse soutenait que les traitements de coupe forestière entraînent une modification de l'utilisation de l'habitat par la faune du sol et se propagent à travers 
le réseau trophique des salamandres vers les collemboles, en passant par les carabes. 
Les principales conclusions tirées de cette étude sont :

\begin{enumerate}
  \item Les coupes forestières ont un effet significatif sur les variables environnementales influençant l'utilisation de l'habitat par la faune du sol. Ces changements suivent de manière générale le niveau d'intensité du traitement appliqué, avec des perturbations plus importantes dans les coupes totales, suivis des coupes partielles.
  \item La probabilité d'occupation des salamandres et des carabes ne semble globalement pas affectée par les différents traitements de coupes forestières comparativement aux sites témoins. Toutefois, un faible effet a été mesuré chez les salamandres où la probabilité d'occupation de celles-ci était marginalement plus faible dans les coupes totales par rapport aux coupes partielles.
  \item La biomasse des collemboles ne variait pas de façon significative entre les différents traitements sylvicoles.
  \item La relation de cooccurence entre les trois taxons ne semble pas subir de changement en fonction du type de coupe. Ainsi, la biomasse des collemboles n'est pas affectée par le changement de probabilité d'occupation des salamandres et des carabes, et la présence des salamandres ne modifie pas la probabilité d'occupation des carabes de petite taille.
\end{enumerate}


\subsection{Résultat, limitations et ouvertures}

Notre étude a révélé que les coupes forestières peuvent avoir un impact négatif sur certaines variables environnementales importantes pour la faune du sol. 
Ces effets sont particulièrement prononcés dans les traitements de coupe totale, où l'on observe une réduction du volume de débris ligneux, une augmentation significative de l'ouverture de la canopée et une diminution de la profondeur de la litière. 
L'étude a également permis d'observer que les traitements de coupe partielle permettent une meilleure rétention des attributs environnementaux, puisque les effets pour les trois variables mesurées étaient moindres dans ce type de traitement comparativement aux témoins.
Malgré les changements environnementaux observés, notre étude ne permet pas de conclure que les traitements de coupe forestière influencent directement ou indirectement l'occupation de l'habitat par la salamandre cendrée et les carabes ou la biomasse des collemboles. 
Notre hypothèse concernant la relation entre les taxons et le modèle d'équation structurelle qui en découle suggérait que l'effet des perturbations causées par les coupes forestières se propage à travers le réseau trophique des prédateurs vers les proies. 
Toutefois, la récolte forestière ne semble pas affecter la cooccurrence entre ces trois groupes d'espèces, selon ce type de relation. 

Une autre approche serait d'envisager que les perturbations environnementales affectent en premier lieu le bas de la chaîne alimentaire, avant de se propager à travers le réseau trophique jusqu'aux prédateurs \citep{Laigle2021Directindirect}. 
En effet, les espèces les plus sensibles aux perturbations sont généralement celles ayant une grande taille corporelle et un niveau trophique élevé \citep{Seibold2015Associationextinction,Nolte2019Habitatspecialization,Hagge2021Whatdoes}. 
En tant que consommateurs au sommet de la chaîne alimentaire au niveau du sol, les salamandres seraient ainsi particulièrement vulnérables aux changements dans les niveaux trophiques inférieurs. 
Un modèle alternatif d'équations structurelles aurait pu être élaboré pour examiner si les traitements sylvicoles peuvent influencer les collemboles, avec des effets se propageant ensuite vers les carabes et les salamandres \citep{Laigle2021Directindirect}. 
Nous ne pensons cependant pas que ce type de dynamique puisse avoir lieu dans notre cas, étant donné que nous n'avons mesuré aucun effet direct des traitements de coupes forestières sur la biomasse des collemboles. 

Les SEM (modèle d’équations structurelles) sont des outils puissants pour étudier la dynamique des milieux naturels, analyser les relations entre diverses variables sous forme de réseau et observer la propagation des effets dans des systèmes complexes, 
tels que les écosystèmes forestiers \citep{graceSpecificationStructuralEquation2010}.  
Plusieurs études ont souligné l’importance d’inclure les interactions biotiques, notamment entre espèces appartenant à différents niveaux trophiques, pour mieux comprendre les impacts des perturbations 
sur le fonctionnement des écosystèmes \citep{Thebault2003Foodwebconstraints,Seibold2018necessitymultitrophic,Laigle2021Directindirect}.  
Dans cette étude, l’intégration de modèles d’occupation et de modèles linéaires mixtes au SEM a ainsi permis d’adopter une approche multi-trophique sous forme de réseau, 
afin d’analyser les interactions entre différents taxons en réponse aux traitements de coupe forestière \citep{josephIntegratingOccupancyModels2016}.  
Cette méthodologie nous a permis de quantifier les effets directs et indirects des coupes forestières sur la faune du sol, tout en évaluant les modifications de plusieurs variables environnementales. 
Cependant, la complexité des SEM dépend fortement de la quantité de données disponibles, car l’utilisation d’un modèle complexe avec un jeu de données limité peut entraîner des problèmes d’estimation lors des analyses.  
Dans cette étude, les données recueillies se limitent à une seule année, ce qui restreint le niveau de complexité possible dans l'élaboration du SEM. 
Afin de refléter les potentielles différences annuelles, il serait essentiel de reproduire cette recherche pour intégrer la dimension temporelle dans l’analyse.

Nous avions l’intention d’utiliser des modèles \textit{N}-mélange pour quantifier l'impact des traitements sylvicoles sur l'abondance des différents taxons \citep{Royle2004Nmixturemodels,Mazerolle2021Woodlandsalamander}. 
Intégré au SEM, cette approche aurait permis d'acquérir une compréhension plus fine de l'impact des traitements sylvicoles sur nos espèces. 
Malheureusement, le nombre de salamandres observées n'a pas été suffisant pour utiliser ce type d'analyse. 
Il est possible que le manque de vieillissement des planches servant de refuge artificiel explique en partie le faible descompte des salamandres. 
En effet, la salamandre cendrée favorise habituellement les débris ligneux ayant un stade de décomposition avancé pour s'abriter \citep{Otto2011ComparingCover,hedrickEffectsCoverboardAge2021}. 
Dans cette étude, les planches ont été installées en milieu naturel seulement trois mois avant le début de l'échantillonnage, ce qui pourrait limiter leur utilisation par les salamandres. 
Concernant les carabes et les collemboles, les étapes d’identification et de pesée requièrent la collecte des individus. 
Le retrait d’individus ne respecte pas la condition de fermeture démographique de la population requise par les modèles \textit{N}-mélange  \citep{Royle2004Nmixturemodels}. 
Il n’était donc pas possible d’utiliser ces modèles d’abondance avec les données de carabes et de collemboles

Le dispositif développé dans le cadre du projet DREAM (Desired Regeneration through Assisted Migration) et utilisé dans notre étude a été conçu pour examiner la migration assistée en tant que solution pour préserver les forêts face aux changements climatiques \citep{royoDesiredREgenerationAssisted2023}. 
Cette méthode consiste à déplacer des individus ou du matériel génétique d’une région climatique originelle vers une zone mieux adaptée aux conditions futures \citep{Vitt2010Assistedmigration}. 
Cette approche permettrait de modifier rapidement la composition des peuplements pour les adapter aux climats futurs, répondant ainsi à des besoins de conservation \citep{Dumroese2015Considerationsrestoring,Park2018Informationunderload,Park2023Provenancetrials}. 
Cependant, un manque de connaissances et un degré d’incertitude subsistent autour de cette mesure d’adaptation \citep{Klenk2015assistedmigration,Park2018Informationunderload}. 
Il serait donc important d’évaluer si la migration assistée peut influencer la faune présente dans le milieu hôte.
Lors de notre étude, les semis issus de diverses origines géographiques étaient encore trop jeunes pour potentiellement modifier les conditions environnementales et influencer les groupes d’espèces étudiés. 
Il serait néanmoins pertinent de répéter cette expérience lorsque les arbres auront atteint une maturité suffisante pour modifier les conditions au sol. 
De plus, le dispositif que nous avons utiliser dans le cadre de cette étude se trouve en forêt mixte tempérée.
Toutefois, il est possible que la réponse des espèces aux perturbations puisse varier selon le type d'habitat dans lequel elles se trouvent \citep{Kudrin2023metaanalysiseffects}. 
Il serait intéressant de reproduire ce type d'expérience dans d'autres environnements, tels que les forêts feuillues ou conifériennes.

Enfin, une meilleure compréhension des relations entre les perturbations, les variables environnementales et la biodiversité est cruciale pour préserver la pérennité des écosystèmes forestiers. 
L’ensemble de cette recherche avait pour but d'approfondir la compréhension des relations existantes entre les pratiques sylvicoles et la dynamique des communautés du sol forestier. 
L'amélioration des connaissances permettra, à long terme, de fournir de meilleurs outils pour orienter les plans de gestion et adapter les pratiques sylvicoles.


\cleardoublepage

\bibliography{References.bib}
\bibliographystyle{ecologyNewFR.bst}            % conclusion

\appendix                       % annexes le cas échéant

\chapter{Supplementary material for the article}     % numérotée
\label{chap:supp}                   % étiquette pour renvois (à compléter!)

\pagebreak
\section{Carabids and collembola identifications}

\begin{table}[h]
    \centering
    \caption[List of carabid species captured during summer 2021 in Portneuf Wildlife Reserve and classification between salamanders competitor and prey groups.]
    {List of carabids captured during summer 2021 in Portneuf Wildlife Reserve}
    \label{tab:carabid}
    \begin{tabular}{lll} 
        \hline
        Group & Specie & Count \\ [0.5ex] 
        \hline      
        Competitor          & \textit{Chlaenius sericeus}               & 2 \\  
                            & \textit{Harpalus erythropus}              & 1 \\
                            & \textit{Harpalus faunus}                  & 1 \\
                            & \textit{Harpalus rufipes}                 & 1 \\
                            & \textit{Patrobus longicornis}             & 1 \\
                            & \textit{Platynus decentis}                & 2 \\
                            & \textit{Poecilus lucublandus}             & 3 \\
                            & \textit{Pterostichus adstrictus}          & 27 \\
                            & \textit{Pterostichus coracinus}           & 32 \\
                            & \textit{Pterostichus diligendus}          & 3 \\
                            & \textit{Pterostichus mutus}               & 10 \\
                            & \textit{Pterostichus pensylvanicus}       & 18 \\
                            & \textit{Pterostichus punctatissimus}      & 1 \\
                            & \textit{Pterostichus tristis}             & 25 \\
                            & \textit{Sphaeroderus canadensis}          & 6 \\
                            & \textit{Sphaeroderus nitidicollis}        & 1 \\
                            \hline 
        Prey                & \textit{Agonum affine}                    & 8 \\ 
                            & \textit{Agonum palustre}                  & 3 \\
                            & \textit{Agonum punctiforme}               & 2 \\ 
                            & \textit{Agonum retractum}                 & 3 \\ 
                            & \textit{Bradycellus lugubris}             & 1 \\
                            & \textit{Bradycellus semipubescens}        & 1 \\
                            & \textit{Calathus gregarius}               & 1 \\
                            & \textit{Harpalus providens}               & 1 \\
                            & \textit{Loricera pilicornis}              & 1 \\
                            & \textit{Notiobia terminata}               & 5 \\
                            & \textit{Pseudamara arenaria}              & 7 \\
                            & \textit{Pterostichus commutabilis}        & 1 \\
                            & \textit{Synuchus impunctatus}             & 20 \\
                            & \textit{Trechus apicalis}                 & 1 \\
                            \hline 
        \textbf{Total}      &                                           & 189 \\
        \hline
    \end{tabular}
  \end{table}

  \begin{table}[ht]
    \centering
    \caption[List of sringtails order and families]
    {List of sringtails order and families captured during summer 2021 in Portneuf Wildlife Reserve.}
    \label{tab:springtail}
    \begin{tabular}{lll} 
        \hline
        Order & Family & Count \\ [0.5ex] 
        \hline      
        Entomobryomorpha    & Entomobryidae     & 62 \\  
                            & Isotomidae        & 182 \\
                            & Tomoceridae       & 8 \\
        Poduromorpha        & Hypogastruridae   & 92 \\
                            & Neanuridae        & 11 \\
                            & Onychiuridae      & 33 \\
                            & Tullbergidae      & 24 \\
        Symphypleona        & Dicyrtomidae      & 6 \\
                            & Katiannidae       & 6 \\
                            & Neelidae          & 23 \\
                            & Sminthuridae      & 10 \\
                            & Sminthurididae    & 11 \\
                            \hline 
        \textbf{Total}      &                   & 468 \\
        \hline
    \end{tabular}
  \end{table}

  \clearpage

\section{Structural equation models}


\begin{table}[ht!]
\caption[Equations used to estimate impact of overstory treatments on Red-backed salamanders (\textit{Plethodon cinereus}) and ground beetle occupancy, springtail biomass and environmental variables that could effect soil fauna habitat selection.]
{JAGS code used to estimate impact of overstory treatments on Red-backed salamanders (\textit{Plethodon cinereus}) and ground beetle occupancy, springtail biomass and environmental variables that could effect soil fauna habitat selection in Portneuf Wildlife Reserve, Québec, Canada.}
\label{ann:SEM_equation}
\end{table}

\subsubsection{Occupancy model component} 

The occupancy models of salamanders and ground beetles used the observed detections and non-detections at site $i$ during survey $j$, 
represented by $\text{Salamander}_{ij}$, $\text{Large.carabids}_{ij}$ and $\text{Small.carabids}_{ij}$. These detections and non-detections followed a Bernoulli distribution with $z_{i} \times p_{ij}$ as the success parameter, 
where $z_{i}$ represents the latent occupancy state of a given species group at site $i$ and $p_{ij}$ represents the probability of detecting the same species group at site $i$ during survey $j$ :


\begin{align}
  \text{Salamander}_{ij} &\sim \text{Bernoulli}(z_{\text{Salamander}_i} \times p_{\text{Salamander}_{ij}}) \nonumber \\
  \text{Large.carabids}_{ij} &\sim \text{Bernoulli}(z_{\text{Large.carabids}_i} \times p_{\text{Large.carabids}_{ij}})  \\
  \text{Small.carabids}_{ij} &\sim \text{Bernoulli}(z_{\text{Small.carabids}_i} \times p_{\text{Small.carabids}_{ij}}) \nonumber
\end{align}


The latent occupancy state of salamanders and large ground beetles at a site ($z_{i}$) followed a Bernoulli distribution 
with probability of occupancy ($\psi$) of the given species for a given harvest treatment (control, partial cut, clearcut):


\begin{align}
  z_{\text{Salamander}_i} &\sim \text{Bernoulli}(\psi_{\text{Salamander}_{\text{Treat}_i}}) \nonumber \\
  z_{\text{Large.carabids}_i} &\sim \text{Bernoulli}(\psi_{\text{Large.carabids}_{\text{Treat}_i}})
\end{align}


We used uniform distribution priors ($\text{U}(0, 1)$) for the occupancy probabilities of salamanders and large ground beetles in a given treatment. 

For small ground beetles, the latent occupancy state was drawn from a Bernoulli distribution, where the occupancy probability ($\psi_{Small.carabids_{i}}$) 
depended on salamander latent occupancy ($z_{\text{Salamander}_i}$) and overstory treatments (control as reference):


\begin{align}
  \text{logit}(\psi_{\text{Small.carabids}_i}) &= \beta_{0[\text{Small.carabids}]} + \beta_{z_{\text{Salamander}}[\text{Small.carabids}]} \times z_{\text{Salamander}_i} + \nonumber \\
  &\beta_{\text{Cutpartial}[\text{Small.carabids}]} \times \text{Cutpartial}_i + \\
  &\beta_{\text{Cutclear}[\text{Small.carabids}]} \times \text{Cutclear}_i \nonumber
\end{align}

We assumed vague normal priors for the $\beta$ parameters, $\text{N}(0, \sigma = \sqrt{10})$. 
We allowed the detection probability of a given species groups ($\text{logit}(p_{ij})$) with the volume of coarse woody debris ($\text{CWD}_i$) and the precipitation levels 
($\text{Precipitation}_{ij}$) as explanatory variables, and a block random effect ($\alpha_{Block}$) to reflect the experimental design:


\begin{align}
  \text{logit}(p_{\text{Salamander}_{ij}}) &= \alpha_{0[\text{Salamander}]} + \alpha_{\text{CWD}[\text{Salamander}]} \times \text{CWD}_i + \nonumber \\
  &\alpha_{\text{Precipitation}[\text{Salamander}]} \times \text{Precipitation}_{ij} + \alpha_{\text{Block}[\text{Salamander}]_{\text{Block}_i}} \nonumber
\end{align}

\begin{align}
  \text{logit}(p_{\text{Large.carabids}_{ij}}) &= \alpha_{0[\text{Large.carabids}]} + \alpha_{\text{CWD}[\text{Large.carabids}]} \times \text{CWD}_i + \\
  &\alpha_{\text{Precipitation}[\text{Large.carabids}]} \times \text{Precipitation}_{ij} + \alpha_{\text{Block}[\text{Large.carabids}]_{\text{Block}_i}} \nonumber 
\end{align}

\begin{align}
  \text{logit}(p_{\text{Small.carabids}_{ij}}) &= \alpha_{0[\text{Small.carabids}]} + \alpha_{\text{CWD}[\text{Small.carabids}]} \times \text{CWD}_i + \nonumber \\
  &\alpha_{\text{Precipitation}[\text{Small.carabids}]} \times \text{Precipitation}_{ij} + \alpha_{\text{Block}[\text{Small.carabids}]_{\text{Block}_i}} \nonumber 
\end{align}


We used vague normal priors for CWD and precipitation levels $\text{N}(0, \sigma = \sqrt{10})$. 
The priors for block random effects were $\text{N}(0, \alpha_{Block})$, where $\alpha_{Block}$ was drawn from a uniform distribution $\text{U}(0, 10)$. 
Precipitation was highly correlated with air temperature and relative humidity. 
For this reason, we did not include these two variables in our models.


\subsubsection{Linear mixed model component} 

We used linear mixed models to assess how overstory treatments affect springtail biomass ($\text{Springtail}_{i}$) and 
environmental variables ($\text{CWD}_{i}$, $\text{Canopy}_{i}$, $\text{Litter}_{i}$) at site $i$ :

\begin{align}
  \text{Springtail}_{i} &\sim \text{N} (\mu_{\text{Springtail}_i}, \sigma_{\text{Springtail}_{\text{Group}_i}}) \nonumber \\
  \text{CWD}_{i} &\sim \text{N} (\mu_{\text{CWD}_i}, \sigma_{\text{CWD}_{\text{Group}_i}}) \\
  \text{Canopy}_{i} &\sim \text{N} (\mu_{\text{Canopy}_i}, \sigma_{\text{Canopy}_{\text{Group}_i}}) \nonumber \\
  \text{Litter}_{i} &\sim \text{N} (\mu_{\text{Litter}_i}, \sigma_{\text{Litter}_{i}}) \nonumber 
\end{align}

Due to heteroscedasticity in the Springtail biomass, we allowed each treatment group $j$ to have their own variances ($\sigma_j \sim \text{U}(0,150)$). 
Springtail biomass was drawn from a normal distribution $\text{N} (\mu_{\text{Springtail}_i}, \sigma_{\text{Springtail}_{\text{Group}_i}})$, where $\sigma_{\text{Springtail}_{\text{Group}_i}}$ denotes the 
residual variance of a given treatment group and ($\mu_{i}$) corresponds to the linear predictor including overstory treatments ($\text{Cutpartial}_i$, $\text{Cutclear}_i$), 
a block random effect ($\alpha_{\text{Block}}$), as well as the latent occupancy state of salamanders and of each ground beetle group 
($z_{Salamander}$, $z_{Large.carabids}$, $z_{Small.carabids}$):


\begin{align}
  \mu_{\text{Springtail}_i} &= \beta_{0[\text{Springtail}]} + \beta_{\text{Cutpartial}[\text{Springtail}]} \times \text{Cutpartial}_i + \nonumber\\
  &\beta_{\text{Cutclear}[Springtail]} \times \text{Cutclear}_i + \beta_{z_{\text{Salamander}}[\text{Springtail}]} \times z_{Salamander} +  \nonumber\\
  &\beta_{z_{\text{Small.carabids}}[\text{Springtail}]} \times z_{Small.carabids} + \beta_{z_{\text{Large.carabids}}[\text{Springtail}]} \times z_{Large.carabids} + \nonumber\\
  &\alpha_{\text{Block}[\text{Springtail}]_{\text{Block}_i}} \nonumber
\end{align}


Again, we assumed vague normal priors for the coefficients $\text{N}(0, \sigma = \sqrt{10})$. 
We used $\text{N}(0, \alpha_{Block})$ priors for the block random effects, where $\alpha_{Block})$ is drawn from a uniform distribution $\text{U}(0, 50)$. 


\begin{align}
  \mu_{\text{CWD}_i} &= \beta_{0[\text{CWD}]} + \beta_{\text{Cutpartial}[\text{CWD}]} \times \text{Cutpartial}_{i} + \nonumber\\
  & \beta_{\text{Cutclear}[\text{CWD}]} \times \text{Cutclear}_{i} + \alpha_{\text{Block}[\text{CWD}]_{\text{block}_i}} 
\end{align}

\begin{align}
  \mu_{\text{Canopy}_i} &= \beta_{0[\text{Canopy}]} + \beta_{\text{Cutpartial}[\text{Canopy}]} \times \text{Cutpartial}_{i} + \nonumber \\
  & \beta_{\text{Cutclear}[\text{Canopy}]} \times \text{Cutclear}_{i} + \alpha_{\text{Block}[\text{Canopy}]_{\text{block}_i}} \nonumber
\end{align}

\begin{align}
  \mu_{\text{Litter}_i} &= \beta_{0[\text{Litter}]} + \beta_{\text{Cutpartial}[\text{Litter}]} \times \text{Cutpartial}_{i} + \nonumber\\
  & \beta_{\text{Cutclear}[\text{Litter}]} \times \text{Cutclear}_{i} + \alpha_{\text{Block}[\text{Litter}]_{\text{block}_i}} \nonumber
\end{align}


\clearpage

\begin{table}[ht]
\caption[JAGS code used to estimate impact of overstory treatments on Red-backed salamanders (\textit{Plethodon cinereus}) and ground beetle occupancy, springtail biomass and environmental variables that could effect soil fauna habitat selection.]
    {JAGS code used to estimate impact of overstory treatments on Red-backed salamanders (\textit{Plethodon cinereus}) and ground beetle occupancy, springtail biomass and environmental variables that could effect soil fauna habitat selection in Portneuf Wildlife Reserve, Québec, Canada.}
    \label{ann:SEM_script}
\end{table}

\begin{lstlisting}
##SEM
model {

##CWD
##priors
beta0.cwd ~ dnorm(0, 0.01)
beta.Cutpartial.cwd ~ dnorm(0, 0.01)
beta.Cutclear.cwd ~ dnorm(0, 0.01)

##block random effect
for(m in 1:nblocks) {
   alpha.block.cwd[m] ~ dnorm(0, tau.block.cwd)
}

##variance of block
tau.block.cwd <- pow(sigma.block.cwd, -2)
sigma.block.cwd ~ dunif(0, 50)

##allow each group to have different variance
for(j in 1:3) {
   tau.cwd[j] <- pow(sigma.cwd[j], -2)
   sigma.cwd[j] ~ dunif(0, 150)
}

##iterate over each observation
for (i in 1:nsites) {

    ##linear predictor  
    mu.cwd[i] <- beta0.cwd + beta.Cutpartial.cwd * Cutpartial[i] + 
       beta.Cutclear.cwd * Cutclear[i] + alpha.block.cwd[Block[i]]

    ##response
    CWD_tot[i] ~ dnorm(mu.cwd[i], tau.cwd[Group[i]])
}

##derived parameters
for(i in 1:nsites) {
    pred.cwd[i] <- mu.cwd[i]
    res.cwd[i] <- CWD[i] - mu.cwd[i]
    res.pearson.cwd[i] <- res.cwd[i]/sigma.cwd[Group[i]]
}

##Canopy openness
##priors
beta0.can ~ dnorm(0, 0.01)
beta.Cutpartial.can ~ dnorm(0, 0.01)
beta.Cutclear.can ~ dnorm(0, 0.01)

##block random effect
for(m in 1:nblocks) {
   alpha.block.can[m] ~ dnorm(0, tau.block.can)
}

##variance of block
tau.block.can <- pow(sigma.block.can, -2)
sigma.block.can ~ dunif(0, 50)

##allow each group to have different variance
for(j in 1:3) {
   tau.can[j] <- pow(sigma.can[j], -2)
   sigma.can[j] ~ dunif(0, 150)
}

##iterate over each observation
for (i in 1:nsites) {

    ##linear predictor  
    mu.can[i] <- beta0.can + beta.Cutpartial.can * Cutpartial[i] + 
       beta.Cutclear.can * Cutclear[i] + alpha.block.can[Block[i]]

    ##response
    Canopy[i] ~ dnorm(mu.can[i], tau.can[Group[i]])
}

##derived parameters
for(i in 1:nsites) {
    pred.can[i] <- mu.can[i]
    res.can[i] <- Canopy[i] - mu.can[i]
    res.pearson.can[i] <- res.can[i]/sigma.can[Group[i]]
}

##Litter depth
##priors
beta0.lit ~ dnorm(0, 0.01)
beta.Cutpartial.lit ~ dnorm(0, 0.01)
beta.Cutclear.lit ~ dnorm(0, 0.01)

##block random effect
for(m in 1:nblocks) {
   alpha.block.lit[m] ~ dnorm(0, tau.block.lit)
}

##variance of block
tau.block.lit <- pow(sigma.block.lit, -2)
sigma.block.lit ~ dunif(0, 50)

##allow each group to have different variance
tau.lit <- pow(sigma.lit, -2)
sigma.lit ~ dunif(0, 150)

##iterate over each observation
for (i in 1:nsites) {

    ##linear predictor  
    mu.lit[i] <- beta0.lit + beta.Cutpartial.lit * Cutpartial[i] + 
       beta.Cutclear.lit * Cutclear[i] + alpha.block.lit[Block[i]]

       ##response
       ##Litter[i] ~ dnorm(mu.lit[i], tau.lit[Group[i]])
       Litter[i] ~ dnorm(mu.lit[i], tau.lit)
}

##derived parameters
for(i in 1:nsites) {
    pred.lit[i] <- mu.lit[i]
    res.lit[i] <- Litter[i] - mu.lit[i]
    res.pearson.lit[i] <- res.lit[i]/sigma.lit
}

##salamander component
##priors for psi
for(k in 1:ngroups) {
   psi.sal[k] ~ dunif(0, 1)
}

##priors for p
alpha0.sal ~ dnorm(0, 0.01)
alpha.precip.sal ~ dnorm(0, 0.01)
alpha.cwd.sal ~ dnorm(0, 0.01)

##block random effect
for(m in 1:nblocks) {
   alpha.block.sal[m] ~ dnorm(0, tau.block.sal)
}

tau.block.sal <- pow(sigma.block.sal, -2)
sigma.block.sal ~ dunif(0, 10)

##salamander model
for(i in 1:nsites) {
   ##occupancy
   z.sal[i] ~ dbern(psi.sal[Group[i]])
   for(j in 1:nvisits) {
      ##p
      logit.p.sal[i, j] <- alpha0.sal + alpha.cwd.sal * CWD[i] + alpha.precip.sal * Precip[i, j] + alpha.block.sal[Block[i]]
      p.sal[i, j] <- exp(logit.p.sal[i, j])/(1 + exp(logit.p.sal[i, j]))

      eff.p.sal[i, j] <- z.sal[i] * p.sal[i, j]
      y.sal[i, j] ~ dbern(eff.p.sal[i, j])
   }
}

finiteOcc.sal <- sum(z.sal[])

####carab beetles - prey
##psi priors
psi.beta0.car.prey ~ dnorm(0, 0.01)
psi.beta.z.sal.car.prey ~ dnorm(0, 0.01)
psi.beta.Cutpartial.car.prey ~ dnorm(0, 0.01)
psi.beta.Cutclear.car.prey ~ dnorm(0, 0.01)

##p priors
alpha0.car.prey ~ dnorm(0, 0.01)
alpha.precip.car.prey ~ dnorm(0, 0.01)
alpha.cwd.car.prey ~ dnorm(0, 0.01)

##block random effects
for(m in 1:nblocks) {
   alpha.block.car.prey[m] ~ dnorm(0, tau.block.car.prey)
}

tau.block.car.prey <- pow(sigma.block.car.prey, -2)
sigma.block.car.prey ~ dunif(0, 10)

for(i in 1:nsites) {
   ##occupancy
   logit(psi.car.prey[i]) <- psi.beta0.car.prey + psi.beta.z.sal.car.prey * z.sal[i] +
    	psi.beta.Cutpartial.car.prey*Cutpartial[i] + psi.beta.Cutclear.car.prey*Cutclear[i]
    	
   z.car.prey[i] ~ dbern(psi.car.prey[i])

   for(j in 1:nvisits) {
      ##p
      logit.p.car.prey[i, j] <- alpha0.car.prey + alpha.cwd.car.prey * CWD[i] + alpha.precip.car.prey * Precip[i, j]
      + alpha.block.car.prey[Block[i]]
      p.car.prey[i, j] <- exp(logit.p.car.prey[i, j])/(1 + exp(logit.p.car.prey[i, j]))

      eff.p.car.prey[i, j] <- z.car.prey[i] * p.car.prey[i, j]
      y.car.prey[i, j] ~ dbern(eff.p.car.prey[i, j])
   }
}

finiteOcc.car.prey <- sum(z.car.prey[])

####carab beetles - comp
##psi priors
for(k in 1:ngroups) {
   psi.car.comp[k] ~ dunif(0, 1)
}

##p priors
alpha0.car.comp ~ dnorm(0, 0.01)
alpha.precip.car.comp ~ dnorm(0, 0.01)
alpha.cwd.car.comp ~ dnorm(0, 0.01)

##block random effects
for(m in 1:nblocks) {
   alpha.block.car.comp[m] ~ dnorm(0, tau.block.car.comp)
}

tau.block.car.comp <- pow(sigma.block.car.comp, -2)
sigma.block.car.comp ~ dunif(0, 10)

for(i in 1:nsites) {
      ##occupancy
      z.car.comp[i] ~ dbern(psi.car.comp[Group[i]])

   	for(j in 1:nvisits) {
      	   ##detection
      	   logit.p.car.comp[i, j] <- alpha0.car.comp + alpha.cwd.car.comp * CWD[i] + alpha.precip.car.comp * Precip[i, j]
           + alpha.block.car.comp[Block[i]]

      	   ##p
      	   p.car.comp[i, j] <- exp(logit.p.car.comp[i, j])/(1 + exp(logit.p.car.comp[i, j]))
      	   eff.p.car.comp[i, j] <- z.car.comp[i] * p.car.comp[i, j]
      	   y.car.comp[i, j] ~ dbern(eff.p.car.comp[i, j])
      	}
}

finiteOcc.car.comp <- sum(z.car.comp[])

##collembola data
##priors
beta0.coll ~ dnorm(0, 0.01)
beta.Cutpartial.coll ~ dnorm(0, 0.01)
beta.Cutclear.coll ~ dnorm(0, 0.01)
beta.z.sal.coll ~ dnorm(0, 0.01)
beta.z.car.prey.coll ~ dnorm(0, 0.01)
beta.z.car.comp.coll ~ dnorm(0, 0.01)

##block random effect
for(m in 1:nblocks) {
   alpha.block.coll[m] ~ dnorm(0, tau.block.coll)
}

tau.block.coll <- pow(sigma.block.coll, -2)
sigma.block.coll ~ dunif(0, 50)

##allow each group to have different variance
for(j in 1:3) {
   tau.coll[j] <- pow(sigma.coll[j], -2)
   sigma.coll[j] ~ dunif(0, 150)
}

##iterate over each observation
for (i in 1:nsites) {

    ##linear predictor  
    mu.coll[i] <- beta0.coll +
    beta.Cutpartial.coll*Cutpartial[i] +
    beta.Cutclear.coll*Cutclear[i] + beta.z.sal.coll*z.sal[i] +
    beta.z.car.prey.coll*z.car.prey[i] +
    beta.z.car.comp.coll*z.car.comp[i] + alpha.block.coll[Block[i]]

    ##response
    y.coll[i] ~ dnorm(mu.coll[i], tau.coll[Group[i]])
}

##derived parameters
for(i in 1:nsites) {
    pred.coll[i] <- mu.coll[i]
    res.coll[i] <- y.coll[i] - mu.coll[i]
    res.pearson[i] <- res.coll[i]/sigma.coll[Group[i]]
}
}
\end{lstlisting}

\clearpage

\chapter*{Bibliographie}         
\label{chap:biblio}         
\phantomsection\addcontentsline{toc}{chapter}{\nameref{chap:biblio}} 

\nocite{*}
\renewcommand{\bibsection}{}
\begin{otherlanguage}{english}
\bibliography{References}
\bibliographystyle{ecologyNewEN}
\end{otherlanguage}
                % annexe A

%<!articles>\bibliography{}                 % production de la bibliographie
%<!articles>
\end{document}
%</gabarit>
%
% ^^A Gabarits des parties du document
%<*resume>
\chapter*{Résumé}               % ne pas numéroter
\label{chap-resume}             % étiquette pour renvois
\phantomsection\addcontentsline{toc}{chapter}{\nameref{chap-resume}} % inclure dans TdM

\begin{otherlanguage*}{french}
  <Texte du résumé en français. Obligatoire.>
\end{otherlanguage*}
%</resume>
%
%<*abstract>
\chapter*{Abstract}             % ne pas numéroter
\label{chap-abstract}           % étiquette pour renvois
\phantomsection\addcontentsline{toc}{chapter}{\nameref{chap-abstract}} % inclure dans TdM

\begin{otherlanguage*}{english}
  <Text of English abstract. Optional, but recommended.>
\end{otherlanguage*}
%</abstract>
%
%<*remerciements>
\chapter*{Remerciements}        % ne pas numéroter
\label{chap-remerciements}      % étiquette pour renvois
\phantomsection\addcontentsline{toc}{chapter}{\nameref{chap-remerciements}} % inclure dans TdM

<Texte des remerciements en prose.>
%</remerciements>
%
%<*avantpropos>
\chapter*{Avant-propos}         % ne pas numéroter
\label{chap-avantpropos}        % étiquette pour renvois
\phantomsection\addcontentsline{toc}{chapter}{\nameref{chap-avantpropos}} % inclure dans TdM

<Texte de l'avant-propos. Obligatoire dans une thèse ou un mémoire par
articles.>
%</avantpropos>
%
%<*introduction>
\chapter*{Introduction}         % ne pas numéroter
\label{chap-introduction}       % étiquette pour renvois
\phantomsection\addcontentsline{toc}{chapter}{\nameref{chap-introduction}} % inclure dans TdM

<Texte de l'introduction. La thèse ou le mémoire devrait normalement
débuter par une introduction. Celle-ci est traitée comme un chapitre
normal, sauf qu'elle n'est pas numérotée.>
%</introduction>
%
%<*chapitre>
\chapter{<Titre du chapitre ou de l'article>}     % numéroté
\label{chap-}                   % étiquette pour renvois (à compléter!)

%<articles>\section{Résumé}
%<articles>
%<articles>\begin{otherlanguage*}{french}
%<articles>  <Résumé de l'article en français. Obligatoire.>
%<articles>\end{otherlanguage*}
%<articles>
%<articles>\section{Abstract}
%<articles>
%<articles>\begin{otherlanguage*}{english}
%<articles>  <English abstract of the paper. Optional, but recommended.>
%<articles>\end{otherlanguage*}
%<articles>
<Texte du chapitre ou de l'article.>
%<articles>
%<articles>\bibliographystyle{}              % style de la bibliographie
%<articles>\bibliography{}                   % production de la bibliographie
%</chapitre>
%
%<*conclusion>
\chapter*{Conclusion}           % ne pas numéroter
\label{chap-conclusion}         % étiquette pour renvois
\phantomsection\addcontentsline{toc}{chapter}{\nameref{chap-conclusion}} % inclure dans TdM

<Texte de la conclusion. Une thèse ou un mémoire devrait normalement
se terminer par une conclusion placée avant les annexes, le cas
échéant. La conclusion est traitée comme un chapitre normal, sauf
qu'elle n'est pas numérotée.>
%</conclusion>
%
%<*annexe>
\chapter{<Titre de l'annexe>}     % numérotée
\label{chap-}                   % étiquette pour renvois (à compléter!)

<Texte de l'annexe.>
%</annexe>
% \fi
